
\medterm{T Cell} Abbreviation for T lymphocyte, a type of white blood cell that is produced in the bone marrow and is part of the body\&rsquo;s immune system.

\medterm{T Wave} The electrocardiographic deflection representing ventricular repolarization. Its shape and direction
vary by lead and clinical context. Abnormal T waves may occur with ischemia, electrolyte disorders, ventricular strain,
or other conditions.

\medterm{T-Cell Count} A laboratory measurement of T lymphocytes in blood, often including CD4 cells, used to assess immune function and monitor conditions such as HIV infection or immunosuppression.

\medterm{T-Cell Lymphoma} A cancer of T lymphocytes, a type of white blood cell. It includes several different diseases that may affect lymph nodes, skin, blood, or organs and can cause swollen nodes, fever, weight loss, rash, or other symptoms.
\medterm{T-Cell Lymphomas} Lymphomas arising from mature or immature T lymphocytes. They include several biologically distinct diseases that may
involve lymph nodes, skin, blood, or other organs; behavior and treatment depend on the
subtype.

\medterm{T-Cell Receptor} A clonally distributed receptor on T lymphocytes that recognizes a peptide antigen
presented by a major histocompatibility complex molecule. TCR signaling requires
accessory and co-stimulatory signals for effective T-cell activation.

\medterm{T-Cells (T-Lymphocytes)} A type of lymphocyte that matures in the thymus and is the central effector of
cell-mediated immunity. T-cells express the T-cell receptor (TCR) on their surface,
which recognises peptide antigens presented by major histocompatibility complex (MHC)
molecules on antigen-presenting cells.

\synonyms
T-Lymphocytes

\medterm{T-Lymphocyte} A white blood cell that matures in the thymus and directs cell-based immune responses. Different T cells kill infected cells, coordinate immunity, remember previous threats, or reduce excessive immune activity.
\medterm{T-Wave Alternans} T-wave alternans is a beat-to-beat change in the size or shape of the heart's T wave and may indicate increased risk of dangerous ventricular rhythms.

\medterm{T3} Triiodothyronine, the most biologically active thyroid hormone, produced by deiodination
of T4 in peripheral tissues (liver, kidney) and secreted by the thyroid. Regulates
metabolism (BMR), protein/carbohydrate/lipid metabolism, growth, and development.

\medterm{T3 Test} A blood test measuring triiodothyronine, used with thyroid-stimulating hormone and thyroxine to evaluate thyroid function, particularly suspected hyperthyroidism.

\medterm{T3RU Test} An older test estimating the proportion of thyroid-hormone-binding sites occupied in serum; it is used indirectly to assess thyroid-hormone binding and has largely been replaced by free-thyroxine measurements.

\medterm{T4} Thyroxine, the main hormone secreted by the thyroid gland (thyroxine, 80-90 percent of
secretion), converted to the active T3 in tissues. Bound largely to thyroxine-binding
globulin (TBG); free T4 is the active unbound fraction. Regulates metabolism, growth,
and development.

\medterm{TAB Vaccine} An older vaccine against typhoid fever and paratyphoid A and B. It has largely been replaced by newer typhoid vaccines and does not protect against every cause of fever or every type of paratyphoid infection.
\medterm{Tabes} An older term for tabes dorsalis, a late neurologic complication of untreated syphilis.
\medterm{Tabes Dorsalis} A late manifestation of untreated syphilis in which degeneration of dorsal spinal-cord columns and roots causes sensory ataxia, lightning pains, loss of position sense, areflexia, and bladder or sexual dysfunction.

\medterm{Tablet} A solid dosage form of compressed drug plus excipients (binders, fillers, disintegrants,
lubricants, coatings). Types: immediate-release, enteric-coated (resistant to gastric
acid), sustained/controlled-release, sublingual, buccal, effervescent, chewable, and
orally disintegrating.

\medterm{Tablet Splitting} Dividing a tablet to take a smaller dose; only tablets designed for splitting should be divided, because extended-release, enteric-coated, and some specialized tablets may not work safely when split.

\medterm{TAC Regimen} A chemotherapy combination of docetaxel, doxorubicin, and cyclophosphamide, used for selected breast cancers; it can cause infection risk, nausea, hair loss, heart effects, and other toxicities.

\medterm{Tacalcitol} A vitamin-D-like medicine applied to skin affected by psoriasis or selected scaling disorders. It helps regulate skin-cell growth and inflammation but can irritate skin or raise calcium if overused.
\medterm{Tacedinaline} An investigational medicine that inhibits histone deacetylases, enzymes that influence gene activity. It has been studied for slowing cancer-cell growth but is not an established routine treatment.
\medterm{Tache Noire} Tache noire is a black eschar at the site of a tick bite, classically associated with Mediterranean spotted fever and related rickettsial infections.

\medterm{Tachyarrhythmia} An abnormally fast heart rhythm, which may arise from the atria, ventricles, or the heart's conduction
system. Symptoms can include palpitations, dizziness, chest discomfort, shortness of breath,
or fainting; fainting or severe symptoms require urgent assessment.

\medterm{Tachycardia} A faster-than-expected heart rate that may be a normal response or a sign of disease. See Fast Heart Rate.

\synonyms
Rapid Heartbeat

\medterm{Tachycardia Paroxysmal Atrial (Paroxysmal Supraventricular Tachycardia)} Alternate index label for Paroxysmal Supraventricular Tachycardia.

\medterm{Tachycardia Paroxysmal Supraventricular (Paroxysmal Supraventricular Tachycardia)} Alternate index label for Paroxysmal Supraventricular Tachycardia.

\medterm{Tachycardia-Induced Cardiomyopathy} Weakening and dilation of the heart muscle caused by a persistent or frequently recurring rapid heart rhythm, often improving when the arrhythmia is controlled.

\medterm{Tachykinin} A family of signaling peptides, including substance P, that influence pain, inflammation, smooth-muscle activity, and nerve communication.

\medterm{Tachyphrenia} Tachyphrenia is rapid flow of thoughts or accelerated mental activity, often described in association with mania, stimulants, or other psychiatric states.

\medterm{Tachyphylaxis} A rapid, acute decrease in response to a drug after repeated administration over a short
period, distinct from tolerance which develops more gradually. Tachyphylaxis results
from rapid depletion of neurotransmitter stores, receptor desensitization, or receptor
downregulation.

\medterm{Tachypnea} Abnormally rapid respiratory rate (adult >20 breaths/min). Physiologic (exercise,
anxiety) or pathologic: hypoxia, fever, sepsis, metabolic acidosis (Kussmaul), pulmonary
embolism, pneumonia, heart failure, pain, and compensation for metabolic derangement.

\medterm{Tachypnoea} Abnormally rapid breathing, defined as a respiratory rate >20 breaths per minute in
adults (age-dependent: >60 in neonates, >40 in 1--12 months, >30 in 1--5 years, >20 in
>5 years). Tachypnoea is a non-specific sign of respiratory or metabolic distress.

\medterm{Tachysystole (Uterine Hyperstimulation)} More than 5 contractions in 10 minutes averaged over 30 minutes, or contractions lasting
\(\geq 2\) minutes. Tachysystole with fetal heart rate changes (late decelerations, prolonged
deceleration, loss of variability) is more concerning. Causes: oxytocin (most common),
prostaglandins (misoprostol > dinoprostone), and cocaine.

\synonyms
Uterine Hyperstimulation

\medterm{TACO} Transfusion-associated circulatory overload is breathing difficulty caused by too much fluid entering the circulation during or soon after a blood transfusion. It can cause lung fluid, high blood pressure, and swelling, especially in people with heart or kidney disease or those receiving rapid or large transfusions.

\medterm{Tacrine} An older acetylcholinesterase inhibitor formerly used for Alzheimer disease; liver toxicity and frequent dosing limited its use.

\medterm{Tacrolimus} An immune-suppressing medicine used after transplants and on skin for eczema; it can increase infection risk.
\medterm{Tactile} Relating to the sense of touch, including feeling pressure, texture, vibration, temperature, or pain. Touch signals travel through nerves to the brain; reduced or abnormal touch sensation may result from nerve, spinal-cord, brain, or skin disease.
\medterm{Tactile Agnosia} Tactile agnosia is inability to recognize a familiar object by touch despite intact primary sensation, usually from parietal-cortex dysfunction.

\medterm{Tactile Fremitus} Vibrations transmitted through the lung to the chest wall, felt during palpation. It is
increased in consolidation (pneumonia) due to better sound transmission through solid
tissue. It is decreased in pleural effusion, pneumothorax, and atelectasis due to
impaired transmission.

\medterm{Tactile Sense} The ability to perceive touch, pressure, vibration, texture, and related mechanical stimuli through receptors in the skin and deeper tissues.

\medterm{Tadalafil} Tadalafil is a phosphodiesterase-5 inhibitor used for erectile dysfunction, benign prostatic hyperplasia, and pulmonary hypertension; headache, flushing, hypotension, priapism, and dangerous nitrate interactions are important risks.

\medterm{Taenia} A genus of tapeworms, including T. solium and T. saginata.
\medterm{Taenia Saginata} Taenia saginata is the beef tapeworm acquired from undercooked infected beef; intestinal infection may cause abdominal symptoms or passage of motile segments but does not usually cause human cysticercosis.

\medterm{Taenia Solium} A pork tapeworm that infects people who eat undercooked pork or swallow its eggs. Adult worms live in the intestine, while eggs can cause cysticercosis, including dangerous infection of the brain.
\medterm{Taeniasis} Intestinal infection caused by adult Taenia tapeworms, usually acquired by eating undercooked beef or pork containing larval cysts. It may cause abdominal symptoms or passage of worm segments; swallowing Taenia solium eggs causes cysticercosis in tissues, which is a different infection.
\medterm{Taeniasis (Tapeworm Infection)} Intestinal infection caused by adult Taenia tapeworms, usually acquired by eating undercooked
meat containing larval cysts. Taenia solium and Taenia saginata are important species;
Taenia solium eggs can also cause cysticercosis in tissues.

\synonyms
Tapeworm Infection

\medterm{Tail} A tapering appendage extending from the end of an animal's body. In human anatomy, the term may refer to a tail-like congenital structure or the embryologic caudal end; evaluation depends on the specific finding.

\medterm{Tailbone Trauma} Injury to the coccyx, the small bone at the bottom of the spine, usually after a fall or prolonged pressure. It causes pain when sitting, standing, or passing stool. Most injuries are bruises and improve with time, cushions, and pain relief; persistent or severe pain needs medical assessment.


\medterm{Taint} A taint is contamination or an undesirable influence; in medicine or laboratory work it may refer to contamination that affects a specimen, product, or result.

\medterm{Takayasu Arteritis} A chronic granulomatous vasculitis primarily affecting the aorta and its major branches,
most commonly diagnosed in women under 50 years of age, particularly in Asian
populations.

\medterm{Takayasu Disease} A large-vessel vasculitis that primarily affects the aorta and its major branches. It may cause limb claudication,
weak or unequal pulses, blood-pressure differences, bruits, fever, fatigue, or arterial narrowing and aneurysms.
Imaging and inflammatory assessment guide immunosuppressive and vascular treatment.

\medterm{Takayasu’S Arteritis} A chronic granulomatous vasculitis of the aorta and its major branches, predominantly
affecting young women of Asian descent. Causes arterial stenosis, occlusion, or aneurysm
formation. Presents with limb claudication, absent pulses, bruit, hypertension, and
constitutional symptoms. Treatment: corticosteroids and immunosuppressives.

\medterm{Takotsubo Cardiomyopathy} Transient left ventricular apical ballooning resembling a Japanese octopus trap
(takotsubo), typically triggered by emotional or physical stress. ECG may show ST
elevation mimicking STEMI. Coronary angiography shows no obstructive disease. Most
patients recover normal ventricular function within weeks.

\medterm{Talabostat} Talabostat is an investigational inhibitor of fibroblast activation protein and related enzymes studied for immune and anticancer effects.

\medterm{Talactoferrin Alfa} Talactoferrin alfa is a recombinant lactoferrin-related immunomodulatory protein studied as a treatment or adjunct in cancer and infectious disease.

\medterm{Talaporfin Sodium} Talaporfin sodium is a photosensitizer used or studied in photodynamic therapy, where light activation produces local cytotoxic effects.

\medterm{Talc} A soft mineral; inhalation of talc dust can cause talcosis, and contaminated injected talc can injure vessels and lungs.
\medterm{Talcosis} Pneumoconiosis from talc dust inhalation or IV injection of crushed oral medications.
Causes upper lobe nodular fibrosis. HRCT may show ``galaxy sign'' of calcified nodules.

\medterm{Talcum Powder Poisoning} Harmful effects from swallowing or inhaling talcum powder. Swallowing usually causes mild stomach upset, but inhaling it can injure the lungs and make breathing difficult.

\medterm{Talipes} A deformity of the foot and ankle present at birth or acquired later. It may affect position, flexibility, walking, and shoe fit, and treatment depends on severity, cause, age, and function.
\medterm{Talipes Calcaneovalgus} A foot deformity in which the heel is turned outward and the ankle is dorsiflexed so the foot points upward; flexible newborn cases often improve, while rigid cases need evaluation.

\medterm{Talipes Equinovarus} Clubfoot, a congenital deformity combining ankle equinus, hindfoot varus, forefoot adduction, and cavus; early serial casting and bracing are standard initial treatment.

\medterm{Tall Stature} Tall stature means height substantially above the expected range for age and sex and may be familial, endocrine, genetic, or associated with another condition.

\medterm{Talotrexin} Talotrexin is an investigational antifolate antineoplastic drug designed to inhibit folate-dependent cellular metabolism. An experimental cancer medicine designed to block folate-dependent cell processes needed for uncontrolled tumor growth.
\medterm{Talus} Second largest tarsal bone articulating with tibia/fibula proximally at the ankle.
Approximately 60 percent of its surface is articular cartilage with no muscular
attachments. Transmits entire body weight to the calcaneus. Talar neck fractures may
disrupt blood supply causing avascular necrosis in up to 50 percent.

\medterm{Tamoxifen (Nolvadex)} A medicine, biologic, or pharmacologic term related to tamoxifen (nolvadex); its mechanism, indication, interactions, and adverse effects depend on the specific agent.

\medterm{Tamoxifen Citrate} The citrate salt of tamoxifen, a selective estrogen-receptor modulator used to treat or prevent some breast cancers; blood clots and endometrial cancer are important uncommon risks.

\medterm{Tampon} A material inserted into the vagina to absorb menstrual blood; prolonged use can rarely lead to toxic shock syndrome.
\medterm{Tamponade} Compression of an organ or vessel by accumulated fluid, blood, or gas; cardiac tamponade impairs cardiac filling.
\medterm{Tamsulosin Hydrochloride} The hydrochloride salt of tamsulosin, an alpha-1 blocker that relaxes prostate and bladder-neck muscle to improve urinary symptoms from benign prostatic enlargement.

\medterm{Tandutinib} An experimental medicine studied for some blood disorders and cancers. It blocks signals that help certain cells grow, but it is not a routine treatment and its benefits and risks remain dependent on research findings.

\medterm{Tanespimycin} An investigational derivative of 17-AAG that inhibits heat-shock protein 90 and has been studied as an anticancer drug.

\medterm{Tangier Disease} A rare inherited disorder causing extremely low HDL cholesterol and buildup of cholesterol in tissues. It may cause enlarged yellow-orange tonsils, enlarged liver or spleen, nerve symptoms, or clouding of the cornea; cardiovascular risk varies and specialist testing is needed.

\medterm{Tanner Staging} A scale for assessing sexual maturity during puberty. Stage 1: prepubertal. Stage 2: breast buds (girls), testicular enlargement (boys). Stage 3: breast enlargement, pubic hair. Stage 4: areolar development, more pubic hair.

\medterm{Tannic Acid} A plant-derived polyphenol with astringent and protein-binding properties, used in some industrial or historical preparations; concentrated ingestion can irritate the gastrointestinal tract and affect absorption.

\medterm{Tanomastat} Tanomastat is an experimental medicine that blocks enzymes involved in tissue breakdown and remodeling. It has been studied for limiting tumor spread, but its safety, effectiveness, and clinical role remain unestablished.

\medterm{Tantalum} A corrosion-resistant metal used in some surgical implants and medical devices; exposure risk depends on the form, dose, and route.

\medterm{Tantulum} A misspelling of tantalum, a metal used in selected implants and medical devices.
\medterm{Tap, Joint (Aspiration)} A needle procedure that removes fluid from a joint for testing or symptom relief, helping assess infection, bleeding, crystals, or inflammation.

\medterm{Tapentadol} Tapentadol is a centrally acting analgesic with mu-opioid agonist and norepinephrine-reuptake-inhibitory effects, used for moderate-to-severe pain; respiratory depression, dependence, seizures, and serotonin or sedative interactions are possible.

\medterm{Tapeworm} A parasitic flatworm with a segmented body that may live in the intestine or form larval cysts in tissues.
People acquire different species by eating infected meat or fish or by ingesting eggs from
contaminated sources. Symptoms and complications depend on the species and site.

\medterm{Tapeworm Infection} Infection caused by parasitic tapeworms. Taenia species may live as adult worms in the intestine, while some tapeworm larvae can form cysts in tissues.

\medterm{Tapeworm Infection - Beef Or Pork} Intestinal infection with Taenia saginata or Taenia solium acquired from undercooked beef or pork; T. solium eggs can also cause cysticercosis in tissues.

\medterm{Tapeworm Infection - Hymenolepis} Intestinal infection with Hymenolepis nana or H. diminuta, acquired by ingesting eggs or infected arthropods and causing variable gastrointestinal symptoms.

\medterm{Tapotement} A massage technique using rhythmic tapping, patting, or percussion. It should be avoided over injured areas or in people at risk of bleeding unless a trained clinician recommends it.
\medterm{Tapping} Repeated light striking or percussion used during an examination, treatment, or test. Tapping may help assess reflexes, pain, body sounds, or fluid and may also be used in chest physiotherapy or other treatments.
\medterm{TAPSE} Alternate name; see \textit{Tricuspid Annular Plane Systolic Excursion}.

\medterm{TAPVR} Alternate index label; see \textit{Total anomalous pulmonary venous return}.

\medterm{Taq Polymerase} A heat-stable DNA-copying enzyme from Thermus aquaticus, widely used in polymerase chain reaction testing and research.

\medterm{Tar Remover Poisoning} Poisoning from swallowing, inhaling, or skin contact with tar-removing chemicals. Effects vary by product and may include irritation, dizziness, vomiting, or breathing problems; contact poison-control services promptly.

\medterm{Tarantism} Tarantism is a historical term for a dancing or behavioral syndrome once attributed to spider bite and now understood as a cultural or neurologic phenomenon rather than a specific disease.

\medterm{Tarantula Spider Bite} A bite from a tarantula that usually causes local pain, redness, and swelling. Some species also release irritating hairs that can affect the skin or eyes.

\medterm{Tardive Dyskinesia} Involuntary, repetitive, purposeless movements caused by prolonged use of antipsychotic
medications, typically developing after months to years of exposure. Movements include
choreoathetoid movements of the tongue, face, and extremities (lip smacking, tongue
protrusion, grimacing, choreiform movements of limbs).

\medterm{Tarenflurbil} Tarenflurbil is an investigational derivative of flurbiprofen studied for effects on amyloid processing and Alzheimer disease.

\medterm{Target Cell} A cell specifically affected by or responding to a hormone, drug, toxin, antigen, or other signal; target cells may also describe abnormal erythrocytes.
\medterm{Target Organ} An organ especially affected by a disease, hormone, toxin, medicine, or other exposure. Effects may occur because the organ has specific receptors, receives the substance, or is especially vulnerable to damage.
\medterm{Target Population} The defined group of people for whom a medical intervention, screening program, or research study is intended. It is characterized by criteria such as age, diagnosis, risk, or other relevant features.

\medterm{Target Tissue} A tissue containing receptors or other molecular targets on which a hormone, drug, toxin, or signaling molecule acts.

\medterm{Target-Controlled Infusion} A computerized infusion pump system that maintains a predetermined drug concentration in
plasma or effect site using pharmacokinetic modeling. TCI systems use three-compartment
pharmacokinetic models to calculate infusion rates that achieve and maintain target
concentration.

\medterm{Targeted Radiation Therapy} Radiation directed precisely at a tumor or other abnormal tissue to damage its DNA while limiting exposure to nearby healthy tissue. It may be delivered from outside the body or through an implanted source.
\medterm{Targeted Temperature Management} Controlled therapeutic hypothermia for comatose survivors of cardiac arrest. Target:
32-36 degrees C for 24-72 hours followed by gradual rewarming. Mechanism: reduces
metabolic demand, limits reperfusion injury, and reduces free radical generation. TTM2
trial showed 33 degrees C was not superior to 36 degrees C.

\medterm{Targeted Therapies For Cancer} Cancer medicines designed to act on specific features of cancer cells or their surroundings. They may be more selective than traditional chemotherapy, but can still cause important side effects.

\medterm{Targeted Therapy} Treatment that acts on a specific change or feature helping cancer cells grow, survive, or spread. Testing the tumor may identify a target and predict benefit. These treatments can be more selective than traditional chemotherapy but may still cause serious side effects and resistance.


\medterm{Tariquidar} Tariquidar is an investigational inhibitor of the P-glycoprotein drug transporter studied to reverse multidrug resistance and alter medicine distribution.

\medterm{Tarlov Cyst} A cerebrospinal-fluid-filled perineural cyst, usually located near the dorsal root
ganglia in the sacral spine. Many are incidental, but larger or symptomatic cysts can
cause sacral pain, radiculopathy, sensory change, bowel or bladder dysfunction, or
sexual dysfunction.

\medterm{Tars} A thick black substance from organic material; tobacco tar contains chemicals that can damage lungs and cause cancer.
\medterm{Tarsal} Relating to the tarsus, the group of bones forming the ankle and hindfoot. Tarsal bones help support body weight and allow foot movement.
\medterm{Tarsal Bones} The seven bones of the hindfoot and midfoot that connect the ankle to the metatarsals and form the arches and load-bearing framework of the foot.

\synonyms
Tarsals

\medterm{Tarsal Coalition} An abnormal connection between two bones in the rear or middle foot, usually present from birth. It can
cause a stiff, painful flatfoot, recurrent ankle sprains, and limited foot motion, often becoming
more noticeable during adolescence.

\medterm{Tarsal Joints} The articulations among the tarsal bones and between the tarsals and metatarsals, providing stability and controlled movement for walking and adapting the foot to surfaces.

\medterm{Tarsal Tunnel Syndrome} Pressure on the tibial nerve as it passes through a narrow space near the inner ankle. It can cause burning pain, tingling, or numbness in the sole and may result from injury, swelling, diabetes, or a mass. Treatment depends on the cause.

\medterm{Tarsorrhaphy} A procedure that partially or completely closes the eyelids to protect the cornea or reduce exposure in conditions such as severe dry eye or facial nerve palsy.

\medterm{Tarsus} The collection of seven bones forming the posterior part of the foot, between the tibia/fibula and the metatarsals. The group of seven bones forming the hindfoot and midfoot.
\medterm{Tartar} A hard deposit of mineralized dental plaque, also called calculus.
\medterm{Tarui Disease} A rare inherited muscle-energy disorder caused by deficiency of an enzyme needed to break down sugar for energy. Exercise can trigger early fatigue, cramps, muscle breakdown, and dark urine; some people also have chronic anemia or high uric acid.

\synonyms
Glycogen Storage Disease Type VII (GSD VII); Muscle Phosphofructokinase Deficiency

\medterm{Tassinari Syndrome} An older name associated with epileptic encephalopathy featuring frequent abnormal electrical activity during sleep, seizures, and possible effects on learning or development. The term is now usually described more specifically as electrical status epilepticus in sleep.

\medterm{Taste} The sensation produced by the interaction of chemical substances with taste receptor
cells (TRCs) on the tongue, palate, pharynx, epiglottis, and upper esophagus. Taste is
one of the five primary senses (gustation).

\medterm{Taste - Impaired} Impaired taste is reduced, distorted, or absent ability to taste and may result from infection, medicines, oral disease, nerve injury, or changes in smell.

\medterm{Taste Bud} A taste bud is a cluster of specialized sensory cells in papillae of the tongue and some oral tissues that detects sweet, sour, salty, bitter, and umami stimuli and transmits signals through cranial nerves.

\medterm{Taste Disorder} Abnormal reduction, loss, distortion, or unpleasant alteration of taste, caused by infection, medications, oral disease, nerve injury, neurologic disease, or other factors.



\medterm{Tau PET} Positron-emission tomography using a tracer that binds to abnormal tau protein deposits. It is used mainly in research and selected clinical evaluations of neurodegenerative disease.

\medterm{Taurodontism} A tooth-development difference with an enlarged pulp chamber and shortened roots, usually affecting molars. It may occur alone or with genetic conditions and can complicate dental treatment or root-canal procedures.

\medterm{Taussig-Bing Anomaly} A rare birth defect in which both major arteries leave the right lower heart chamber and a hole lies between the lower chambers below the pulmonary artery. It can cause severe low oxygen and heart failure in infancy and usually requires early specialist surgery.

\synonyms
Taussig-Bing Syndrome; DORV with Subpulmonic VSD

\medterm{TAVR} Alternate name; see \textit{Transcatheter Aortic Valve Replacement}.

\medterm{Taxanes} Chemotherapy medicines, including paclitaxel and docetaxel, that disrupt cell division by stabilizing tiny internal structures. They treat several cancers but may cause low blood counts, nerve symptoms, hair loss, allergic reactions, or fluid retention.
\medterm{Taxis} Directed movement toward or away from a stimulus; in surgery it may also mean orderly repositioning.

\medterm{Tay-Sachs Disease} A rare inherited disorder in which a fatty substance builds up in nerve cells because of an enzyme deficiency. The infant form causes loss of skills, weakness, an exaggerated startle, and seizures; later forms progress more slowly. Supportive care and genetic counseling are important.

\medterm{Tazarotene} A topical retinoid used for acne, plaque psoriasis, and some sun-damaged skin. It can irritate skin and increases sun sensitivity; pregnancy precautions are important.
\medterm{TB} Alternate index label; see \textit{Tuberculosis}.

\medterm{TB Test} A tuberculosis test looks for infection with Mycobacterium tuberculosis using a skin test, blood test, or examination of respiratory specimens.

\medterm{TBG Blood Test} A blood test measuring thyroxine-binding globulin, the principal thyroid-hormone transport protein; it helps interpret abnormal total thyroid-hormone levels when binding-protein changes are suspected.

\medterm{TBI} Alternate index label; see \textit{Traumatic brain injury}.

\medterm{TCA Cycle} The tricarboxylic acid cycle, also called the Krebs cycle, is a series of reactions inside mitochondria that breaks down products of carbohydrates, fats, and proteins. It produces carbon dioxide and energy-carrying molecules used to make cellular energy.

\medterm{TCAs (Tricyclic Antidepressants)} Antidepressants that block reuptake of serotonin and norepinephrine, with additional
anticholinergic, antihistaminic, and alpha-adrenergic blocking effects. Examples include
amitriptyline, nortriptyline, imipramine, and desipramine. They are effective for
depression, neuropathic pain, and migraine prophylaxis.

\synonyms
Tricyclic Antidepressants

\medterm{Td} Td is a vaccine protecting against tetanus and diphtheria, generally used as a booster in adolescents and adults according to local schedules.


\medterm{Td Immunization} Vaccination against tetanus and reduced-dose diphtheria, used for routine boosters, catch-up immunization, and wound-related protection according to age and vaccination history.


\medterm{Tdap Vaccination In Pregnancy} Administration of tetanus, diphtheria, and acellular pertussis (Tdap) vaccine during
each pregnancy, optimally at 27-36 weeks. The goal: boost maternal antibodies to protect
the newborn from pertussis (whooping cough) before the infant can be vaccinated at 2
months.

\medterm{Tear} A split or rupture in tissue, such as a muscle, tendon, ligament, meniscus, skin, or organ, caused by trauma, overuse, degeneration, or pressure.

\medterm{Tearing} Excessive production or inadequate drainage of tears, commonly caused by irritation, allergy, dry-eye reflex tearing, infection, eyelid malposition, or lacrimal-duct obstruction.

\medterm{Tears} Fluid made mainly by the lacrimal glands that lubricates, nourishes, cleans, and protects the eye's surface. Tears contain water, salts, oils, mucus, and immune substances that help maintain clear, comfortable vision.
\medterm{Tears, Decreased Production} Alternate index label; see \textit{Dry eyes}.

\medterm{Technetium-99} A radioactive isotope related to technetium-99m; technetium isotopes are used in nuclear medicine and radioactive decay studies.
\medterm{Technetium-99M} Most widely used radionuclide in nuclear medicine. Half-life: 6 hours. Emits 140 keV
gamma rays. Produced from molybdenum-99/technetium-99m generator. Used with various
ligands: MDP (bone), sestamibi (myocardium), DTPA (renal), nanocolloid (lymphatic),
HMPAO (brain perfusion).

\medterm{Technique} The method, equipment, and sequence of steps used to perform a medical examination, procedure, or imaging study. Correct technique improves safety, accuracy, and reproducibility.
\medterm{Tedizolid Pharmacokinetics} The study of how tedizolid is absorbed, distributed, metabolized, and eliminated.
Tedizolid is available orally and intravenously and has pharmacokinetic properties that support once-daily treatment;
its clinical use depends on the infection and patient factors.

\medterm{TEE} TEE, or transesophageal echocardiography, uses an ultrasound probe in the esophagus to obtain detailed images of heart valves, chambers, clots, and nearby vessels.

\medterm{Teen Depression} A depressive disorder in an adolescent that causes persistent low mood or loss of interest and impairs daily functioning.

\synonyms
Depression, Teen

\medterm{Teenage Pregnancy} Pregnancy in adolescents (usually 10-19 years), associated with higher maternal and
fetal risks: preterm birth, low birth weight, preeclampsia, anemia, and increased
neonatal mortality; also socioeconomic and educational consequences. Causes: unprotected
intercourse, lack of contraception/education, poverty, and peer/family factors.

\medterm{Teenagers And Drugs} Drug use during adolescence can affect brain development, school and relationships, safety, and mental health; early, nonjudgmental screening and support can reduce harm.

\medterm{Teenagers And Sleep} Teenagers commonly need about eight to ten hours of sleep, but schedules, screens, caffeine, mood, school demands, and sleep disorders can reduce sleep quality.

\medterm{TEER} Alternate name; see \textit{Transcatheter Edge-to-Edge Repair}.

\medterm{Teeth} Hard, mineralized structures in the mouth used for cutting, tearing, and grinding food. Hard mineralized structures in the mouth that cut, tear, and grind food, supported by gums and jawbones.
\medterm{Teeth Grinding} Alternate index label; see \textit{Teeth grinding (bruxism)}.

\medterm{Teeth Grinding (Bruxism)} Alternate index label for Bruxism.

\medterm{Teeth, Disorders Of} Conditions affecting teeth, including decay, developmental differences, infection, injury, poor alignment, wear, and tumors. Problems may cause pain, sensitivity, swelling, difficulty chewing, altered appearance, or tooth loss.
\medterm{Teeth-Grinding} Alternate index label; see Teeth grinding (bruxism). Bruxism, repetitive clenching or grinding of teeth, often during sleep and capable of causing tooth wear and jaw pain.
\medterm{Teething} Teething is eruption of an infant’s primary teeth through the gums and may cause gum discomfort, drooling, and chewing; high fever or severe illness suggests another cause.

\medterm{Teichopsia} A visual disturbance involving shimmering zigzags, flashing patterns, or fortress-like shapes. It commonly occurs as a migraine aura but can have other causes, especially when new, persistent, or associated with vision loss.
\medterm{Teicoplanin} A glycopeptide antibiotic used for selected serious infections caused by susceptible Gram-positive bacteria, including some resistant staphylococci. Kidney function and blood levels may need monitoring.
\medterm{Telangiectasia} A visible, permanently dilated small blood vessel near the skin or mucosal surface, occurring alone or with conditions such as rosacea, liver disease, or inherited vascular syndromes.

\medterm{Telangiectasis} Widening of small blood vessels near the skin or a moist body lining, producing visible red or purple lines or patches. Causes include sun exposure, rosacea, hormones, liver disease, and inherited conditions.
\medterm{Tele Radiology (Teleradiology)} The remote review of medical images, such as X-rays, CT scans, or MRI scans, by a qualified radiologist who is not at the imaging site. It can provide urgent or specialist interpretation when local expertise is unavailable and must protect patient information.

\medterm{Tele-Rehabilitation} Remote rehabilitation via technology. Video consultations, remote monitoring, home exercise programs. Increases access, especially rural.
\medterm{Telehealth} Delivery of healthcare through secure electronic communication or remote monitoring, including video visits, telephone care, messaging, and transmission of clinical data.

\medterm{Telemedicine} Delivery of medical care at a distance using information and communication technology. It may include video or telephone consultations, secure messages, remote diagnosis, treatment, and review of clinical information.

\medterm{Telogen Effluvium} A common, non-scarring form of diffuse hair shedding that occurs when more hairs than usual enter the resting (telogen) phase. Shedding often begins 2--4 months after illness, fever, surgery, childbirth, major stress, rapid weight loss, nutritional deficiency, or certain medicines. It is usually temporary, with regrowth after the trigger is corrected, although recovery may take several months.

\medterm{Telomere} A protective region at the end of each chromosome. It helps keep genetic material stable when cells divide and usually becomes shorter with repeated cell division and aging.
\medterm{Temazepam} A benzodiazepine hypnotic used short term for insomnia; sedation, dependence, and respiratory depression are risks.
\medterm{Temperature} A measurement of body heat. The result varies with the measurement site, time, activity, age, and illness.

\medterm{Temperature (Body Temperature)} The balance between heat production (metabolism, muscle activity) and heat loss
(radiation, conduction, convection, evaporation). Normal core body temperature:
36.5--37.5°C (97.7--99.5°F), regulated by the hypothalamus (preoptic anterior
hypothalamus—thermoregulatory center).

\synonyms
Body Temperature

\medterm{Temperature Measurement} Measurement of body or environmental temperature using a thermometer or sensor at a specified site, interpreted with method, age, timing, and clinical context.

\medterm{Temperature Scales} Systems for expressing temperature, including Celsius, Fahrenheit, and Kelvin. Clinical temperatures are commonly reported in Celsius or Fahrenheit.
\medterm{Temple} The region of the skull at the side of the head between the forehead and ear.
\medterm{Temple Syndrome} A genomic-imprinting disorder involving chromosome 14q32 that causes prenatal and
postnatal growth restriction, hypotonia, developmental delay, feeding difficulty, short
stature, and abnormal pubertal development.

\medterm{Temporal} Relating either to time or to the temple area on the side of the skull, including the temporal lobe.
\medterm{Temporal Arteritis} A vasculitic disorder affecting medium and large arteries, particularly the temporal
artery, predominantly occurring in adults over 50 years of age.

\medterm{Temporal Artery} An artery running along the side of the scalp, supplying nearby tissues; inflammation can cause temporal arteritis.
\medterm{Temporal Bones} Paired bones forming the cranial floor and lateral wall, with squamous, petrous,
mastoid, tympanic, and styloid parts. Contain the external auditory canal, middle and
inner ear, and foramina: carotid canal (internal carotid), jugular foramen (CN IX-XI,
IJV), and foramen lacerum.

\medterm{Temporal Lobe} One of the four major lobes of the cerebral hemisphere, located beneath the lateral
fissure (Sylvian fissure) on the lateral surface of the brain, extending medially to the
parahippocampal gyrus. The temporal lobe is involved in auditory processing, language
comprehension, memory, and emotion.

\medterm{Temporal Lobe Epilepsy} The most common form of focal (partial) epilepsy, arising from the medial temporal lobe
structures (hippocampus, amygdala, parahippocampal gyrus).

\medterm{Temporal Lobe Seizure} A seizure that begins in the temporal lobe of the brain and may cause altered awareness, unusual sensations, or repetitive movements.

\synonyms
Seizure, Temporal Lobe

\medterm{Temporary Loss Of Consciousness} Temporary loss of consciousness is a brief episode of unresponsiveness with loss of postural tone, commonly from syncope but sometimes from seizure or another disorder.

\medterm{Temporomandibular Disorders} Alternate index label; see \textit{TMJ disorders}.

\medterm{Temporomandibular Joint} Synovial ginglymoarthrodial joint between mandibular condyle and temporal bone's
mandibular fossa, separated by an articular disc. Only bilateral body joint, allowing
hinge and sliding movements. Disc divides into superior/inferior compartments.
Temporomandibular disorders cause pain, clicking, limited movement, common with bruxism.

\medterm{Temporomandibular Joint Disorder} Pain or impaired movement involving the jaw joint or chewing muscles. It may cause jaw pain, clicking, stiffness, headache, or difficulty opening the mouth. Clenching, grinding, injury, arthritis, and muscle tension can contribute; treatment usually begins with simple pain relief and jaw-rest measures.


\medterm{Temporomandibular Joint Replacement} Surgery that removes a severely damaged jaw joint and replaces it with an artificial joint. It may be considered for ankylosis, advanced joint destruction, or persistent severe symptoms after other treatments have failed.
\medterm{Temporomandibular Joint Syndrome} Pain or dysfunction involving the jaw joint and chewing muscles, often causing jaw pain, clicking, restricted
opening, headache, or tooth-related symptoms. Causes include clenching, trauma, arthritis, and stress-related muscle
tension. Persistent locking, severe pain, or inability to eat requires dental or medical evaluation.

\medterm{TEN} Alternate index label; see \textit{Toxic epidermal necrolysis}.

\medterm{Tenascin} A group of proteins found around cells that help organize tissues and influence growth, healing, and inflammation. Some forms are increased around tumors or damaged tissue.

\medterm{Tender} Painful or sensitive when touched or pressed, a finding that may indicate inflammation, infection, trauma, ischemia, or another underlying process.

\medterm{Tenderness} Pain or discomfort caused by pressing or touching a body part. It may indicate inflammation, injury, infection, or another local or systemic problem.
\medterm{Tendinitis} Tendinitis is painful inflammation or degeneration of a tendon, often from repetitive loading, causing localized tenderness and pain with movement.

\medterm{Tendinitis Rotator Cuff (Rotator Cuff)} Alternate index label for Rotator Cuff.

\medterm{Tendinitis, Achilles} Alternate index label; see \textit{Achilles tendinitis}.

\medterm{Tendinitis, Patellar} Alternate index label; see \textit{Patellar tendinitis}.

\medterm{Tendinopathy} Pain and abnormal structure or function of a tendon, often related to repeated loading, aging, or injury. It can cause tenderness and weakness; treatment usually combines activity adjustment with rehabilitation.
\medterm{Tendon} Dense regular connective tissue connecting muscle to bone, composed primarily of type I
collagen fibers arranged in parallel bundles. Transmits muscle force to bone to produce
movement. Poor vascularity (slower healing). Enclosed in tendon sheath (synovial) at
friction points.

\medterm{Tendon Injury} Damage to a tendon ranging from microscopic overuse injury to partial or complete rupture, causing localized pain, weakness, swelling, or loss of function.

\medterm{Tendon Repair} Surgery to reconnect a torn tendon to itself or to bone so that movement can be restored. It is followed by protected movement and rehabilitation; re-tearing, stiffness, weakness, infection, and nerve or vessel injury can occur.

\medterm{Tendon Sheath} A thin lubricated covering around some tendons that reduces friction as the tendon moves; inflammation is called tenosynovitis.

\medterm{Tendon Transfer} Surgical relocation of a functioning tendon to restore movement or balance when the
original muscle-tendon unit is absent, paralyzed, or irreparable. Successful transfer
requires an adequate donor, appropriate excursion, and a stable recipient joint.

\medterm{Tendonitis} Tendinopathy: inflammation or degeneration of a tendon from overuse, repetitive strain,
or trauma, causing localized pain, swelling, and tenderness with activity. Common sites:
rotator cuff (supraspinatus), lateral elbow (tennis elbow), medial elbow (golfer's
elbow), Achilles, patellar (jumper's knee), biceps, and De Quervain.

\medterm{Tenecteplase} A modified tissue plasminogen activator (TNK-tPA) thrombolytic with longer half-life,
higher fibrin specificity, and resistance to PAI-1. Used for acute ischemic stroke and
STEMI. Given as a single IV bolus (weight-based), simplifying administration. Adverse
effects: hemorrhage (intracranial—major concern), allergic reactions.

\medterm{Tenericutes} A former bacterial group characterized by organisms that lack a typical cell wall, including Mycoplasma and related bacteria.

\medterm{Tenesmus} A distressing sensation of needing to empty the bowel or bladder despite little output, commonly caused by rectal inflammation, infection, tumor, constipation, or urinary irritation.

\medterm{Tenia} Tenia is an older or variant spelling related to Taenia, a genus of tapeworms that can infect humans or animals.


\medterm{Teniposide} Teniposide is a podophyllotoxin-derived antineoplastic drug that inhibits topoisomerase II and is used in selected leukemia or lymphoma regimens.

\medterm{Tennis Elbow} Lateral epicondylitis, a tendinopathy of the common extensor tendon at the lateral
epicondyle (extensor carpi radialis brevis). Caused by repetitive wrist extension and
overuse, not limited to tennis. Presents with lateral elbow pain worsened by wrist
extension and grip. Tenderness over the lateral epicondyle.

\synonyms
Epicondylitis, Lateral
Lateral Elbow Tendinopathy
Lateral Epicondylitis
Lateral Epicondylosis

\medterm{Tennis Elbow (Lateral Epicondylitis)} Alternate index label for Lateral Epicondylitis.

\medterm{Tennis Elbow Surgery} An operation to remove or repair damaged tendon tissue at the outer elbow when persistent tennis-elbow pain has not improved with nonsurgical care. Recovery requires rehabilitation, and pain, stiffness, nerve injury, or continued symptoms are possible.


\medterm{Tenodesis} Surgical fixation of a tendon to bone or another tendon to restore function, stabilize a joint, or treat tendon imbalance.

\medterm{Tenofovir} Tenofovir is a nucleotide reverse-transcriptase inhibitor used with other medicines for HIV and as treatment or prevention for hepatitis B; renal and bone toxicity and formulation-specific monitoring are important.

\medterm{Tenosynovitis} Tenosynovitis is inflammation of the sheath surrounding a tendon, producing pain, swelling, and restricted gliding, commonly in the hand or wrist.

\medterm{Tenosynovitis, De Quervain'S} Alternate index label; see \textit{De Quervain tenosynovitis}.

\medterm{Tenotomy} Surgery that cuts or releases a tendon to reduce tightness, correct a deformity, relieve pressure, or change movement. It may be performed through a small opening or a larger incision; weakness, bleeding, infection, nerve injury, or recurrence can occur.
\medterm{Tenovaginitis} Inflammation of a tendon and its surrounding sheath, also called tenosynovitis.
\medterm{Tensile Strength} The ability of a tissue or material to resist being pulled apart; it is important in wound healing and the integrity of sutures, tendons, and other connective tissues.

\medterm{Tensilon Test} An older pharmacologic test using edrophonium to look for brief improvement in muscle weakness in suspected myasthenia gravis; it is rarely used because safer tests are available.

\medterm{Tension} A pulling force or state of being stretched; in medicine it may describe tissue tension, blood-vessel wall tension, muscle tightness, or psychological strain depending on context.

\medterm{Tension Headache} The most common headache type, characterized by bilateral, pressing/tightening head pain
of mild-moderate intensity (like a band around the head), lasting 30 min to days,
without nausea or vomiting; photophobia or phonophobia may be present but not both.

\medterm{Tension Pneumothorax} Life-threatening pneumothorax with mediastinal shift. Emergency needle decompression
(2nd intercostal space, midclavicular line) followed by chest tube. Signs: respiratory
distress, hypotension, tracheal deviation, distended neck veins, unilateral absent
breath sounds.

\medterm{Tension-Type Headache} A usually bilateral, pressing or tightening headache of mild-to-moderate intensity without prominent nausea, often associated with stress, muscle tenderness, or poor sleep.

\medterm{Tentorium} A strong fold of the brain’s protective covering separating the cerebellum from the back parts of the cerebrum. It supports brain structures and contains an opening through which the brainstem passes.
\medterm{Tenuazonic Acid} A mycotoxin produced by some fungi that contaminates certain plants and foods and can inhibit protein synthesis; toxicity depends on exposure and the producing organism.

\medterm{Teprotumumab} A monoclonal antibody targeting the insulin-like growth factor 1 receptor and used to treat thyroid eye disease in selected patients.

\medterm{Ter In Die} Ter in die means three times a day; explicit plain-language timing is safer than relying on Latin abbreviations.

\medterm{Terameprocol} An experimental anticancer compound studied for effects on gene activity and cell-survival pathways; it is not routine treatment.

\medterm{Teratocarcinoma} A teratocarcinoma is a malignant germ-cell tumor containing tissues from more than one embryonic layer, often arising in the testis or ovary.

\medterm{Teratogen} A medicine, chemical, infection, radiation exposure, or other influence that can harm an embryo or fetus and cause birth defects, growth problems, abnormal development, or pregnancy loss. Risk depends on the exposure, amount, timing, and individual factors; exposure should be discussed promptly with a clinician.
\medterm{Teratogenesis} The development of a birth defect, impaired growth, abnormal function, or pregnancy loss after harmful exposure during pregnancy. Risk depends on the substance or infection, amount, timing, route, and genetic factors; exposure does not automatically mean that harm will occur.
\medterm{Teratogenic} Capable of causing congenital malformations, growth impairment, functional deficits, or pregnancy loss when an embryo or fetus is exposed during a vulnerable developmental period.

\medterm{Teratogenic Drugs} Medicines that can cause birth defects, pregnancy loss, or impaired fetal development; risk depends on the drug, dose, timing, and pregnancy.

\medterm{Teratogenicity} The ability of a drug, chemical, or infectious agent to cause birth defects or
developmental abnormalities when exposure occurs during pregnancy. The risk of
teratogenic effects depends on the timing of exposure, with the embryonic period from 3
to 8 weeks after fertilization being the most critical for organogenesis.

\medterm{Teratoma} A type of germ cell tumor composed of tissues derived from two or three germ layers
(ectoderm, endoderm, mesoderm), reflecting pluripotent differentiation.

\medterm{Teratozoospermia} Teratozoospermia is a high proportion of abnormally shaped sperm and may contribute to reduced fertility, though semen results must be interpreted as a whole.

\medterm{Terazosin Hydrochloride} The hydrochloride salt of terazosin, an alpha-1 blocker used for urinary symptoms from benign prostatic enlargement and sometimes for hypertension; dizziness and first-dose fainting can occur.

\medterm{Terbinafine} An allylamine antifungal agent that inhibits squalene epoxidase, an early enzyme in the
ergosterol biosynthesis pathway, causing accumulation of squalene and depletion of
ergosterol in the fungal cell membrane. The accumulation of squalene is fungicidal,
while ergosterol depletion disrupts membrane function.

\medterm{Terbufos} Terbufos is a highly toxic organophosphate pesticide that inhibits acetylcholinesterase and can cause severe cholinergic poisoning.

\medterm{Terbutaline} A beta-2 agonist used as a bronchodilator; it has limited and specialist use as a tocolytic.
\medterm{Terfenadine} An older antihistamine withdrawn or restricted in many places because interactions and excessive blood levels can cause dangerous heart-rhythm abnormalities.

\medterm{Teriparatide For Osteoporosis} Recombinant human parathyroid hormone fragments one through thirty-four that stimulate
osteoblast-mediated bone formation. Indicated for severe osteoporosis with fractures or
very low T-scores and for patients who have failed or are intolerant to bisphosphonates.
Daily subcutaneous injection for up to two years.

\medterm{Terlipressin} Terlipressin is a vasopressin analogue that constricts splanchnic vessels and is used in selected cases of variceal bleeding or hepatorenal syndrome.

\medterm{Term} The end of pregnancy at 37 completed weeks of gestation; births before this are classified as preterm.

\medterm{Term Birth} Term birth occurs at or after 37 completed weeks of pregnancy; early term, full term, and late or post-term categories describe more specific gestational ranges.

\medterm{Terminal Care} Care focused on comfort, symptom relief, dignity, and support for a person with an advanced illness who is approaching the end of life.

\medterm{Terminal Delirium} Delirium occurring near the end of life, often caused by advanced organ failure,
medication effects, metabolic disturbance, infection, or reduced cerebral perfusion.
Management prioritizes comfort, reversible contributors when consistent with goals, and
support for the patient and family.

\medterm{Terminal Ileitis} Inflammation of the distal ileum, commonly associated with Crohn disease but also caused by infection, ischemia, medication, or other inflammatory conditions and producing pain, diarrhea, or bleeding.

\medterm{Terminal Illness} An advanced disease that is expected to lead to death, although the timing and response to available treatment may vary.

\medterm{Terminal Insomnia} Waking too early and being unable to return to sleep, a pattern that can occur in depression, circadian-rhythm disorders, aging, or other insomnia conditions.

\medterm{Terms Medical Abbreviations (Common Medical Abbreviations And Terms)} Alternate index label for Common Medical Abbreviations and Terms.

\medterm{Terpene} A terpene is a hydrocarbon built from isoprene units and found in plant oils, fragrances, medicines, and other natural products.

\medterm{Terry'S Nails} A nail change in which most of the nail plate appears white with a narrow pink or brown distal band; it may be associated with liver disease, heart failure, diabetes, or ageing.

\medterm{Tertian Fever} A fever pattern recurring approximately every 48 hours, classically associated with Plasmodium vivax or P. falciparum malaria.
\medterm{Tertiary Hyperparathyroidism} Autonomous overproduction of parathyroid hormone after long-standing secondary hyperparathyroidism, usually in chronic kidney disease, causing high calcium and high PTH.

\medterm{Tertiary Prevention} Care that reduces complications, disability, or suffering from an established disease. Examples include rehabilitation, long-term disease management, preventing relapses, controlling symptoms, restoring function, and providing palliative care.

\medterm{Tertiary Protein} The three-dimensional folding of a single protein chain, created by interactions among amino-acid side chains and determining the protein’s functional shape.

\medterm{Tesetaxel} Tesetaxel is an investigational oral taxane that disrupts microtubules and has been studied as an anticancer treatment.

\medterm{Test Meal} A standardized meal given before measuring a physiologic response, such as gastric acid or glucose handling.
\medterm{Test Strips} Narrow strips containing chemicals that react with a sample to measure or detect a substance. Common examples test urine for blood or infection, or test blood for glucose or ketones. Results can be affected by storage, timing, contamination, and incorrect use.

\medterm{Test-Tube} A small glass or plastic laboratory container used to hold, mix, heat, or test samples. Test tubes may contain blood, urine, chemicals, cells, cultures, or other materials during medical investigations.
\medterm{Test-Tube Baby} A lay term for a baby conceived through in-vitro fertilization.
\medterm{Testes} The paired male reproductive organs in the scrotum. They produce sperm and testosterone and may be affected by infection, injury, torsion, tumors, or disorders of development.
\medterm{Testicle} One of the male reproductive glands in the scrotum. It produces sperm and testosterone; sudden severe testicular pain is an emergency because it may indicate torsion.
\medterm{Testicle Lump} A localized mass or enlargement in or around a testis caused by cyst, infection, inflammation, torsion, trauma, or tumor; a new intratesticular mass requires prompt evaluation.

\medterm{Testicle Pain} Pain in one or both testes or the surrounding structures, with urgent causes including torsion, infection, trauma, and incarcerated hernia.

\medterm{Testicle, Diseases Of} Conditions affecting the testis, including twisting, infection, injury, tumors, blood-vessel problems, and abnormal development. Symptoms may include pain, swelling, a lump, fertility problems, or changes in hormone production.
\medterm{Testicle, Retractile} Alternate index label; see \textit{Retractile testicle}.

\medterm{Testicle, Undescended} Alternate index label; see \textit{Undescended testicle}.

\medterm{Testicles} The paired male reproductive glands in the scrotum that produce sperm and testosterone; pain, enlargement, a new lump, or sudden change requires examination.

\medterm{Testicular Atrophy} Shrinkage and loss of functional tissue in a testicle caused by undescended testis, infection, torsion, trauma, varicocele, hormones, aging, or other injury.

\medterm{Testicular Biopsy} Removal of a small sample of testicular tissue for microscopic examination, used selectively to evaluate infertility, atypical lesions, or suspected testicular disease.

\medterm{Testicular Cancer} Cancer arising in a testicle, most often in younger adult men. It commonly appears as a painless lump, enlargement, heaviness, or ache. Ultrasound and blood tests help assess a suspected tumor; treatment is highly effective for many people, especially when found early.

\synonyms
Cancer, Testicular

\medterm{Testicular Disorders} Conditions affecting the testes or surrounding structures, including torsion, infection, inflammation, trauma, tumors,
and hormonal disorders. Symptoms may include pain, swelling, a mass, infertility, or
changes in sexual development and require appropriate evaluation.

\synonyms
Infection Testicle (Testicular Disorders)
Pain Scrotum (Testicular Disorders)
Torsion Testicle (Testicular Disorders)
Tumor Testicle (Testicular Disorders)



\medterm{Testicular Gonadoblastoma} A rare mixed germ-cell and sex-cord tumor arising in a dysgenetic or undescended gonad, often associated with Y-chromosome material and risk of another germ-cell tumor.

\medterm{Testicular Hydrocele} A collection of clear fluid around a testicle within the tunica vaginalis, causing painless scrotal enlargement that transilluminates; it may be congenital, reactive, or secondary to disease.

\medterm{Testicular Mass} A lump or enlargement within or around a testicle caused by tumor, infection, inflammation, cyst, hematoma, or torsion; a solid intratesticular mass requires urgent ultrasound and specialist assessment.

\medterm{Testicular Microlithiasis} Multiple tiny calcifications within the testicular tissue seen on ultrasound. The finding alone does not establish cancer, and follow-up depends on other risk factors and clinical findings.

\medterm{Testicular Self-Exam} A testicular self-exam is a gentle check for new lumps, swelling, heaviness, or changes in size or shape; a new abnormality should be medically assessed.

\medterm{Testicular Torsion} Surgical emergency caused by twisting of the spermatic cord around its longitudinal
axis, cutting off venous return and leading to arterial occlusion, ischemia, and
infarction of the testicle.

\synonyms
Torsion, Testicular

\medterm{Testicular Torsion Emergency} Rotation of the spermatic cord that compromises testicular blood flow, causing sudden
unilateral scrotal pain, nausea, a high-riding testis, and an abnormal cremasteric
reflex. Immediate surgical exploration and detorsion are required; imaging must not
delay treatment when suspicion is high.

\medterm{Testicular Torsion Repair} Urgent surgery to untwist a spermatic cord and restore testicular blood flow, followed by fixation of the affected and usually opposite testis to prevent recurrence.

\medterm{Testis} The male reproductive gland that produces spermatozoa in seminiferous tubules and
secretes testosterone from interstitial Leydig cells. It is normally located in the
scrotum.

\medterm{Testis (Testicle)} The male gonad, a paired ovoid organ (4--5 cm \(\times\) 2--3 cm, weight 15--20 g) located in
the scrotum, responsible for sperm production (spermatogenesis) and testosterone
secretion. Each testis is covered by the tunica vaginalis (peritoneal remnant) and
tunica albuginea (dense fibrous capsule).

\synonyms
Testicle

\medterm{Testis Disorder} A condition affecting a testicle, including torsion, infection, tumor, trauma, infertility-related damage, hormonal dysfunction, or abnormal development.

\medterm{Testis Neoplasm} A tumor arising in testicular tissue, most commonly a germ-cell tumor in younger adults; painless enlargement, heaviness, or a solid mass requires prompt ultrasound and evaluation.

\medterm{Testolactone} An aromatase inhibitor formerly used to reduce estrogen production in selected endocrine conditions; current use is limited by availability and newer alternatives.

\medterm{Testosterone} A steroid sex hormone produced mainly by the testes and, in smaller amounts, by the ovaries and adrenal glands. It supports reproductive development, muscle and bone maintenance, red-blood-cell production, libido, and other physiologic functions; levels vary with age, sex, time of day, illness, and medicines.

\medterm{Testosterone Deficiency} Alternate index label; see \textit{Male hypogonadism}.

\medterm{Testosterone Low (Low Testosterone (Low T))} Alternate index label for Low Testosterone (Low T).

\medterm{Testosterone Propionate} A short-acting injectable testosterone ester used in selected hormone-replacement or androgen-deficiency settings; risks include erythrocytosis, infertility, acne, mood effects, and prostate-related monitoring needs.

\medterm{Tests And Visits Before Surgery} Preoperative tests and visits review health conditions, medicines, anesthesia risk, needed laboratory or imaging tests, and instructions for the operation.

\medterm{Tests For H Pylori} Tests for Helicobacter pylori infection, including urea breath testing, stool antigen testing, biopsy-based testing, or selected blood tests; choice depends on treatment and clinical context.

\medterm{Tet Spell} A sudden episode of severe bluish coloring and low oxygen in a child with unrepaired tetralogy of Fallot. Crying, feeding, or exertion may trigger it; the child may breathe rapidly, become limp, or faint. It is an emergency requiring calming, positioning, oxygen, and urgent specialist care.

\synonyms
Hypercyanotic Spell; Paroxysmal Dyspneic Attack; Blue Spell

\medterm{Tetanus} A serious infection caused by bacteria entering a wound and producing a toxin that causes severe muscle stiffness and painful spasms. It can cause lockjaw, breathing problems, and death without urgent treatment.

\synonyms
Lockjaw

\medterm{Tetanus Toxin} A neurotoxin made by the bacterium Clostridium tetani. It blocks inhibitory signals in the nervous system, causing painful muscle stiffness and spasms, including lockjaw; vaccination and urgent wound care help prevent tetanus.
\medterm{Tetany} A clinical syndrome of neuromuscular hyperexcitability caused by hypocalcaemia (most
common—total calcium <2.1 mmol/L, ionised calcium <1.1 mmol/L), but also caused by
hypomagnesaemia, hypokalaemia, and alkalosis (respiratory or metabolic).

\medterm{Tethered Cord Syndrome} A congenital or acquired condition where the spinal cord is abnormally attached
(tethered) to surrounding tissues, restricting its normal upward movement within the
spinal canal. Caused by a thickened filum terminale, lipomeningomyelocele,
diastematomyelia, or previous myelomeningocele repair.

\medterm{Tetrabenazine} A medicine that reduces storage and release of certain brain chemicals involved in movement. It treats selected involuntary movement disorders but may cause sleepiness, depression, restlessness, stiffness, swallowing problems, or abnormal movements.
\medterm{Tetracaine} Tetracaine is a potent ester local anaesthetic used topically or for spinal anaesthesia; excessive dose can cause seizures, cardiovascular toxicity, and methemoglobinaemia.

\medterm{Tetrachloroethylene} Tetrachloroethylene, or perchloroethylene, is an industrial solvent used in dry cleaning; substantial exposure can affect the nervous system, liver, and kidneys.

\medterm{Tetrachlorvinphos} An organophosphate insecticide that can block acetylcholinesterase. Significant exposure may cause sweating, vomiting, pinpoint pupils, muscle twitching, slow breathing, seizures, or coma and requires urgent treatment.
\medterm{Tetracycline} A medicine, biologic, or pharmacologic term related to tetracycline; its mechanism, indication, interactions, and adverse effects depend on the specific agent.

\medterm{Tetracycline Stain} Intrinsic discoloration from tetracycline antibiotics during tooth development. Gray or
brown bands. Treatment: bleaching, veneers, crowns.

\medterm{Tetracyclines} A class of bacteriostatic antibiotics that bind reversibly to the 30S ribosomal subunit
and block the attachment of aminoacyl-tRNA to the acceptor site, inhibiting protein
synthesis. They include tetracycline, doxycycline, and minocycline.

\medterm{Tetrahydrozoline Poisoning} Poisoning from swallowing eye drops or nasal drops containing tetrahydrozoline. It can cause extreme sleepiness, slow heart rate, low blood pressure, and breathing problems, especially in children.

\medterm{Tetralogy Of Fallot} A congenital heart defect comprising ventricular septal defect, right-ventricular outflow obstruction, overriding aorta, and right-ventricular hypertrophy.
\medterm{Tetraplegia} Paralysis or severe weakness of all four limbs, usually from cervical spinal-cord disease or injury.
\medterm{Tetraplegia (Quadriplegia)} Paralysis or severe weakness affecting all four limbs and usually the trunk, commonly caused by cervical spinal-cord injury or disease.

\synonyms
Quadriplegia

\medterm{TGA} Alternate index label; see \textit{Transient global amnesia}.

\medterm{Th1/Th2/Th17/Tfh/Treg Balance} A simplified description of how different helper T-cell groups coordinate immune responses. Some promote defense against infections, others support antibody production or inflammation, and regulatory cells limit excessive immune activity. Imbalance may contribute to allergy, autoimmune disease, or persistent infection.

\medterm{Thalamus} Bilateral ovoid gray matter structure forming the third ventricle's lateral walls, the
brain's relay station for nearly all sensory (except olfaction), motor, and limbic
information to the cortex. Numerous nuclei: VPL (spinothalamic, dorsal column), VPM
(trigeminal), LGN (visual), MGN (auditory), VL/VA (motor).

\medterm{Thalassaemia} The British spelling of thalassemia, inherited reduced production of alpha- or beta-globin chains causing anemia.
\medterm{Thalassemia} A group of inherited hemoglobin disorders caused by reduced or absent synthesis of one
or more globin chains, leading to microcytic hypochromic anemia.

\synonyms
Cooley's Anemia
Mediterranean Anemia

\medterm{Thalassemia Major} A severe inherited blood disorder in which the body makes too little hemoglobin. It causes serious anemia beginning in childhood and may require regular blood transfusions, with monitoring for iron buildup and organ damage.

\medterm{Beta-Thalassemia Trait} The heterozygous carrier state for a beta-globin gene variant. It is usually asymptomatic or causes mild microcytosis with mild or no anemia; hemoglobin electrophoresis often shows increased HbA2. It is also called beta-thalassemia minor.

\synonyms
Beta Thalassemia Trait
Beta-Thalassemia Minor
Thalassemia Minor

\medterm{Thalidomide} A medicine that changes immune activity and is used for selected cancers and inflammatory conditions. It can cause severe birth defects, nerve damage, blood clots, sleepiness, and other harms, requiring strict safety controls.
\medterm{Thallium} Thallium is a toxic heavy metal that can cause gastrointestinal illness followed by painful peripheral neuropathy, hair loss, skin changes, kidney or liver injury, seizures, and cardiovascular or neurologic toxicity. Suspected exposure requires immediate poison-control or emergency guidance.

\synonyms
Poisoning Thallium (Thallium)

\medterm{Thallium Scan} A nuclear-medicine test using thallium-201 to show blood flow to heart muscle or selected tumors. It is now less common because other tracers are often preferred, but may still have specific uses.

\medterm{Thallium-201} SPECT radionuclide, half-life 73 hours. Used for myocardial perfusion imaging.
Mechanism: K+ analog, taken up by Na+/K+ ATPase.

\medterm{Thanatology} The scientific study of death, dying, postmortem changes, and the psychological, social, and ethical aspects
of bereavement. In forensic medicine, it also examines the circumstances and evidence
surrounding death.

\medterm{Thanatophoric Dysplasia} A severe skeletal dysplasia, usually caused by an FGFR3 variant, with markedly shortened
limbs, narrow thorax, macrocephaly, and characteristic skull and vertebral
abnormalities. The severe thoracic restriction commonly causes neonatal respiratory
failure.





\medterm{Theca} A connective-tissue sheath, especially the theca surrounding an ovarian follicle.
\medterm{Theca Lutein Cyst} A usually temporary ovarian cyst caused by strong stimulation from pregnancy-related hormones or fertility treatment. It is often present in both ovaries and may cause abdominal pressure or pain; it commonly shrinks when hormone levels return to normal.

\medterm{Thecoma} A usually benign ovarian sex-cord stromal tumor composed of theca-like cells, often producing estrogen and presenting with abnormal uterine bleeding or an incidental adnexal mass.
\medterm{Thelarche} The onset of breast development, usually the first visible sign of female puberty. It commonly begins between ages 8 and 13, with substantial normal variation.

\medterm{Thenar Eminence} The fleshy muscular prominence at the base of the thumb.
\medterm{Theophylline} A medicine that can relax airway muscles and improve breathing in selected lung diseases. Its safe and harmful doses are close together, so blood levels, medicine interactions, heart rhythm, and side effects need monitoring.
\medterm{Theranostic Pair} A matched diagnostic and treatment approach using related agents that target the same biological feature. One agent images the target and the other delivers treatment, often radiation.
theranostics.

\medterm{Theranostics} Combined therapy and diagnostics using matched radiopharmaceutical pairs. Diagnostic
isotope (e.g., Ga-68) identifies target, therapeutic isotope (e.g., Lu-177) delivers
treatment. Examples: Ga-68/Lu-177 DOTATATE (neuroendocrine tumors), PSMA theranostics
(prostate cancer).


\medterm{Therapeutic Abortion} Intentional termination of a pregnancy because of medical, fetal, or other legally recognized indications, according to applicable law and clinical guidance.

\medterm{Therapeutic Alliance} The collaborative, trusting relationship between therapist and patient that is
foundational to effective psychotherapy. It includes agreement on goals, tasks, and the
bond between therapist and patient. It is the strongest predictor of psychotherapy
outcome across all therapeutic modalities.

\medterm{Therapeutic Cloning} Creating a genetically matched embryonic stem-cell line for research or potential tissue replacement; it is distinct from reproductive cloning and remains ethically and technically limited.

\medterm{Therapeutic Exercise} Planned physical activity used to improve strength, flexibility, balance, movement, endurance, or daily function and to reduce pain or disability. A clinician selects and gradually adjusts exercises according to the person’s condition, abilities, symptoms, and goals.

\medterm{Therapeutic Hypothermia} Deliberate, closely monitored cooling of the body to reduce tissue damage in selected emergencies, such as after some cardiac arrests or in severe overheating-related illness. Cooling is used with treatment of the underlying cause and can cause abnormal heart rhythms, shivering, or electrolyte changes.

\medterm{Therapeutic Index} The therapeutic index is a quantitative measure of the relative safety of a drug,
defined as the ratio of the dose that produces toxicity in 50 percent of the population
to the dose that produces the desired therapeutic effect in 50 percent of the
population.

\medterm{Therapeutic Window} The range of plasma drug concentrations between the minimum effective concentration
(producing a therapeutic effect) and the minimum toxic concentration (producing
toxicity), in which efficacy is achieved with acceptable toxicity.

\medterm{Therapeutics} The branch of medicine concerned with the treatment of disease, encompassing all
interventions used to prevent, manage, or cure illness.

\medterm{Therapy} Treatment intended to prevent, relieve, or cure disease or to improve function and quality of life.

\medterm{Therapy Monitoring} Repeated clinical assessment, laboratory testing, or imaging used to evaluate response, adverse effects, and ongoing need for treatment.

\medterm{Thermal Injury} Tissue damage caused by heat, flame, hot liquid, steam, or cold. Burn severity depends
on depth, total body surface area, location, inhalation injury, associated trauma, and
patient factors; assessment should include airway and circulatory complications.

\medterm{Thermal Injury (Burns)} Tissue damage caused by heat. Initial care includes stopping the exposure, removing loose
constricting material, cooling briefly with clean water, covering the injury, and assessing
depth and extent. Larger or deeper burns require specialized evaluation, fluid management,
pain control, wound care, and tetanus assessment.

\synonyms
Burns

\medterm{Thermal Therapy} Treatment using heat or cold to influence pain, swelling, muscle tension, or tissue temperature. Effects and safety depend on the condition and method used.

\medterm{Thermogenesis} The production of heat by the body. It comes mainly from normal metabolism, muscle activity such as shivering, digestion, and specialized fat tissue that releases stored energy as heat.
\medterm{Thermography} Imaging or recording of heat patterns from the body surface.
\medterm{Thermoluminescent Dosimeter} A radiation-monitoring device that measures accumulated exposure by light released from an irradiated crystal when heated.
\medterm{Thermometer} A device used to measure body temperature. Common types measure temperature in the mouth, ear, armpit, or across the forehead. Readings vary with the site and method, so the measurement site and scale should be recorded.

\medterm{Thermometer Scales} Numbering systems used to report temperature, such as Celsius, Fahrenheit, and Kelvin. The scale should be identified when a temperature is communicated.
\medterm{Thermoneutral Zone} The range of ambient temperatures over which body temperature is maintained by basal
metabolic heat production alone, without additional thermoregulatory responses such as
shivering or sweating. In humans the zone is approximately 27-32 degrees C for a resting
nude adult. Within the zone, vasomotor tone adjusts heat loss.

\medterm{Thermoplasty} A bronchoscopic procedure that uses radiofrequency energy to ablate smooth muscle in the
airways, reducing bronchospasm. It is approved for severe, persistent asthma in adults
not controlled with ICS/LABA. It reduces exacerbations and emergency department visits
but does not improve lung function.

\medterm{Thermoreceptor} A nerve ending that senses warmth or cold in the skin and other tissues. Signals travel through the spinal cord to the brain, helping a person recognize temperature and helping the body control heat through sweating, shivering, and changes in skin blood flow. Damage to these nerves can reduce temperature sensation and increase the risk of burns or cold injury.
\medterm{Thermoregulation} The physiological maintenance of core body temperature near the set point (~37 degrees
C) via a balance of heat production and heat loss, controlled by the preoptic area of
the hypothalamus.

\medterm{Thermostat} A thermostat regulates temperature by sensing heat and controlling a heating or cooling system; accurate environmental temperature can matter for patients, specimens, and medicines.

\medterm{Thiabendazole} An antiparasitic medicine used for selected worm infections and, historically, some fungal infections. It can cause nausea, dizziness, liver problems, or nervous-system effects, so safer alternatives are often preferred.
\medterm{Thiamin} Vitamin B1, a water-soluble vitamin required for carbohydrate metabolism and neural function; severe deficiency can cause beriberi or Wernicke encephalopathy.

\medterm{Thiamine} Vitamin B1, a water-soluble vitamin required for carbohydrate metabolism and nerve function. Severe deficiency
can cause beriberi or Wernicke encephalopathy.

\medterm{Thiamine (Vitamin B1)} A water-soluble vitamin whose active coenzyme form, thiamine pyrophosphate, is required
by pyruvate dehydrogenase, alpha-ketoglutarate dehydrogenase, transketolase, and
branched-chain ketoacid dehydrogenase. Deficiency causes beriberi and Wernicke
encephalopathy.

\synonyms
Vitamin B1

\medterm{Thiazide Diuretics} Thiazide diuretics act on the distal convoluted tubule by inhibiting the sodium-chloride
cotransporter, reducing sodium and chloride reabsorption and producing natriuresis and
diuresis. They are less potent than loop diuretics but have a longer duration of action,
typically 6 to 12 hours.

\medterm{Thiazide Overdose} Taking too much of a thiazide diuretic can cause dehydration, low blood pressure, abnormal sodium or potassium levels, confusion, or abnormal heart rhythms.

\medterm{Thiazides} Diuretics that inhibit sodium-chloride reabsorption in the distal convoluted tubule.
\medterm{Thiazolidinedione Drugs} A group of diabetes medicines that make body tissues more sensitive to insulin by activating a protein called PPAR-gamma. Pioglitazone is the main example; possible effects include weight gain, swelling, bone fractures, and worsening heart failure.
\medterm{Thiazolidinediones} Insulin sensitizers activating peroxisome proliferator-activated receptors primarily in
adipose tissue and skeletal muscle. Agents include pioglitazone and rosiglitazone.
Mechanisms include enhanced insulin-mediated glucose uptake, altered adipokine
secretion, and reduced hepatic glucose output.

\medterm{Thigh} The part of the lower limb between the hip and knee. It contains the femur, large muscles, major blood vessels, and nerves that help move and support the leg.
\medterm{ThinPrep Pap Test} A cervical screening test in which cells collected from the cervix are placed in a liquid vial and examined for precancerous changes, cancer, or sometimes HPV-related abnormalities.
\synonyms
Thin Prep Pap Screen

\medterm{Thiopental Sodium} An ultra-short-acting barbiturate, used for induction of general anesthesia (IV, onset
30--60 seconds, duration 5--10 minutes) and as a third-line agent for raised
intracranial pressure (cerebroprotective effect: reduces cerebral metabolism, cerebral
blood flow, and intracranial pressure).

\medterm{Thiopentone Sodium} A short-acting barbiturate anesthetic used for induction in selected settings.
\medterm{Thioridazine} An older phenothiazine antipsychotic with important cardiac and neurologic toxicity.
\medterm{Thioridazine Overdose} Poisoning from taking too much thioridazine, an antipsychotic medicine. It can cause severe sleepiness, low blood pressure, seizures, and dangerous heart rhythms.

\medterm{Thiotepa} A chemotherapy medicine that damages DNA and stops cells from multiplying. It is used for selected cancers and sometimes before a stem-cell transplant; side effects may include low blood counts, infection, nausea, mouth sores, and organ injury.
\medterm{Third Stage Of Labor} From delivery of the baby to delivery of the placenta. The placenta separates from the
uterine wall and is expelled spontaneously or with gentle cord traction. Normal
duration: 5-30 minutes.

\medterm{Third Trimester} The period of pregnancy from 28 weeks until birth. The fetus grows rapidly and continues developing its lungs, brain, and other organs. Prenatal visits monitor blood pressure, growth, movement, and symptoms; tests and screening are scheduled according to local guidance and individual risk.

\medterm{Third-Degree AV Block} Complete heart block with no conduction between atria and ventricles. Atria and
ventricles beat independently (AV dissociation). Ventricular escape rhythm: junctional
(40-60, narrow QRS) or ventricular (20-40, wide QRS).

\medterm{Third-Degree Burn} A full-thickness burn that destroys the epidermis and dermis and may extend into subcutaneous tissue.
The skin may appear white, brown, or charred and sensation may be reduced. Assessment
includes burn depth, extent, inhalation injury, and associated trauma; deep or extensive
burns require specialist treatment and may need grafting.

\medterm{Third-Degree Perineal Tear} A severe tear during vaginal birth that extends from the vaginal opening into the muscles controlling the anus. It needs prompt examination and careful repair, followed by bowel-softening treatment and follow-up; pain, infection, difficulty controlling stool, or a fistula can result.

\synonyms
Third-Degree Tear; Obstetric Anal Sphincter Tear

\medterm{Thirst} The conscious desire to drink, regulated mainly by body-fluid concentration and circulating blood volume. Excessive thirst may occur with dehydration, diabetes, medicines, kidney problems, or other conditions and may need evaluation.
\medterm{Thirst - Absent} Reduced or absent desire to drink despite a potential need for fluids, which can occur with hypothalamic disease, aging, altered consciousness, medications, or severe illness.

\medterm{Thirst - Excessive} Polydipsia, an abnormally strong or persistent desire to drink, commonly associated with dehydration, diabetes mellitus, diabetes insipidus, medicines, or psychiatric conditions.

\medterm{Thomsen Disease} An inherited muscle disorder causing stiffness because muscles relax slowly after they contract. Symptoms often begin in childhood, affect the hands, legs, eyelids, or jaw, and improve after repeated movement. Strength and life expectancy are usually preserved, although stiffness can interfere with activity.

\synonyms
Autosomal Dominant Myotonia Congenita; Myotonia Congenita (Thomsen Type)

\medterm{Thoracentesis} Insertion of a needle or catheter into the pleural space to remove fluid for diagnosis or to relieve respiratory symptoms caused by a pleural effusion.

\medterm{Thoracic} Thoracic means relating to the chest, the thoracic spine, or structures within the thorax such as the heart and lungs.

\medterm{Thoracic Aortic Aneurysm} A weakened, widened area of the aorta within the chest. It may cause no symptoms but can press on nearby structures or rupture, with risks increased by hypertension, inherited disorders, or a bicuspid valve.

\medterm{Thoracic Cage} The bony and cartilaginous framework formed by the thoracic vertebrae, ribs, costal
cartilages, and sternum. It protects thoracic organs and assists respiratory movements.

\medterm{Thoracic Cavity} The chest cavity, enclosed by the ribs and chest muscles, containing the lungs, heart,
esophagus, and trachea.

\medterm{Thoracic Duct} Body's largest lymphatic vessel, approximately 40 cm long, draining lymph from the
entire body below the diaphragm and left side above it (approximately 75 percent of
total lymph).

\medterm{Thoracic Endometriosis} Endometrial tissue within the thoracic cavity causing catamenial pneumothorax,
hemothorax, or lung nodules, recurring cyclically with menstruation and often requiring
hormonal suppression.

\medterm{Thoracic Outlet Syndrome} Pressure on nerves or blood vessels as they pass from the neck into the arm. It may cause neck, shoulder, or arm pain, tingling, weakness, swelling, color change, or a cold hand. Treatment depends on the cause and may include physical therapy or surgery.

\synonyms
TOS (Thoracic Outlet Syndrome)

\medterm{Thoracic Spine CT Scan} Computed tomography of the thoracic spine used to evaluate vertebral fractures, alignment, bone lesions, spinal canal narrowing, and other bony abnormalities.

\medterm{Thoracic Spine X-Ray} Plain radiographic imaging of the thoracic spine used to assess vertebral alignment, fractures, curvature, degenerative change, and selected bone lesions.

\medterm{Thoracic Surgeon} Surgeon specializing in operations on the chest, including lungs, heart, esophagus, and mediastinum. A surgeon who performs operations on chest organs, including the lungs, esophagus, heart, and mediastinum.
\medterm{Thoracic Vertebrae} Have costal facets on bodies and transverse processes for rib articulation. Heart-shaped
bodies, long downward spinous processes, coronally-oriented facets allowing rotation and
lateral flexion but limiting flexion/extension. Small circular vertebral foramen.

\medterm{Thoracic Vertebrae (T1-T12)} The twelve vertebrae of the chest region, each articulating with ribs and protecting the thoracic spinal cord while supporting the trunk and rib cage.
\medterm{Thoracolumbar Spine Trauma} Injury to the thoracic or lumbar spine. Assessment includes fracture morphology,
posterior ligamentous complex integrity, and neurologic status. Classification systems such
as AO Spine and TLICS help guide management, which depends on stability and the clinical
context.

\medterm{Thoracoscopy} A minimally invasive procedure using a camera inserted through the chest wall to examine or treat the pleura, lungs, or mediastinum.

\medterm{Thoracotomy} A surgical incision into the chest wall to access the thoracic cavity (pleural space,
lungs, heart, great vessels, esophagus, diaphragm, mediastinum).

\medterm{Thorax} The thorax is the chest region between the neck and abdomen, enclosed by ribs and containing the heart, lungs, major vessels, and mediastinal structures.

\medterm{Thought Blocking} An abrupt interruption in the train of thought before it is completed, as if the thought
was removed or pushed out of the patient's mind. The patient may pause mid-sentence,
appear confused, and report that the thought was ``taken out'' or ``blocked.'' It is a
formal thought disorder seen in schizophrenia.

\medterm{Thought Disorders} Disturbances of the form, flow, or content of thinking, occurring in psychosis and other conditions.
\medterm{Threadworm} An intestinal infection or the worm Enterobius vermicularis, also called pinworm.
\medterm{Threadworm (Enterobiasis)} An intestinal parasitic infection caused by Enterobius vermicularis, also called
pinworm. Infection spreads when eggs are swallowed, often after contact with contaminated hands, bedding, food, or
surfaces.

\synonyms
Enterobiasis

\medterm{Threadworms (Pinworms)} Small intestinal worms, Enterobius vermicularis, whose females lay eggs around the anus and commonly cause nocturnal itching.

\synonyms
Pinworms

\medterm{Threatened Miscarriage} Vaginal bleeding during early pregnancy with a closed cervix and no passage of pregnancy tissue; the pregnancy may remain viable.
Evaluation commonly includes clinical assessment, serial human chorionic gonadotropin testing when indicated, and ultrasound.

\synonyms
Threatened abortion.

\medterm{Three-Day Measles} Alternate index label; see \textit{Rubella}.

\medterm{Three-Dimensional Ultrasound} Ultrasound that combines many sound-wave measurements to create a three-dimensional image. It is used in obstetrics, gynecology, and other settings to show anatomy from several angles.

\medterm{Threonine} An essential amino acid used to build proteins and other body molecules. The body cannot make enough of it, so it must come from food.
\medterm{Threshold} The lowest amount of a stimulus or measurement needed to produce a detectable response or effect.
\medterm{Thrill} A palpable vibration over the chest or a blood vessel caused by turbulent blood flow.
\medterm{Throat} The pharynx and adjacent structures behind the mouth and nose.
\medterm{Throat Cancer} Cancer arising in the pharynx or nearby structures of the throat. Persistent hoarseness, difficulty or pain
with swallowing, a neck lump, unexplained weight loss, or persistent sore throat requires
medical evaluation, especially in people who use tobacco or alcohol.

\synonyms
Cancer, Throat

\medterm{Throat Cancer (Larynx Cancer)} Alternate index label for Larynx Cancer.

\medterm{Throat Or Larynx Cancer} Throat or larynx cancer is malignant growth in the pharynx or voice box and may cause persistent hoarseness, swallowing difficulty, pain, or a neck lump.

\medterm{Throat Swab Culture} Culture of a throat specimen to detect bacterial pathogens, especially group A Streptococcus, when evaluating pharyngitis.

\synonyms
Throat Culture

\medterm{Throbbing} A pulsating or beating quality of pain or another sensation. It can occur with migraine, inflammation, injury, infection, or increased awareness of a pulse.
\medterm{Thrombectomy} A surgical or interventional radiology procedure to remove a blood clot (thrombus) from
a blood vessel, restoring blood flow. Types: (1) Endovascular thrombectomy (mechanical
thrombectomy)—for acute ischemic stroke due to large vessel occlusion (LVO—internal
carotid artery, middle cerebral artery M1 segment).

\medterm{Thrombectomy (Vascular)} Removal of a blood clot from a blood vessel, using an open surgical or catheter-based procedure.

\synonyms
Vascular

\medterm{Thrombin} A blood-clotting substance that changes fibrinogen into fibrin, the protein threads that form a clot, and also activates platelets. The body makes thrombin after a blood vessel is injured; too much or too little clotting can cause disease.
\medterm{Thromboangiitis Obliterans} A tobacco-related inflammatory disease blocking small and medium arteries and veins, often causing painful fingers, toes, ulcers, or gangrene.
\medterm{Thromboangiitis Obliterans (Buerger Disease)} An inflammatory disease of small and medium blood vessels, strongly associated with tobacco
use. It can cause pain, color changes, ulcers, or gangrene in the hands or feet. Complete
tobacco cessation is essential; further treatment depends on ischemia, tissue loss, and
other complications.

\synonyms
Buerger Disease

\medterm{Thrombocyte (Platelet)} A small blood-cell fragment made in the bone marrow. Platelets stick to damaged blood vessels, become activated, and join together to form the first plug that stops bleeding. Too few platelets or poorly working platelets can cause easy bruising, tiny skin spots, nosebleeds, or other bleeding. Too many or overly active platelets can increase the risk of unwanted clots.
\medterm{Thrombocytes} Another name for platelets: small blood-cell fragments made in the bone marrow that form plugs to stop bleeding. Too few can cause bruising or bleeding, while too many or overly active platelets can increase the risk of unwanted clots.

\medterm{Thrombocytopathy} A disorder of platelet function that impairs primary hemostasis despite a normal or
near-normal platelet count. It may be inherited or acquired and can cause mucocutaneous
bleeding, easy bruising, prolonged bleeding after procedures, or menorrhagia.

\medterm{Thrombocytopenia} Thrombocytopenia is a low platelet count caused by reduced production, increased destruction or consumption, sequestration, or dilution and may increase bleeding risk.

\synonyms
Low Platelet Count

\medterm{Thrombocytopenia (Low Platelet Count)} Alternate index label for Low Platelet Count.

\medterm{Thrombocytopenic} Characterized by too few circulating platelets. It may cause easy bruising, nosebleeds, heavy menstrual bleeding, or other bleeding, although some people have no symptoms.
\medterm{Thrombocytosis} An elevated platelet count (>450,000/$\mu$L). Reactive (secondary): infection, inflammation,
iron deficiency, malignancy, post-splenectomy, hemorrhage, exercise—usually mild, no
thrombotic risk.

\medterm{Thromboelastography} A viscoelastic whole-blood assay that measures the initiation, formation, strength, and
breakdown of a clot. It provides a functional overview of coagulation and fibrinolysis
and may guide targeted blood-component therapy in major bleeding.

\medterm{Thromboembolism} Obstruction of a blood vessel by a clot that formed elsewhere and traveled through the circulation.

\medterm{Thrombolysis} The dissolution of a blood clot using a medicine such as a fibrinolytic agent. It is used only in selected conditions because it can cause serious bleeding.

\medterm{Thrombolysis For PE} Emergency treatment using a clot-dissolving medicine for a life-threatening lung clot causing low blood pressure or shock. It can restore blood flow quickly but carries a serious risk of bleeding.

\medterm{Thrombolysis For Stroke} Emergency treatment for selected strokes caused by a blood clot. A clot-dissolving medicine is given through a vein within a limited time after symptoms begin, after brain imaging excludes bleeding; serious bleeding is a possible complication.

\medterm{Thrombolytic Agents} Medicines that activate fibrinolysis and dissolve an established blood clot; they are used selectively for time-critical conditions because they can cause serious bleeding.
\medterm{Thrombolytic Therapy} Treatment with a medicine that breaks down a blood clot and may restore blood flow in selected emergencies such as certain strokes, heart attacks, or lung clots. The main danger is serious bleeding, including bleeding in the brain.

\medterm{Thrombophilia} A tendency for blood to form clots too easily. It may be inherited or develop with cancer, pregnancy, some medicines, inflammation, or other conditions and can increase the risk of clots in veins.

\medterm{Thrombophilia Testing} Thrombophilia testing evaluates for inherited and acquired hypercoagulable states in
patients with venous thromboembolism, recurrent pregnancy loss, or strong family history
of VTE.

\medterm{Thrombophlebitis} Thrombophlebitis is inflammation of a vein associated with a blood clot and may be superficial or involve deep veins with risk of pulmonary embolism.

\medterm{Thrombopoiesis} The production of platelets from megakaryocytes in the bone marrow. It is regulated
principally by thrombopoietin and involves megakaryocyte maturation, proplatelet
formation, and release of circulating platelets.

\medterm{Thrombosed Intracranial Aneurysm} An intracranial arterial aneurysm that has undergone partial or complete thrombosis of
its lumen, which may reduce rupture risk but can cause mass effect on adjacent brain
structures. Presents with headache, cranial nerve palsies, or seizures depending on
location. Diagnosis: CT, MRI, or conventional angiography.

\medterm{Thrombosis} The formation of a blood clot within the cardiovascular system during life, potentially
obstructing blood flow and causing tissue ischemia or infarction. Arterial thrombi
typically form on ruptured atherosclerotic plaques, are composed primarily of platelets
and fibrin, and appear pale and grossly attached to the vessel wall.

\medterm{Thrombosis Of The Inferior Vena Cava} Formation of a blood clot within the inferior vena cava, which may occur in the setting
of IVC filters, malignancy, hypercoagulable states, or extension from iliofemoral deep
vein thrombosis. May cause bilateral lower extremity edema, abdominal pain, and hepatic
or renal dysfunction. Requires anticoagulation and possible thrombectomy.

\medterm{Thrombotic Microangiopathy} A group of disorders in which tiny clots form in small blood vessels, damaging red blood cells and lowering the platelet count. The clots can injure the kidneys, brain, heart, or other organs. Causes include immune disorders, infections, medicines, pregnancy, and severe illness; urgent specialist treatment is often needed.

\medterm{Thrombotic Thrombocytopenic Purpura} A rare, life-threatening blood disorder in which small clots form in tiny blood vessels and use up platelets. It can damage the brain, kidneys, heart, and other organs.
\medterm{Thrombotic Vasculopathy} A spectrum of disorders characterized by thrombus formation within blood vessels of
various sizes, often associated with antiphospholipid syndrome, disseminated
intravascular coagulation, thrombotic thrombocytopenic purpura, or malignancy.

\medterm{Thromboxane} A fat-derived chemical messenger made mainly by activated platelets. It encourages platelets to stick together and narrows blood vessels, helping clots form but also contributing to unwanted clotting.
\medterm{Thrombus} A blood clot that forms within a blood vessel or heart and remains attached at its site of formation, potentially obstructing blood flow or embolizing.
\medterm{Thrombus Organization} The process by which a thrombus is gradually replaced by granulation tissue and
eventually fibrous scar tissue. The process begins when inflammatory cells, particularly
macrophages, infiltrate the thrombus and begin digesting the fibrin and dead cells
through fibrinolysis.

\medterm{Thrush} A Candida infection, usually of the mouth or throat, causing creamy white plaques, soreness, altered taste, or discomfort with eating or swallowing. Vaginal candidiasis is a separate condition. Risk increases with antibiotics, inhaled corticosteroids, diabetes, or immune suppression.

\synonyms
Oral Thrush

\medterm{Thrush - Children And Adults} Oral candidiasis caused by Candida overgrowth, producing removable white plaques, soreness, redness, altered taste, or cracks at the mouth corners; risk rises with antibiotics, inhaled steroids, diabetes, or immune suppression.

\medterm{Thrush And Other Yeast Infections In Children} Thrush is Candida overgrowth in the mouth, producing white plaques, soreness, or feeding difficulty; Candida can also affect skin folds or diaper areas. Persistent, recurrent, severe, or immune-associated infection needs pediatric evaluation and appropriate antifungal treatment.

\medterm{Thrush In Newborns} Oral Candida infection in an infant, causing white plaques on the tongue or mouth lining that may interfere with feeding and can be passed between infant and breastfeeding parent.

\medterm{Thrush, Oral} Alternate index label; see \textit{Oral thrush}.

\medterm{Thumb Arthritis} Degeneration or inflammation of a thumb joint, especially the joint at the base of the thumb, causing pain, stiffness, and reduced grip.

\synonyms
Basal Joint Arthritis

\medterm{Thumb CMC Arthritis} Thumb carpometacarpal arthritis: Eaton-Littler classification: I (normal), II
(slight narrowing, osteophytes <2mm), III (narrowing, osteophytes >2mm, subluxation), IV
(pantrapezial arthritis). Clinical: base of thumb pain, pinch weakness, thenar wasting,
squaring (subluxation).

\medterm{Thumb-Sucking} A common childhood habit of sucking the thumb; persistent habit can affect dental and facial development.
\medterm{Thymic Hyperplasia} Thymic hyperplasia is enlargement of the thymus from increased lymphoid tissue or rebound growth and may be incidental. It can be associated with autoimmune disease, infection, endocrine conditions, or stress; imaging and clinical context distinguish it from a mass.

\medterm{Thymocyte} An immature T lymphocyte developing in the thymus, where it undergoes selection before becoming a mature T cell.
\medterm{Thymoma} A tumor of thymic epithelial cells, associated with myasthenia gravis, pure red cell
aplasia, and hypogammaglobulinemia. Malignant behavior is determined by capsular
invasion and stage.

\synonyms
Neoplasm Thymomic (Thymoma)

\medterm{Thymus} Primary lymphoid organ in the anterior mediastinum, largest in childhood, involutes with
age. Site of T-cell maturation (thymopoiesis): thymocytes undergo positive and negative
selection to produce self-tolerant, MHC-restricted T cells. Epithelial cells produce
thymic hormones (thymosin, thymopoietin).

\medterm{Thymus Gland} An immune-system organ behind the breastbone where immature T lymphocytes learn to recognize harmful foreign material without attacking the body’s own tissues. It is most active during childhood and gradually becomes smaller after puberty.
\medterm{Thyrocalcitonin} Another name for calcitonin, a hormone made by thyroid C cells that lowers blood calcium mainly by reducing bone resorption.

\medterm{Thyroglobulin Surveillance} Thyroglobulin produced only by thyroid follicular cells serves as a tumor marker after
total thyroidectomy for differentiated thyroid carcinoma. In patients who have undergone
complete thyroidectomy and radioactive iodine ablation thyroglobulin should be
undetectable.

\medterm{Thyroglossal Duct Cyst} Congenital midline neck mass resulting from persistence of the thyroglossal duct the
embryological tract of thyroid descent. Most common congenital neck mass presenting in
childhood. Located at or near the hyoid bone and moves with swallowing and tongue
protrusion. Diagnosis is by ultrasound demonstrating a cystic midline neck mass.

\medterm{Thyroid} An endocrine gland in the front of the neck that produces thyroid hormones, which regulate metabolism, growth, and development.

\medterm{Thyroid Cancer} A malignant tumor of the thyroid. Major types include papillary, follicular, medullary, and anaplastic carcinoma, which differ in behavior and treatment.

\synonyms
Cancer, Thyroid

\medterm{Thyroid Cancer - Medullary Carcinoma} Medullary thyroid carcinoma is cancer of thyroid C cells that produce calcitonin and may be inherited as part of a multiple endocrine neoplasia syndrome.

\medterm{Thyroid Cancer - Papillary Carcinoma} Papillary thyroid carcinoma is the most common thyroid cancer and usually grows slowly, although it can spread to nearby lymph nodes.

\medterm{Thyroid Cancer Staging} American Joint Committee on Cancer TNM staging for differentiated thyroid carcinoma
incorporates tumor size, extent of invasion, lymph node metastases, and distant
metastases. Age is a critical prognostic factor with patients under fifty having more
favorable outcomes. Stage I is any tumor without metastases in patients under fifty.

\medterm{Thyroid Cartilage} The largest cartilage of the larynx, formed by two plates that meet in front and create the laryngeal prominence. It
protects the vocal folds and provides attachment for muscles and ligaments.


\medterm{Thyroid Disorders} Conditions affecting the thyroid gland, a butterfly-shaped endocrine organ in the neck
that produces hormones regulating metabolism, growth, and development.

\medterm{Thyroid Eye Disease} Autoimmune inflammatory disorder of the orbital tissues associated with Graves disease
characterized by expansion of extraocular muscles and orbital fat. Pathogenesis involves
orbital fibroblasts expressing the TSH receptor targeted by the same thyroid-stimulating
immunoglobulins affecting the thyroid.

\medterm{Thyroid Fine-Needle Aspiration} Image-guided sampling of thyroid-nodule cells for cytologic classification. Results are
commonly reported using the Bethesda system and guide surveillance, molecular testing,
surgery, or other management.

\medterm{Thyroid Follicular Carcinoma} A cancer arising from thyroid follicle cells and the second most common differentiated thyroid cancer. It tends to spread through blood vessels to the lungs or bones; diagnosis requires showing invasion of the capsule or blood vessels.

\medterm{Thyroid Function Tests} Laboratory assessment of the hypothalamic-pituitary-thyroid axis. TSH is the most
sensitive screening test with elevated levels indicating primary hypothyroidism and
suppressed levels indicating hyperthyroidism. Free T4 measures the biologically active
unbound thyroxine.

\medterm{Thyroid Gland} An endocrine gland located in the anterior neck, consisting of two lobes (right and
left) connected by an isthmus, lying anterior to the trachea between the cricoid
cartilage and the suprasternal notch.

\medterm{Thyroid Gland Enlargement} Alternate index label; see \textit{Goiter}.

\medterm{Thyroid Gland Removal} Surgery to remove part or all of the thyroid gland for cancer, overactivity, a large goiter, or a suspicious nodule. Possible effects include bleeding, voice changes from nerve injury, low calcium, and the need for lifelong thyroid hormone.


\medterm{Thyroid Gland, Diseases Of} Conditions affecting the thyroid gland, including underactivity, overactivity, inflammation, enlargement, nodules, and cancer. They may change energy, weight, temperature tolerance, heart rate, bowel habits, mood, or neck appearance and are assessed with examination and blood tests.
\medterm{Thyroid Hormone} The hormones T3 and T4, made by the thyroid gland, that help control energy use, temperature, heart rate, growth, and brain development. Too little can slow body functions, while too much can make them overactive.

\medterm{Thyroid Hormones} Thyroxine (T4) and triiodothyronine (T3), hormones that regulate metabolic rate, growth, brain development, heat production, and cardiovascular and gastrointestinal activity.

\medterm{Thyroid Nodule} A discrete lump within the thyroid gland that may be solid, cystic, benign, or malignant. Most are benign, but evaluation
uses clinical assessment, thyroid testing, ultrasound, and biopsy when indicated.

\medterm{Thyroid Nodule Evaluation} Systematic approach to evaluation of thyroid nodules detected by physical examination or
imaging. Thyroid ultrasound characterizes nodule features including size, composition,
echogenicity, margins, calcifications, and shape. TI-RADS classification stratifies
malignancy risk and guides biopsy decisions.

\medterm{Thyroid Nodules} A thyroid nodule is a discrete growth within the thyroid and is often benign. Ultrasound, thyroid-function testing, and selected needle biopsy assess cancer risk; rapid growth, hoarseness, swallowing difficulty, or breathing problems needs prompt review.

\synonyms
Nodule Thyroid (Thyroid Nodules)

\medterm{Thyroid Papillary Carcinoma} The most common differentiated thyroid cancer, arising from thyroid follicular cells. It often
grows slowly and has a favorable prognosis, although outcome depends on tumor features,
spread, age, and response to treatment.

\medterm{Thyroid Peroxidase Antibody} An autoantibody against thyroid peroxidase, an enzyme required for thyroid-hormone
synthesis. It supports autoimmune thyroid disease, especially Hashimoto thyroiditis, but
may be present in some euthyroid individuals.

\medterm{Thyroid Peroxidase Test} A blood test measuring antibodies against thyroid peroxidase, an enzyme needed to make thyroid hormones. Elevated levels support autoimmune thyroid disease but do not by themselves show how well the thyroid is working.

\medterm{Thyroid Preparation Overdose} Taking too much thyroid hormone medicine can cause a fast or irregular heartbeat, shaking, sweating, anxiety, vomiting, fever, or, in severe cases, heart failure or a dangerous thyroid crisis.

\medterm{Thyroid Scan} Nuclear medicine imaging using I-123 or Tc-99m. Evaluates thyroid function and nodules.
Hot nodules rarely malignant. Cold nodules require further evaluation.

\medterm{Thyroid Scintigraphy} Thyroid imaging using I-123, I-131, or Tc-99m pertechnetate. Assesses function, nodules,
ectopic tissue. Hot nodules (functioning) rarely malignant; cold nodules require
ultrasound/FNA.

\medterm{Thyroid Stimulating Hormone} TSH (thyrotropin): the anterior pituitary hormone that stimulates the thyroid to produce
and release T4 and T3, controlled by TRH (hypothalamus) and negative feedback from
thyroid hormones.

\medterm{Thyroid Storm} Life-threatening exacerbation of thyrotoxicosis presenting with hyperpyrexia,
tachycardia, heart failure, altered mental status, and multi-organ dysfunction.
Precipitated by infection, surgery, trauma, or iodine load in patients with untreated
hyperthyroidism. Mortality rates range from twenty to thirty percent.

\medterm{Thyroid Surgery} Operations that remove part or all of the thyroid gland. They may be performed for selected cancers, enlarging or
compressive nodules, goiter, hyperthyroidism, or diagnostic uncertainty; the procedure
depends on the disease and its extent.

\medterm{Thyroid Ultrasound} High-frequency sonographic imaging used to evaluate thyroid size, nodules, cysts,
cervical lymph nodes, and structural features associated with malignancy. Suspicious
nodules may require ultrasound-guided fine-needle aspiration.

\medterm{Thyroid, Overactive} Alternate index label; see \textit{Hyperthyroidism (overactive thyroid)}.

\medterm{Thyroid-Stimulating Hormone} A hormone made by the pituitary gland that tells the thyroid to make thyroid hormones. Blood TSH testing is commonly used to check thyroid function, although results must be interpreted with other hormone levels and the person's symptoms.

\synonyms
TSH

\medterm{Thyroidectomy} Surgery to remove part or all of the thyroid gland. It may treat thyroid cancer, an overactive thyroid, or a large goiter; possible complications include bleeding, voice-nerve injury, and low calcium.
\medterm{Thyroiditis} Inflammation of the thyroid gland with various etiologies. Hashimoto thyroiditis:
autoimmune, most common cause of hypothyroidism in iodine-sufficient areas. Subacute (De
Quervain) thyroiditis: post-viral, presents with painful thyroid and transient
hyperthyroidism followed by hypothyroidism.

\medterm{Thyroiditis Classification} Inflammatory conditions of the thyroid classified by clinical presentation and etiology.
Hashimoto thyroiditis is chronic autoimmune with hypothyroidism. Graves disease is
autoimmune with hyperthyroidism. Subacute thyroiditis is painful and likely viral.

\medterm{Thyroiditis, Chronic Lymphocytic} Alternate index label; see \textit{Hashimoto's disease}.

\medterm{Thyrotoxic Myopathy} Muscle weakness caused by an overactive thyroid, usually affecting the hips, thighs, shoulders, and upper arms. Weakness may occur without pain and can improve when thyroid hormone levels return to normal; severe weakness or swallowing or breathing difficulty needs urgent care.

\synonyms
Hyperthyroid Myopathy

\medterm{Thyrotoxic Periodic Paralysis} Recurrent episodes of transient flaccid weakness associated with thyrotoxicosis, usually
accompanied by hypokalemia caused by intracellular potassium shift. Attacks often affect
the proximal limbs and are triggered by rest after exercise or carbohydrate-rich meals.

\medterm{Thyrotoxicosis} A condition caused by too much thyroid hormone activity in the body, whatever the source. It may cause weight loss, fast heartbeat, tremor, sweating, anxiety, diarrhea, heat intolerance, or severe illness.
\medterm{Thyrotroph Adenoma} A thyrotroph adenoma is a pituitary tumor arising from thyroid-stimulating-hormone-producing cells and may cause excess TSH and abnormal thyroid function.

\medterm{Thyrotrophin-Releasing Hormone} A hormone made in the hypothalamus, a part of the brain, that tells the pituitary gland to release thyroid-stimulating hormone and prolactin. It helps link brain signals with thyroid activity and some reproductive functions.
\medterm{Thyrotrophin-Stimulating Hormone} Thyroid-stimulating hormone, a pituitary hormone that stimulates thyroid hormone production.
\medterm{Thyrotropin} Thyroid-stimulating hormone, a pituitary hormone that stimulates thyroid growth and thyroid-hormone production; it is also called TSH.

\medterm{Thyrotropin Alfa} A recombinant form of thyroid-stimulating hormone used to stimulate residual thyroid tissue or differentiated thyroid cancer cells during selected diagnostic or treatment protocols.

\medterm{Thyroxine} The thyroid hormone T4, which contains four iodine atoms. The thyroid releases it into the blood, and much of it is converted to the more active hormone triiodothyronine (T3).

\synonyms
T4

\medterm{Thyroxine-Binding Globulin} The principal transport protein for thyroid hormones in the blood, binding most circulating thyroxine and triiodothyronine.

\medterm{TI-RADS} The Thyroid Imaging Reporting and Data System is an ultrasound-based system that assigns a risk category to thyroid nodules
using features such as composition, echogenicity, shape, margins, and calcification. It
helps guide follow-up and biopsy decisions.

\medterm{TIA} A transient ischemic attack, a temporary episode of neurological dysfunction from reduced brain blood flow.
\synonyms
Transient ischemic attack; mini-stroke

\medterm{Tiazofurin} Tiazofurin is an investigational nucleoside analogue that interferes with NAD synthesis and has been studied as an antineoplastic agent.

\medterm{Tibia} The tibia is the larger medial bone of the lower leg, extending from the knee to the ankle; it bears much of the body's weight and forms parts of both joints.
Proximal: tibial plateau (medial/lateral condyles) and tibial tuberosity (patellar
ligament insertion). Shaft: sharp subcutaneous anterior border. Distal: medial malleolus
articulating with the talus at the ankle.

\medterm{Tibia Vara} Tibia vara is inward angulation of the upper shinbone causing bowing of the leg, classically associated with Blount disease or other growth disorders.

\medterm{Tibial Artery} An artery of the lower leg, chiefly the anterior and posterior tibial arteries, supplying the leg, ankle, and foot and contributing to pedal pulses.

\medterm{Tibial Fracture} A break in the tibia, the major weight-bearing bone of the lower leg, often caused by high-energy trauma, twisting injury, or repetitive stress.

\medterm{Tibial Nerve} Sciatic division (L4-S3) continuing from posterior thigh through the popliteal fossa
into the deep posterior leg, innervating gastrocnemius, soleus, tibialis posterior,
flexor digitorum longus, and flexor hallucis longus. Passes behind the medial malleolus
as the tarsal tunnel, dividing into medial/lateral plantar nerves for the foot.

\medterm{Tibial Shaft Fractures} Breaks in the middle portion of the shinbone, often caused by a high-energy injury but sometimes by a twist or lower-force event. Pain, swelling, deformity, or inability to bear weight are common. Treatment depends on alignment, skin and muscle damage, and whether the fracture is open.

\medterm{Tibialcrural Artery} A lower-leg arterial vessel or branching region supplying the ankle and foot, usually discussed in relation to tibial and fibular arterial anatomy.

\medterm{Tibialis Anterior Muscle} Superficial anterior leg muscle from lateral tibial condyle and shaft to medial
cuneiform and first metatarsal base. Innervated by deep peroneal nerve (L4-L5). Primary
dorsiflexor for foot clearance during gait swing phase; inverts foot. Weakness causes
foot drop with steppage gait.

\medterm{Tibialis Cranialis} The animal-anatomy term corresponding broadly to the tibialis anterior muscle, which dorsiflexes the ankle and helps invert the foot.

\medterm{Tibialis Posterior Muscle} Deepest posterior leg muscle from posterior tibia, fibula, and interosseous membrane to
navicular tuberosity, medial cuneiform, and sustentaculum tali. Innervated by tibial
nerve (L4-L5). Plantarflexes and inverts foot; primary dynamic medial longitudinal arch
stabilizer.

\medterm{Tic} A sudden, rapid, recurrent, nonrhythmic movement or sound that a person may feel an urge to make. Tics can be motor or vocal and may worsen with stress.
\medterm{Tic Disorder} A neurologic disorder characterized by recurrent, sudden, rapid, nonrhythmic motor movements or vocalizations that may be temporarily suppressible.

\medterm{Tic Douloureux} A chronic pain condition characterized by recurrent, sudden, severe, electric shock-like
pain in the distribution of the trigeminal nerve (typically V2/V3). See Trigeminal
Neuralgia.

\medterm{Tic Douloureux (Trigeminal Neuralgia)} Alternate index label for Trigeminal Neuralgia.

\medterm{Ticarcillin} A penicillin antibiotic with activity against many gram-negative bacteria, including Pseudomonas; it is usually given with clavulanate and can cause allergy or electrolyte disturbances.

\medterm{Tick} A small blood-feeding arachnid that can attach to people or animals and transmit infections such as Lyme disease, rickettsial diseases, and babesiosis.

\medterm{Tick Bite} A bite from a tick, which may cause local irritation and can sometimes spread infections such as Lyme disease, ehrlichiosis, or Rocky Mountain spotted fever. Remove the tick promptly and watch for fever or a spreading rash.

\medterm{Tick Paralysis} Acute, progressive flaccid paralysis caused by neurotoxin released by an attached tick.
Weakness usually begins in the legs and ascends; respiratory failure can occur and
symptoms improve after prompt tick removal.

\medterm{Tick Removal} Promptly removing an attached tick with fine-tipped tweezers by pulling steadily close to the skin. Clean the site afterward and monitor for expanding redness, fever, or other illness; do not burn, crush, or cover the tick with chemicals.

\medterm{Tick Typhus} A rickettsial infection transmitted by ticks, usually causing fever, headache, rash, and sometimes severe systemic illness.

\medterm{Tick-Borne Encephalitis} A viral infection transmitted mainly by ticks that can cause fever, meningitis, encephalitis, or lasting neurological complications.

\medterm{Tickborne Relapsing Fever} A vector-borne illness caused by Borrelia hermsii and related species, transmitted by
Ornithodoros ticks. It presents with recurrent episodes of fever, headache, myalgia, and
arthralgia, separated by afebrile intervals. The relapsing pattern results from
antigenic variation of the spirochete's surface lipoproteins.

\medterm{Ticks} Small arachnids that attach to skin and may transmit infections or cause tick paralysis.

\synonyms
Ticks Bite (Ticks)

\medterm{Ticks Bite (Ticks)} Alternate index label for Ticks.

\medterm{Ticlopidine} An antiplatelet medicine that reduces platelet aggregation and clot formation.

\medterm{Tics} Sudden, rapid, recurrent, nonrhythmic stereotyped movements (motor tics) or
vocalizations (vocal tics) that are usually brief and suppressible. Transient tics (most
common in childhood) vs. chronic tic disorder; Tourette syndrome: multiple motor + at
least one vocal tic for >1 year beginning before age 18.

\medterm{Tidal Volume} The volume of air inspired or expired during each normal breath, approximately 500 mL
(6-8 mL/kg) in a resting adult. Tidal volume is measured by spirometry. Minute
ventilation = tidal volume x respiratory rate (~6 L/min). Only a portion reaches the
alveoli (alveolar ventilation) after subtracting dead space (~150 mL).

\medterm{Tidy Wound} A clean, sharply incised wound with minimal contamination, crushing, or tissue devitalization and clearly defined viable edges. After appropriate assessment and cleaning, it may often be closed primarily.
\medterm{Tietze Syndrome (Costochondritis And Tietze Syndrome)} Alternate index label for Costochondritis and Tietze Syndrome.

\medterm{Tigecycline} Tigecycline is a glycylcycline antibiotic that binds the bacterial 30S ribosomal subunit and inhibits protein synthesis; it treats selected complicated infections but commonly causes nausea and has a warning for increased mortality in some severe infections.

\medterm{Tight Junction} A cell-to-cell junction that seals the space between epithelial cells and controls paracellular movement while helping maintain distinct apical and basal cell surfaces.

\medterm{Tilorone} Tilorone is a synthetic antiviral and immune-modulating compound used or studied in some countries; evidence and indications vary by setting.

\medterm{Tilt Table Test} Evaluation of unexplained syncope by reproducing orthostatic stress. Patient is strapped
to a motorized table tilted from supine to 60-80 degrees for 20-45 minutes. Positive:
reproduction of symptoms with hypotension and/or bradycardia (vasovagal syncope, delayed
orthostatic hypotension).

\medterm{Tilt Table Testing} A supervised test in which a patient is moved from a lying to an upright or head-up tilt position while blood pressure, heart rate, and symptoms are monitored. It may help evaluate reflex syncope, orthostatic hypotension, or selected autonomic disorders when the diagnosis remains uncertain; it does not replace assessment for cardiac or neurologic causes.

\medterm{Tilted Disc Syndrome} A congenital appearance in which the optic disc is angled or rotated, often with an unusual shape of nearby retina and blood vessels. Vision may be normal, but some people have visual-field changes or other eye conditions.

\medterm{Time In Range} Percentage of time that continuous glucose monitoring readings remain within the target
glucose range of seventy to one hundred eighty milligrams per deciliter for most adults
with diabetes. General target is greater than seventy percent corresponding to
approximately seventeen hours per day.

\medterm{Time Of Death} The moment or estimated time at which biological death occurred, a key medicolegal
question. Estimated from postmortem changes: body cooling (algor mortis), rigor mortis,
lividity (livor mortis), decomposition/putrefaction, insect colonization, and chemical
changes (vitreous potassium, vitreous glucose, body temperature).

\medterm{Time-Of-Flight PET} A PET acquisition and reconstruction technique that uses the measured arrival-time
difference of coincident annihilation photons to improve localization of the
annihilation event. It can improve signal-to-noise performance, especially in larger
patients.

\medterm{Timed Up-And-Go Test} A functional mobility test measuring the time required to stand from a chair, walk a
short distance, turn, return, and sit. Prolonged performance suggests impaired mobility
and increased fall risk but must be interpreted with gait, vision, cognition, and
environment.

\medterm{TIMI Risk Score} A clinical score estimating short-term risk in people with unstable angina or a non-ST-elevation heart attack. It uses age, risk factors, known artery disease, recent symptoms, ECG changes, and heart-injury blood tests to help guide treatment; it supports but does not replace clinical judgment.

\synonyms
TIMI Score; TIMI Risk Score for UA/NSTEMI; TIMI Score for STEMI

\medterm{Timing Of Breastfeeding} The timing and frequency of breastfeeds, usually guided by early hunger cues and effective milk transfer rather than a rigid clock schedule, with special considerations for newborns and preterm infants.

\medterm{Timolol} Timolol is a non-selective beta blocker used in glaucoma eye drops and for selected cardiovascular conditions; topical ocular use can still cause bradycardia, hypotension, bronchospasm, or fatigue.

\medterm{Timolol Maleate} The maleate salt of timolol, a beta blocker used in glaucoma and selected cardiovascular conditions.
\medterm{Timothy Syndrome} A rare multisystem disorder caused by pathogenic variants in CACNA1C, characterized by
marked QT prolongation and congenital heart disease with syndactyly. Developmental,
neurological, immune, and facial features may also occur.

\medterm{Tincture} Alcohol extract of plant or chemical substance. Common concentration: 10-40\% alcohol.
Examples: tincture of iodine (antiseptic), tincture of benzoin (skin protector). Also
used for liquid medications.

\medterm{Tinea} Dermatophytosis: superficial fungal infection of keratinized tissue (skin, hair, nails) by dermatophytes (Trichophyton, Microsporum, Epidermophyton).
\medterm{Tinea (Dermatophytosis)} Superficial fungal infection by dermatophytes. Tinea corporis (ringworm): annular,
scaly, erythematous plaque with central clearing. Tinea pedis (athlete's foot):
interdigital scaling, maceration, and fissuring. Tinea cruris (jock itch): erythematous,
pruritic rash in groin with raised border.

\synonyms
Dermatophytosis

\medterm{Tinea Barbae} Dermatophyte infection of the beard and mustache area, presenting with erythematous,
scaly, follicular papulopustules or deep inflammatory nodules, often acquired from
contact with infected animals or through grooming. It may be inflammatory (kerion-like)
or superficial, and differential diagnosis includes bacterial folliculitis and acne.

\medterm{Tinea Capitis} A dermatophyte infection of the scalp and hair shafts, most common in children. It may
cause scaling, alopecia, broken hairs, black dots, or an inflammatory kerion.

\medterm{Tinea Circinata} An older term for tinea corporis (ringworm) of the body, presenting as annular,
erythematous, scaly plaques with active, raised borders and central clearing. It is
caused by dermatophyte fungi and may be acquired through direct contact with infected
humans, animals, or fomites.

\medterm{Tinea Corporis} A dermatophyte infection of glabrous skin, presenting as annular, erythematous, scaly
plaques with raised, active borders and central clearing. It is contagious through
direct contact with infected persons, animals, or contaminated objects. Diagnosis is by
KOH preparation showing branching septate hyphae.

\medterm{Tinea Cruris} Alternate index label; see \textit{Jock itch}.

\medterm{Tinea Faciei} Dermatophyte infection of the facial skin, often presenting with erythematous, scaly,
annular or ill-defined plaques that may mimic eczema or rosacea, leading to frequent
misdiagnosis. It is commonly caused by Trichophyton or Microsporum species and may be
acquired from pets, other humans, or self-inoculation from other body sites.

\medterm{Tinea Favosa} A chronic dermatophyte infection of the scalp or skin, usually caused by Trichophyton schoenleinii, producing yellow cup-shaped crusts and possible scarring alopecia.

\medterm{Tinea Imbricata} A chronic dermatophyte infection caused by Trichophyton concentricum, endemic in the
South Pacific and Southeast Asia, characterised by widespread, concentric, overlapping,
polycyclic scaly rings producing a "wood-grain" pattern on the trunk and limbs.

\medterm{Tinea Incognito} A dermatophyte infection whose clinical appearance has been altered by inappropriate
topical corticosteroid or immunosuppressive treatment. It may become less scaly and more
extensive or atypical; fungal testing is often needed before directed therapy.

\medterm{Tinea Nigra} A superficial fungal infection of the stratum corneum caused by Hortaea werneckii,
presenting as a painless, brown to black, non-scaly macule or patch, most commonly on
the palms or soles. It is acquired in tropical and subtropical regions and is often
confused with acral melanoma.

\medterm{Tinea Pedis} A dermatophyte infection of the feet presenting in interdigital, moccasin,
vesiculobullous, or ulcerative patterns. It may coexist with tinea unguium. Diagnosis is
clinical or mycologic when uncertain; treatment uses topical antifungals for limited
disease and oral therapy for extensive, refractory, or nail-associated infection.

\medterm{Tinea Pedis (Athlete'S Foot)} Alternate index label for Athlete's Foot.

\medterm{Tinea Unguium} A fungal infection of the nail plate, also called onychomycosis, causing thickening, discoloration, brittleness, or separation of the nail.

\medterm{Tinea Versicolor} A superficial fungal infection caused by Malassezia species, presenting as hypo- or
hyperpigmented macules with fine scale on the trunk and proximal extremities. It is more
noticeable after sun exposure due to unaffected tanning of surrounding skin.

\synonyms
Pityriasis Versicolor

\medterm{Tinel Sign} Tingling or an electric sensation felt along a nerve when the nerve is gently tapped. It may suggest nerve irritation or recovery after injury, but it is not diagnostic by itself and must be interpreted with the examination and other tests.

\medterm{Tinidazole} Tinidazole is a nitroimidazole antimicrobial effective against anaerobic bacteria and protozoa such as Giardia and Trichomonas; metallic taste, gastrointestinal effects, neuropathy, and an alcohol interaction may occur.

\medterm{Tinnitus} Tinnitus is perception of sound without an external source, often related to hearing loss, noise exposure, ear disease, medications, vascular causes, or neurologic factors.

\synonyms
Ringing In The Ear

\medterm{Tinnitus (Ringing In The Ears)} Alternate index label for Ringing in the Ears.

\medterm{Tinzaparin} A low-molecular-weight heparin anticoagulant used for prevention or treatment of venous thromboembolism; bleeding and heparin-induced thrombocytopenia are important risks.

\medterm{Tinzaparin Sodium} The sodium salt of tinzaparin, a low-molecular-weight heparin used to prevent or treat venous blood clots; bleeding and heparin-induced thrombocytopenia are important risks.

\medterm{Tiopronin} A thiol-containing medicine used in selected cases of cystinuria to increase urinary cystine solubility; rash, proteinuria, and blood abnormalities may occur.

\medterm{Tiotropium Bromide} A long-acting muscarinic antagonist bronchodilator used for maintenance therapy
in COPD. Once-daily dosing provides sustained bronchodilation. It reduces exacerbations
and improves lung function. It is often combined with a LABA for additive
bronchodilation.

\medterm{Tipifarnib} An experimental medicine studied for some cancers and blood disorders. It blocks a cell enzyme needed by certain cancer cells to grow, but it is not a routine treatment and its usefulness depends on the disease and research results.

\medterm{TIPS} TIPS, or transjugular intrahepatic portosystemic shunt, connects portal and hepatic veins to reduce portal hypertension, often for variceal bleeding or refractory ascites.

\medterm{Tirapazamine} Tirapazamine is an investigational prodrug activated in low-oxygen tissue and studied to target hypoxic tumor cells.

\medterm{Tirbanibulin} A topical medicine used in short courses to treat actinic keratoses on the face or scalp; local redness, scaling, or irritation may occur.

\medterm{Tired} Feeling a lack of energy or physical or mental stamina, which may result from sleep loss, exertion, infection, anemia, endocrine disease, medication, depression, or other illness.

\medterm{Tiredness} Fatigue or reduced energy and endurance that can be physical, mental, or both; persistent unexplained fatigue warrants assessment for sleep, mood, systemic, medication, and metabolic causes.

\medterm{Tiredness (Fatigue)} A subjective lack of energy or exhaustion that may result from inadequate sleep, illness, psychological stress, anaemia, or many other causes.

\synonyms
Fatigue

\medterm{Tirofiban} An intravenous glycoprotein IIb/IIIa platelet inhibitor used in selected acute coronary syndromes or coronary interventions to reduce platelet aggregation.

\medterm{Tissierella} A genus of anaerobic bacteria found in the environment and sometimes isolated from polymicrobial human infections.

\medterm{Tissierella Praeacuta} An anaerobic bacterium rarely associated with human infection; significance depends on recovery from a compatible sterile-site infection rather than colonization.

\medterm{Tissue} A group of similar cells and their surrounding materials that work together to perform a
specific function.

\medterm{Tissue Adhesions} Bands of fibrous scar tissue that abnormally connect surfaces or organs, often after surgery, inflammation, infection, or bleeding and potentially causing pain, infertility, or obstruction.

\medterm{Tissue Banks} Facilities that recover, screen, process, preserve, store, and distribute donated human tissue for transplantation, research, or reconstructive treatment under safety and traceability standards.

\medterm{Tissue Congestion} Accumulation of blood or fluid within tissue because of impaired venous drainage, inflammation, obstruction, or increased vascular flow, producing swelling, redness, or impaired function.

\medterm{Tissue Culture} Growth of cells or tissue outside the body in a controlled nutrient medium for research, diagnosis, vaccine production, regenerative medicine, or drug testing.

\medterm{Tissue Distribution} Reversible movement of a chemical from blood into tissues and organs. Distribution is
influenced by perfusion, lipid solubility, protein binding, tissue affinity, barriers,
body composition, and time after exposure.

\medterm{Tissue Elastance} The change in pressure required to produce a given change in tissue or organ volume, the inverse of compliance and important in lung and chest-wall mechanics.

\medterm{Tissue Engineering} The use of cells, supportive materials called scaffolds, and biologically active substances to create or repair tissues that restore, maintain, or improve body function.

\medterm{Tissue Expander} Balloon implanted under skin/muscle, gradually filled with saline to expand tissue. Used
for breast reconstruction, scalp reconstruction, burn coverage. Removed after expansion,
replaced with permanent implant or flap.

\medterm{Tissue Expansion} Gradual stretching of skin or soft tissue with an implanted or external expander so additional tissue can be created for reconstruction; it usually requires staged treatment and monitoring for infection, pain, or device problems.

\medterm{Tissue Fixation} Preserving a tissue sample soon after collection so its cells and structure do not break down before laboratory examination. Chemical fixatives are commonly used, and the method must match the planned tests because poor or delayed fixation can affect the results.

\medterm{Tissue Kallikrein} A serine protease that cleaves kininogen to release kinins such as bradykinin, influencing local blood flow, vascular permeability, pain, and inflammation.

\medterm{Tissue Membrane} A thin sheet or boundary of tissue that separates compartments, covers surfaces, or surrounds structures, with composition and function varying by epithelial, connective, muscle, or biologic membrane type.

\medterm{Tissue Plasminogen Activator} An enzyme that converts plasminogen to plasmin, helping dissolve fibrin in blood clots. Recombinant forms are used as thrombolytic medicines in selected emergencies.

\medterm{Tissue Plasminogen Activator-Associated Angioedema} Angioedema occurring as a complication of tissue plasminogen activator therapy for acute
ischemic stroke or myocardial infarction, affecting approximately 1 to 5 percent of
patients. It typically involves the face, tongue, and upper airway and may be
life-threatening.

\medterm{Tissue Processor} A laboratory instrument that moves tissue through fixation, dehydration, clearing, and paraffin infiltration before embedding and microscopic sectioning.

\medterm{Tissue Scaffold} A natural or engineered three-dimensional framework that supports cell attachment, growth, organization, and tissue repair in regenerative medicine.

\medterm{Tissue Transglutaminase (tTG)} An enzyme that cross-links proteins and serves as the principal autoantigen in celiac disease; IgA antibodies against tissue transglutaminase are commonly used in celiac disease evaluation.

\medterm{Tissue Typing} Testing tissue antigens, especially HLA, to assess compatibility for transplantation.
\medterm{Tissues Of The Body} Groups of similar cells and extracellular material performing a common function: epithelial, connective, muscle, and nervous tissue.

\medterm{Titanium} A strong, lightweight, corrosion-resistant metal used in implants, prostheses, dental devices, and surgical instruments. Bone can form a stable connection with some titanium implant surfaces.

\medterm{Titer} The concentration or dilution level at which a measurable antibody, antigen, or other activity remains detectable in a specimen.
expressed as the highest dilution that still produces a detectable reaction (e.g.,
1:128).

\medterm{Titration} Gradual adjustment of a medicine dose or measurement of a substance using a reagent of known concentration.
\medterm{Titre} A measure of antibody or another substance's concentration, determined by repeatedly diluting a sample until detection stops. Titres can help assess infection, immunity, or response, but interpretation depends on the test and timing.
\medterm{Titubation} A rhythmic nodding or tremulous movement, especially of the head or trunk, often associated with cerebellar disease.
\medterm{Tizanidine} A muscle-relaxing medicine used for painful spasms and spasticity; it can cause sleepiness and low blood pressure.
\medterm{Thin-Layer Chromatography} A laboratory method that separates substances on a thin layer of silica, alumina, or another material as a solvent moves across it. It can help identify compounds, compare samples, or measure radiochemical purity.

\synonyms
TLC

\medterm{Tmax} The time after a dose when a medicine reaches its highest measured concentration in blood or plasma.

\medterm{TMD} Alternate index label; see \textit{TMJ disorders}.

\medterm{TMJ Disorders} Temporomandibular joint disorders affect the jaw joints or chewing muscles and may cause jaw pain, clicking, stiffness, headaches, or limited movement.

\synonyms
Temporomandibular Disorders
Temporomandibular Joint Disorders
TMD

\medterm{TNBC} Alternate index label; see \textit{Triple-negative breast cancer}.

\medterm{TNF Alpha} Tumor necrosis factor alpha is a potent pro-inflammatory cytokine produced primarily by
activated macrophages, as well as by T cells, NK cells, mast cells, and endothelial
cells. TNF-alpha exists as a 26-kilodalton transmembrane protein that is cleaved by the
metalloprotease TACE to release a soluble 17-kilodalton active form.

\medterm{TNM Classification} A cancer-staging system describing the primary tumor's size or spread, nearby lymph-node involvement, and distant metastases. Combined findings help estimate disease extent, guide treatment, compare outcomes, and discuss prognosis.
\medterm{TNM Staging} A cancer-staging system describing the primary tumor, regional lymph nodes, and distant metastases to estimate disease extent and prognosis.

\medterm{Tobacco} A plant product containing nicotine and many toxic combustion or processing chemicals. Its use causes dependence and increases the risk of cancer, cardiovascular disease, and lung disease.
\medterm{Tobacco Chewing (Smokeless Tobacco)} Alternate index label for Smokeless Tobacco.

\medterm{Tobacco Use} Tobacco use exposes the body to nicotine and toxic combustion or aerosol chemicals, causing addiction and increased risks of cancer, vascular disease, and lung disease.

\medterm{Tobramycin} An aminoglycoside antibiotic used for selected serious bacterial infections, including some Gram-negative infections. It can damage the kidneys or hearing and is used with dosing and laboratory monitoring when appropriate.
\medterm{Tocolysis} The use of medication to delay preterm birth for a limited time, often to allow corticosteroid treatment
or transfer to a facility with appropriate neonatal care. The drug, gestational-age range,
and contraindications depend on local guidance and the clinical situation.

\medterm{Tocolytic Agents} Tocolytic agents suppress uterine contractions temporarily to delay preterm birth, often allowing time for corticosteroids or transfer to specialist care.

\medterm{Tocolytic Therapy} Medications used to inhibit uterine contractions temporarily to prolong pregnancy,
allowing time for corticosteroid administration and transport to a tertiary care center.
Indications: preterm labor with intact membranes, 24-34 weeks.

\medterm{Tocolytics} Medications used to suppress preterm labor (contractions before 37 weeks gestation), to
delay delivery for corticosteroid administration (fetal lung maturity) and in utero
transport to a tertiary care center.

\medterm{Tocopherol} A form of vitamin E, a fat-soluble antioxidant involved in protecting cell membranes; deficiency is uncommon but can cause neurologic or hematologic abnormalities.

\medterm{Tocotrienols} Vitamin E compounds with unsaturated side chains found in some plant oils and grains; proposed cardiovascular and metabolic benefits remain under study.

\medterm{Toddler} A toddler is a young child, usually about one to three years old, developing rapidly in movement, language, behavior, feeding, and independence.

\medterm{Toddler Test Or Procedure Preparation} Preparing a toddler for a test or procedure uses simple explanations, play, distraction, comfort objects, caregiver support, and any required fasting or medicine instructions.


\medterm{Toddler'S Fracture} A small spiral crack in the lower shinbone of a young child, usually after a minor twist. The child may suddenly limp or refuse to walk, even though there is little swelling or deformity. The first X-ray may look normal; a repeat X-ray after about two weeks may show healing. Treatment commonly uses a supportive boot or cast, based on examination and imaging.

\synonyms
Cast-Tibia Fracture; Childhood Spiral Tibial Fracture

\medterm{Toe Broken (Broken Toe)} Alternate index label for Broken Toe.

\medterm{Toe Joint} Any joint connecting the bones of a toe. Toe joints include the metatarsophalangeal and interphalangeal joints; they permit toe movement and may be affected by arthritis, gout, injury, or deformity.

\synonyms
toe articulation.

\medterm{Toe Pressure} The blood pressure measured in a toe, useful for assessing circulation in the foot. It can remain reliable when ankle arteries are too stiff to compress, as may occur with diabetes or kidney disease.

\medterm{Toe-Brachial Index} The ratio of systolic toe pressure to brachial systolic pressure. It can support
diagnosis of peripheral arterial disease when ankle arteries are noncompressible because
of medial arterial calcification.

\medterm{Toenail} A toenail is the hard keratin plate covering the end of a toe and protecting its underlying tissue; injury, fungus, and ingrowth can cause disease.

\medterm{Toenail Fungus} Alternate index label; see \textit{Nail fungus}.

\medterm{Toes (Digits Of The Foot)} The five digits at the end of each foot. The big toe has two bones and helps with balance and push-off; the other toes usually have three bones each. Injury, arthritis, tight shoes, or nerve and muscle problems can cause pain or deformity.
\medterm{Togaviridae} A family of viruses that includes some mosquito-borne viruses, such as chikungunya and eastern equine encephalitis virus, which can cause fever, rash, joint pain, or brain inflammation.

\medterm{Togaviridae Infections} Infections caused by viruses in the Togaviridae family. Symptoms vary by virus and may include fever, rash, joint pain, or encephalitis.

\medterm{Toilet} A fixture or place used for urination and defecation; in healthcare, toilet access and toileting assistance are important parts of continence and fall prevention.

\medterm{Toilet Bowl Cleaners And Deodorizers Poisoning} Poisoning from swallowing or contacting toilet cleaners or deodorizers. Acidic or alkaline products can burn the mouth, throat, eyes, and skin; mixing cleaners may release dangerous gases.

\medterm{Tolazamide} An older sulfonylurea diabetes medicine that increases insulin release; hypoglycemia and weight gain can occur, especially in older adults or people with kidney disease.

\medterm{Tolazomide} An older diabetes medicine that lowers blood sugar by stimulating the pancreas to release insulin; now rarely used.
\medterm{Tolbutamide} An older first-generation sulfonylurea that lowers blood glucose by stimulating pancreatic insulin release. It is now rarely used and can cause low blood sugar and other adverse effects.
\medterm{Tolcapone} A catechol-O-methyltransferase inhibitor used with levodopa for Parkinson disease; potentially fatal liver injury limits use and requires monitoring.

\medterm{Tolerable Upper Intake Level} The highest average daily intake of a nutrient unlikely to cause harm in almost all healthy people. It includes food, drinks, and supplements and is not a recommended amount to consume.
\synonyms
UL

\medterm{Tolerance} A decreased response to a drug after repeated administration, requiring dose escalation to achieve the same effect. Unlike tachyphylaxis, tolerance develops over weeks to months of chronic use.
\medterm{Toll-Like Receptor} A class of single-pass membrane-spanning pattern recognition receptors (PRRs) expressed
on innate immune cells (macrophages, dendritic cells). Recognize evolutionary conserved
microbial structures (PAMPs).

\medterm{Toll-Like Receptors} Receptors of the innate immune system that recognize characteristic molecules from microbes or damaged tissue. Their activation starts inflammatory and antimicrobial responses.

\synonyms
TLRs

\medterm{Tolmetin} A nonsteroidal anti-inflammatory drug used for pain and arthritis; gastrointestinal bleeding, kidney injury, and cardiovascular events are important risks.

\medterm{Tolmetin Overdose} Taking too much tolmetin, a nonsteroidal anti-inflammatory medicine, can cause stomach bleeding, vomiting, kidney injury, drowsiness, seizures, or other serious effects.

\medterm{Tolnaftate} A topical antifungal medicine used for some dermatophyte infections such as athlete's foot and ringworm; it is not effective for every type of fungal infection.

\medterm{Tolonium Chloride} Tolonium chloride, or toluidine blue, is a dye used in laboratory and selected clinical procedures to highlight cells or tissue changes.

\medterm{Tolterodine Tartrate} Tolterodine tartrate is an antimuscarinic medicine used for overactive bladder and urinary urgency; dry mouth, constipation, and urinary retention can occur.

\medterm{Toluene} Toluene is a volatile solvent that can irritate tissues and, with substantial inhalation or ingestion, injure the nervous system, liver, kidneys, or heart.

\medterm{Toluene And Xylene Poisoning} Poisoning from solvents containing toluene or xylene, often through swallowing or inhalation. It can cause dizziness, confusion, abnormal heart rhythms, breathing problems, or organ injury.

\medterm{Tolvaptan} A vasopressin V2-receptor antagonist used in selected patients with hyponatremia or autosomal dominant polycystic kidney disease; thirst, rapid sodium correction, and liver injury require attention.

\medterm{Tomogram} A tomogram is an image of a selected cross-sectional layer of the body or another object, produced by tomography such as CT or MRI.

\medterm{Tomography} Imaging that produces cross-sectional views, or slices, through the body. Computed tomography uses X-rays, while other forms use ultrasound, magnetic fields, or radioactive tracers.
\medterm{Tomosynthesis} An imaging test that takes several low-dose X-ray pictures from different angles and combines them into thin layers. It is commonly used to examine breast tissue in three dimensions.

\medterm{Tomotherapy} Tomotherapy combines CT imaging with rotating, precisely shaped radiation beams to deliver radiation treatment while adjusting for the target and nearby tissues.

\medterm{Tongue} A muscular organ involved in taste, speech, chewing, and swallowing.
\medterm{Tongue Biopsy} A tongue biopsy removes a small tissue sample from a tongue lesion for microscopic examination of infection, inflammation, precancer, or cancer.

\medterm{Tongue Carcinoma} Tongue carcinoma is cancer of the tongue, usually squamous-cell cancer, and may cause a persistent sore, lump, pain, or difficulty speaking or swallowing.

\medterm{Tongue Disease} Any disorder affecting the tongue, such as infection, inflammation, injury, nerve disease, nutritional deficiency, or cancer. Symptoms may include pain, sores, swelling, altered taste, or difficulty speaking or swallowing.
\medterm{Tongue Neoplasm} A tongue neoplasm is an abnormal growth in the tongue that may be benign, precancerous, or malignant.

\medterm{Tongue Problems} Changes in the tongue, such as pain, swelling, sores, color change, numbness, or difficulty moving or tasting. Causes include injury, infection, inflammation, nutritional deficiency, and other disease.

\medterm{Tongue Tie} A short or tight band of tissue under the tongue that limits tongue movement. It may interfere with breastfeeding, speech, or eating in some people.

\medterm{Tongue, Disorders Of} Conditions affecting the tongue, including infection, inflammation, injury, tumors, ulcers, movement problems, and developmental differences. They may cause pain, swelling, altered taste, speech difficulty, swallowing problems, or breathing obstruction.
\medterm{Tongue-Tie (Ankyloglossia)} Restricted movement of the tongue caused by an unusually short or tight lingual frenulum, sometimes affecting feeding, speech, or oral hygiene.

\synonyms
Ankyloglossia


\medterm{Tonic Neck Reflex} A primitive reflex present from birth to 6 months, also called fencer's pose. When the
head is turned to one side, the arm on that side extends while the opposite arm flexes.
Present in both directions. Asymmetry may indicate neurologic injury. Disappears as
voluntary head control develops.

\synonyms
fencer's pose.

\medterm{Tonic Pupil} A tonic pupil is a pupil with poor light response but a slow constriction to near focus, often caused by injury to postganglionic parasympathetic fibers.

\medterm{Tonic Seizures} Seizures causing sustained stiffening of muscles, usually for seconds to minutes. They may affect one area or the whole body, cause falls or breathing difficulty, and require individualized seizure evaluation.

\medterm{Tonic-Clonic Seizure} A generalized seizure with a tonic stiffening phase followed by rhythmic clonic jerking,
impaired consciousness, and a postictal state.

\medterm{Tonics} Preparations traditionally claimed to improve strength, energy, appetite, or general health. Their ingredients, safety, and evidence vary widely, and some may interact with medicines or contain harmful amounts of active substances.
\medterm{Tonometer} An instrument that measures pressure inside the eye. Eye-pressure readings help detect glaucoma or monitor treatment, but they must be interpreted with the optic nerve, visual field, cornea, and other findings.

\medterm{Tonometry} Measurement of pressure inside the eye, usually to assess for glaucoma or monitor its treatment. Methods include applanation, noncontact air-puff, rebound, and electronic devices; results must be interpreted with corneal thickness and the clinical context.

\medterm{Tonsil} A small mass of immune tissue at the back of the throat. Tonsils help detect germs entering through the mouth and nose but can become enlarged or infected, causing sore throat, difficulty swallowing, or breathing problems.
\synonyms
Tonsils

\medterm{Tonsil Inflammation (Adenoids And Tonsils)} Alternate index label for Adenoids and Tonsils.


\medterm{Tonsil Stones} Small, hard lumps formed when food particles, mucus, dead cells, and bacteria become trapped and calcify in the folds of the tonsils. They are often harmless but may cause bad breath, throat irritation, or a feeling that something is stuck in the throat. Good oral hygiene and gargling may help; recurrent symptoms, bleeding, or one-sided tonsil enlargement requires examination.

\medterm{Tonsillar Cancer} A malignancy arising from the palatine tonsils, most commonly squamous cell carcinoma
associated with HPV infection. It often presents with throat pain, otalgia, and cervical
lymphadenopathy.

\medterm{Tonsillar Neoplasm} A tonsillar neoplasm is an abnormal growth in a tonsil that may be benign or malignant and can cause a persistent sore throat, swallowing difficulty, or neck swelling.


\medterm{Tonsillectomy} Surgical removal of the palatine tonsils, performed for selected recurrent infection, obstructive sleep-disordered breathing, or other tonsillar disease.
apnea, or other selected conditions after assessment of potential benefits and risks.

\medterm{Tonsillitis} Inflammation or infection of the tonsils, commonly viral or streptococcal.

\medterm{Tonus} The normal resting tension or partial contraction of a muscle or tissue. Abnormally low or high tone can occur with nerve, muscle, brain, spinal-cord, or developmental disorders.
\medterm{Tooth} A hard mineralized structure anchored in the jaw, used for biting, tearing, chewing, and grinding food.
\medterm{Tooth - Abnormal Colors} Abnormal tooth color can result from plaque, decay, trauma, fluorosis, medicines, enamel defects, or staining, and the cause determines treatment.

\medterm{Tooth - Abnormal Shape} Abnormal tooth shape may result from developmental defects, genetic conditions, trauma, wear, or tumors and can affect function, appearance, and eruption.

\medterm{Tooth Abrasion} Loss of tooth substance caused by external mechanical friction, such as aggressive brushing, abrasive products, or habits, producing grooves, notches, or sensitivity.

\medterm{Tooth Abscess} A tooth abscess is a localized collection of pus from bacterial infection of the tooth or surrounding tissues, causing pain, swelling, and possible spreading infection.

\medterm{Tooth Anatomy} Tooth anatomy includes enamel, dentin, pulp, cementum, root, crown, periodontal ligament, and supporting bone and gums.

\medterm{Tooth Ankylosis} Fusion of a tooth root to alveolar bone through loss of the periodontal ligament, preventing normal eruption and making the tooth appear infraoccluded as neighboring teeth grow.

\medterm{Tooth Attrition} Wear of tooth surfaces caused by tooth-to-tooth contact, commonly from chewing, bruxism, or malocclusion, producing flattened cusps and possible sensitivity.

\medterm{Tooth Calcification} Deposition of mineral that forms and hardens enamel, dentin, cementum, and other tooth structures during development; abnormal timing or amount can affect tooth strength and appearance.

\medterm{Tooth Component} Tooth components include enamel, dentin, cementum, pulp, and supporting tissues, each contributing to tooth strength, sensation, nutrition, and attachment.

\medterm{Tooth Crowding} Tooth crowding occurs when teeth lack enough space to align normally, often from tooth size, jaw size, eruption pattern, or retained primary teeth.

\medterm{Tooth Crown} The tooth crown is the visible portion above the gum, covered mainly by hard enamel and containing underlying dentin and pulp.

\medterm{Tooth Decay} Gradual destruction of tooth structure caused when mouth bacteria make acids from sugars. It can cause sensitivity, pain, holes, infection, and tooth loss, but fluoride, cleaning, diet, and dental care help prevent progression.

\synonyms
Dental caries; cavities

\medterm{Tooth Decay - Early Childhood} Early-childhood tooth decay is bacterial acid destruction of primary teeth, promoted by frequent sugars, poor brushing, enamel defects, and prolonged bottle or nursing exposure.

\medterm{Tooth Discoloration} Change in tooth color caused by surface stains, dental decay, trauma, medications, fluorosis, aging, or changes within enamel or dentin.

\medterm{Tooth Disease} Any disorder affecting a tooth, including decay, infection, cracks, developmental defects, or wear. It may cause pain, sensitivity, discoloration, swelling, or difficulty chewing and should be assessed by a dental professional.
\medterm{Tooth Enamel} Tooth enamel is the highly mineralized outer covering of the crown that protects dentin and pulp from wear, acid, and bacteria.

\medterm{Tooth Eruption} Process by which tooth moves from intraosseous position to functional occlusion. Primary
(deciduous) dentition: 20 teeth, eruption 6-30 months. Permanent dentition: 32 teeth
(including third molars), eruption 6-21 years. Delayed eruption may indicate systemic
conditions, local obstruction, or normal variant.

\medterm{Tooth Extraction} Removal of a tooth from its socket because of severe decay, infection, periodontal disease, fracture, impaction, orthodontic planning, or another clinical indication.

\medterm{Tooth Formation - Delayed Or Absent} Delayed or absent tooth formation means teeth develop or erupt later than expected or fail to form, due to local, genetic, endocrine, nutritional, or systemic causes.

\medterm{Tooth Fracture} A crack or break in a tooth caused by injury, biting hard material, decay, grinding, or weakened enamel. It may expose the inner nerve, causing pain, sensitivity, infection, or loss of the tooth.
\medterm{Tooth Infection} Infection of the tooth, its pulp, or surrounding tissues, commonly caused by untreated decay, cracks, or gum disease.
It may cause toothache, sensitivity, swelling, fever, or drainage. Facial swelling, difficulty swallowing or
breathing, or spreading infection requires emergency assessment.

\medterm{Tooth Loss} Tooth loss is absence of a tooth due to decay, gum disease, trauma, developmental absence, or extraction and can affect chewing, speech, and oral health.

\medterm{Tooth Luxation} Traumatic displacement of a tooth without complete avulsion. Concussion and subluxation
involve periodontal injury with little or no displacement, whereas extrusive, lateral,
and intrusive luxation involve increasingly marked displacement and require prompt
dental evaluation.

\medterm{Tooth Replantation} Replacement of an accidentally displaced or avulsed tooth into its socket, followed by stabilization
and dental follow-up.

\medterm{Tooth Resorption} Loss of dental hard tissue caused by odontoclast activity, occurring internally or on the root surface and associated with trauma, infection, orthodontics, or idiopathic disease.

\medterm{Tooth Root} The tooth root anchors a tooth in its socket and contains dentin, cementum, pulp, nerves, and blood vessels beneath the gum.

\medterm{Tooth Sensibility Testing} Clinical testing of a tooth's response to thermal, electric, or dentin stimulation. It
assesses neural response rather than directly measuring pulp blood flow, so results must
be interpreted with history, examination, and imaging.

\medterm{Tooth Sensitivity} Pain in response to thermal, sweet, acidic, or tactile stimulus. Caused by dentin
exposure (recession, erosion, abrasion, abfraction). Management: desensitizing
toothpaste, fluoride varnish, restorative coverage, gum graft for recession.

\medterm{Tooth, Supernumerary} An extra tooth beyond the normal number that may crowd neighboring teeth, remain trapped, or interfere with normal eruption.
\medterm{Toothache} Pain arising from a tooth or nearby gum, ligament, bone, or other supporting tissue. Common causes include decay, infection, cracks, trauma, and gum disease.
\medterm{Toothache Overview} Toothache usually reflects decay, pulp inflammation, fracture, gum disease, abscess, trauma, or referred pain. Dental treatment is needed to correct the cause; facial swelling, fever, difficulty swallowing, or breathing difficulty requires urgent care.

\synonyms
Pain Tooth (Toothache Overview)

\medterm{Toothaches} Toothache is pain arising from decay, pulp inflammation, abscess, gum disease, trauma, cracking, or referred facial pain and requires dental assessment if persistent.

\medterm{Toothpaste Overdose} Swallowing too much toothpaste can cause stomach upset, vomiting, or diarrhea. Large amounts of fluoride toothpaste can cause serious salt and calcium disturbances, especially in children.

\medterm{Tophaceous Gout} Chronic tophaceous gout represents advanced gout with deposition of large, visible tophi
consisting of monosodium urate crystals surrounded by granulomatous inflammation. Tophi
typically develop after years of inadequately treated gout and occur on the ears,
fingers, toes, and olecranon bursa.

\medterm{Tophus} A deposit of monosodium urate crystals in gout, often around joints or in soft tissues.
\medterm{Topical} Applied directly to skin or a moist lining, usually to act at the treatment site with limited whole-body exposure.

\medterm{Topical Antibiotics} Antibiotics applied to the skin to treat selected superficial bacterial infections. The choice and
duration depend on the organism, site, extent of infection, resistance patterns, and local
guidance. Unnecessary or prolonged use can cause irritation and antimicrobial resistance.

\medterm{Topical Antifungals} Topical antifungals treat superficial fungal infections. Allylamines (terbinafine,
naftifine) inhibit squalene epoxidase, fungicidal. Azoles (clotrimazole, miconazole,
ketoconazole, econazole) inhibit ergosterol synthesis, fungistatic. Ciclopirox has
broad-spectrum activity. Tolnaftate is effective against dermatophytes.

\medterm{Topical Calcineurin Inhibitors} Non-steroidal anti-inflammatory agents used for atopic dermatitis and facial eczema.
They inhibit T-cell activation by blocking calcineurin, preventing cytokine production.
Advantages over corticosteroids: no skin atrophy, safe for face and intertriginous
areas. Side effects include transient burning and itching at application site.

\medterm{Topical Corticosteroid Potency} The strength of a steroid cream, ointment, lotion, or gel used on the skin. Mild products are used for sensitive areas or children, while stronger products may be needed briefly for thick or severe inflammation. Choice depends on the diagnosis, body site, age, duration, and risk of skin thinning or other effects.

\medterm{Topical Retinoids} Vitamin A derivatives used for acne, photoaging, and hyperpigmentation. Tretinoin
(all-trans retinoic acid) is the most studied. Adapalene is more stable and less
irritating. Tazarotene is the most potent. Mechanism: normalizes follicular
keratinization, increases cell turnover, and stimulates collagen production.

\medterm{Topiramate} Topiramate is an antiseizure medicine also used to prevent migraine; cognitive slowing, paresthesias, weight loss, kidney stones, metabolic acidosis, and fetal oral-cleft risk are important considerations.

\medterm{Topoisomerase} An enzyme that temporarily cuts and rejoins DNA to control its twisting during replication, transcription, and chromosome separation.

\medterm{Topoisomerase Ii} An enzyme that cuts both DNA strands, passes another DNA segment through the break, and rejoins the strands; it is targeted by some anticancer drugs.

\medterm{Topoisomerase Inhibitors} Drugs that interfere with DNA topoisomerases; some are antibiotics and others are anticancer medicines, with toxicity depending on the drug and target enzyme.

\medterm{Topotecan} A topoisomerase-I inhibitor used for selected ovarian, small-cell lung, and cervical cancers; bone-marrow suppression is a major toxicity.

\medterm{Topotecan Hydrochloride} A salt form of topotecan, a chemotherapy medicine that interferes with DNA copying. It treats selected cancers but can cause dangerously low blood counts, infection, bleeding, diarrhea, nausea, or fatigue.

\medterm{TORCH Panel} Serologic screening for congenital infections: Toxoplasma, Rubella, CMV, Herpes simplex
(and others: syphilis, parvovirus B19, VZV, Zika). IgM: indicates recent/active
infection (maternal or fetal). IgG: past infection or passive immunity. Paired maternal
and neonatal testing.

\medterm{TORCH Screen} Testing for selected congenital infections—classically Toxoplasma, Other infections, Rubella, Cytomegalovirus, and Herpes—when fetal or newborn findings suggest an in-utero infection; broad panels require careful interpretation.

\medterm{Toremifene} A selective estrogen-receptor modulator used for metastatic breast cancer in postmenopausal women; hot flashes and, rarely, dangerous QT prolongation may occur.

\medterm{Torn ACL} A torn anterior cruciate ligament is a knee injury commonly caused by sudden pivoting, landing, or contact, producing a pop, swelling, instability, and difficulty changing direction. Rehabilitation, bracing, or reconstruction is selected according to activity, associated injury, and goals.

\medterm{Torn Meniscus} A torn meniscus is a split in the knee's shock-absorbing cartilage caused by twisting or degeneration. Pain, swelling, catching, or locking may occur; physical therapy or arthroscopic treatment depends on tear pattern, symptoms, age, and associated disease.

\synonyms
Meniscus Tear

\medterm{Torpor} A state of greatly reduced activity, alertness, or responsiveness. In a patient, marked or new torpor may signal a serious illness, medicine effect, or metabolic problem.
\medterm{Torsade De Pointes} Torsade de pointes is a potentially fatal polymorphic ventricular tachycardia associated with prolonged QT, electrolyte disturbance, medicines, or inherited channel disorders.

\medterm{Torsades De Pointes} A specific form of polymorphic ventricular tachycardia with QRS complexes that twist
around the baseline, occurring in the setting of QT prolongation. It may degenerate to
ventricular fibrillation.

\medterm{Torsemide} Torsemide is a loop diuretic that inhibits sodium-potassium-chloride reabsorption in the thick ascending limb and treats oedema from heart, kidney, or liver disease; dehydration and electrolyte loss may occur.

\medterm{Torsion} Twisting of a body structure around its supporting axis, which may reduce or stop blood flow. Torsion of the testicle, ovary, or bowel can destroy tissue and often requires emergency surgery.
\medterm{Torsion Dystonia} A movement disorder causing sustained or repetitive muscle contractions that twist the body or limbs into abnormal postures; it may be inherited, acquired, focal, or generalized.

\medterm{Torsion Of The Appendix Testis} Twisting of a tiny leftover structure attached to the upper part of the testicle. It is common in boys before puberty and causes sudden, localized scrotal pain and tenderness; a small blue spot may be visible through the skin. The testicle usually still has normal blood flow. Rest, support, and pain medicine usually allow it to heal, but urgent assessment is needed to exclude testicular torsion.

\synonyms
Hydatid of Morgagni Torsion; Appendix Testis Torsion

\medterm{Torsion Testicle (Testicular Disorders)} Alternate index label for Testicular Disorders.

\medterm{Torsion, Testicular} Alternate index label; see \textit{Testicular torsion}.

\medterm{Torticollis} An abnormal neck position caused by muscle tightening, pain, nerve problems, injury, or birth-related conditions. The head may tilt or rotate, and treatment depends on the cause, age, and severity.
\medterm{Torticollis (Wryneck)} A condition in which the head is tilted to one side (ear toward the shoulder) and
rotated to the opposite side (chin toward opposite shoulder), due to contraction or
shortening of the sternocleidomastoid muscle.

\synonyms
Wryneck

\medterm{Torticollis, Spasmodic} Alternate index label; see \textit{Cervical dystonia}.

\medterm{Torulaspora} A genus of yeasts used in food and fermentation; human infection is rare and usually opportunistic.

\medterm{Torulaspora Delbrueckii} A yeast used in wine and food fermentation; it rarely causes infection in people with major immune compromise.

\medterm{Torus Fracture} A stable incomplete break in a child's bone caused by compression, often in the wrist after a fall. The bone bulges rather than breaking all the way through; it usually heals well with a splint or cast and follow-up as advised.

\synonyms
Buckle Fracture; Metaphyseal Buckle Fracture; Torus Fracture of Radius

\medterm{Torus Mandibularis} A bony outgrowth arising from the lingual surface of the mandible, typically bilateral
and located in the premolar region above the mylohyoid ridge. It is a common,
asymptomatic developmental anomaly, more prevalent in women and certain ethnic
populations.

\medterm{Torus Palatinus} A bony exostosis arising from the midline of the hard palate, typically presenting as a
slow-growing, asymptomatic, hard, bony protuberance discovered incidentally. It is more
common in women and certain ethnic groups.

\medterm{Tositumomab} A monoclonal antibody formerly used with radioactive iodine-131 for some relapsed or refractory B-cell non-Hodgkin lymphomas; it is no longer routinely available in many settings.

\medterm{Total Abdominal Colectomy} Surgical removal of the entire colon through an abdominal approach, with restoration of intestinal continuity or creation of an ileostomy as clinically appropriate.

\medterm{Total Anomalous Pulmonary Venous Connection} Alternate index label; see \textit{Total anomalous pulmonary venous return}.

\medterm{Total Anomalous Pulmonary Venous Return} A birth defect in which the veins returning oxygen-rich blood from the lungs connect to the wrong side of the heart. Babies may have blue skin, rapid breathing, poor feeding, or heart failure. Severe obstruction is an emergency, and surgery is needed to reroute the veins.

\medterm{Total Artificial Heart} A mechanical device that replaces the pumping function of both ventricles and is used in selected patients with severe biventricular heart failure, usually as a bridge to transplantation. It requires specialized center-based management and carries risks including bleeding, infection, thrombosis, and device malfunction.

\medterm{Total Body Surface Area} The percentage of body surface involved by a burn or other extensive skin process. It is
estimated using methods such as the rule of nines, Lund-Browder chart, or the patient's
palm and is used to guide burn severity assessment and resuscitation.

\medterm{Total Breech Extraction} A maneuver where the operator grasps the fetal feet and delivers the entire baby in
breech presentation. The most common indication is delivery of the second twin
(malpresentation). Contraindicated for term singleton breech due to high risk of head
entrapment, nuchal arms, and birth trauma.

\medterm{Total Cholesterol} Total cholesterol is the combined cholesterol carried in major blood lipoproteins; cardiovascular risk is assessed with HDL, LDL, triglycerides, and other factors.


\medterm{Total Daily Energy Expenditure} Total calories burned in 24 hours: BMR + thermic effect of food + physical activity.
Estimation guides dietary recommendations. Activity factor multiplies BMR.

\medterm{Total Haemolytic Complement Test} A blood test measuring the overall activity of the classical complement pathway, commonly reported as CH50. Abnormal results can suggest immune deficiency, autoimmune disease, infection, inflammation, or laboratory interference.

\medterm{Total Hysterectomy} Removal of the uterus and cervix. The fallopian tubes and ovaries are not necessarily removed and should be specified separately, such as salpingectomy or oophorectomy.

\medterm{Total Hip Arthroplasty} Replacement of the hip joint with prosthetic components. Acetabular component: metal
shell + polyethylene liner. Femoral component: metal stem + femoral head (cobalt-chrome
or ceramic). Fixation: cemented (PMMA, for elderly, osteoporotic bone), cementless
(porous-coated, for younger patients, biologic fixation), and hybrid.

\medterm{Total Hip Joint Replacement - Revision} Surgery to remove and replace some or all parts of a previous hip replacement that have loosened, worn out, become infected, dislocated repeatedly, or failed. It is more complex than first-time replacement and may require bone reconstruction.

\medterm{Total Intravenous Anesthesia} General anesthesia given entirely through medicines injected into a vein rather than inhaled anesthetic gases. It allows carefully controlled unconsciousness during surgery, but can lower blood pressure or breathing, so trained anesthesia staff must continuously monitor the patient.

\medterm{Total Iron Binding Capacity} TIBC: a measure of the maximum iron-binding capacity of transferrin in serum, reflecting
total iron-transport capacity. Calculated from transferrin or measured directly.

\medterm{Total Knee Arthroplasty} Replacement of the knee joint with prosthetic components. Femoral component
(cobalt-chrome), tibial component (titanium tray + polyethylene insert), and patellar
component (polyethylene, optional). Fixation: cement (PMMA) most common in older adults;
cementless in younger.

\medterm{Total Knee Joint Replacement - Revision} Surgery to remove and replace parts of a previous knee replacement because of loosening, wear, infection, instability, stiffness, or persistent pain. It may require specialized implants or bone grafting and has higher risks than the first operation.

\medterm{Total Knee Replacement} Surgery that replaces damaged surfaces of the knee with artificial parts to relieve severe pain and improve function. Recovery includes exercises and gradual activity; infection, blood clots, stiffness, loosening, persistent pain, or repeat surgery may occur.

\medterm{Total Laryngectomy} Total laryngectomy removes the entire voice box and creates a permanent opening in the neck for breathing, usually to treat advanced laryngeal cancer.

\medterm{Total Lesion Glycolysis} A measurement from a PET scan that combines the active volume of a tumor with its average tracer uptake. It estimates the tumor’s total metabolic activity and may help assess treatment response or prognosis, but it is not interpreted alone.

\medterm{Total Lung Capacity} The total volume of air in the lungs after maximal inspiration. It is measured by
lung-volume testing, such as body plethysmography, and may be increased in hyperinflation
or decreased in restrictive lung disease.

\medterm{Total Pancreatectomy} Total pancreatectomy removes the entire pancreas and causes permanent loss of insulin and digestive-enzyme production, requiring lifelong replacement treatment.

\medterm{Total Parenteral Nutrition} Intravenous administration of nutrients sufficient to meet essentially all nutritional
requirements when the gastrointestinal tract cannot be used adequately.

\medterm{Total Parenteral Nutrition - Infants} Intravenous nutrition for infants who cannot receive enough nutrients through the gastrointestinal tract; it requires careful monitoring of glucose, electrolytes, liver function, infection, and catheter complications.

\medterm{Total Proctocolectomy And Ileal-Anal Pouch} Total proctocolectomy with an ileal-anal pouch removes the colon and rectum and creates a pouch from the ileum that is connected to the anus.

\medterm{Total Proctocolectomy With Ileostomy} Total proctocolectomy with ileostomy removes the colon and rectum and diverts the end of the small intestine through an abdominal stoma.

\medterm{Total Protein} The combined amount of albumin and other proteins in the blood. A blood test may help assess nutrition, liver or kidney disease, inflammation, or abnormal antibody production.

\medterm{Total Sleep Time} The total amount of time spent asleep during a defined sleep period or sleep study, excluding periods of wakefulness.

\medterm{Touch} The sensory perception of contact, pressure, texture, vibration, or movement across the skin. Signals travel through sensory nerves to the spinal cord and brain.
\medterm{Tourette Disorder} A neurodevelopmental disorder characterized by multiple motor tics and one or more vocal
tics lasting >1 year, with onset before age 18. Motor tics: eye blinking, head jerking,
shoulder shrugging, facial grimacing. Vocal tics: throat clearing, sniffing, grunting,
barking, coprolalia (obscene words, occurs in <10\%).

\medterm{Tourette Syndrome} A neurodevelopmental disorder characterized by multiple motor tics and at least one
vocal tic occurring over at least one year. Tics are sudden, rapid, recurrent,
stereotyped movements or vocalizations. They are often preceded by premonitory urges.
Tics typically worsen in adolescence and may improve in adulthood.

\synonyms
Tourette's

\medterm{Tourette'S} Alternate index label; see \textit{Tourette syndrome}.

\medterm{Tourette'S Syndrome} A neurodevelopmental disorder involving multiple motor tics and at least one vocal tic persisting for more than one year.
\medterm{Tourniquet} A tight band or device placed around a limb to temporarily reduce blood flow. It can control life-threatening bleeding or provide a bloodless surgical field, but prolonged or incorrect use can damage nerves and tissue.
\medterm{Towel Clamp} A surgical clamp used to secure drapes or hold materials in place during an operation. Its sharp points must be applied carefully because they can puncture skin or damage underlying tissue.

\medterm{Townes Syndrome} A rare inherited birth disorder that commonly causes blockage or absence of the anal opening and unusual thumbs. Kidney, heart, hearing, foot, or genital abnormalities may also occur. The specific problems vary, so treatment is planned around each child’s findings.

\medterm{Toxaemia} An obsolete term once used for illness thought to be caused by toxins in the blood. In pregnancy, “toxaemia” was formerly used for pre-eclampsia, which is now diagnosed by high blood pressure and related organ findings.
\medterm{Toxaphene} A persistent organochlorine insecticide mixture that can accumulate in soil, water, and body fat. Significant exposure may overstimulate the nervous system and cause tremors, seizures, liver injury, or other toxic effects.
\medterm{Toxbase} A clinical poison-information service that helps healthcare professionals assess chemical or medicine exposures. It provides evidence-based advice about expected symptoms, necessary tests, treatment, observation, and urgency for particular substances.
\medterm{Toxic} Poisonous or capable of harming living cells, tissues, organs, or the whole body. Toxicity depends on the substance, dose, route, duration, and individual susceptibility.

\medterm{Toxic Dose} The amount of a substance that causes harmful effects. It varies with the substance, route, duration of exposure, age, health, and individual susceptibility.

\medterm{Toxic Effect} Harmful change caused by a chemical, medicine, poison, radiation, or other exposure. The effect may involve one or more organs and depends on the dose, route, duration, metabolism, and individual susceptibility.
\medterm{Toxic Encephalopathy} Brain dysfunction caused by exposure to a toxic substance or medicine. Symptoms may include
confusion, altered behavior, seizures, or reduced consciousness.

\medterm{Toxic Epidermal Necrolysis} A severe, potentially fatal mucocutaneous reaction, usually to a medication,
causing extensive epidermal necrosis and skin detachment. It is on the severe end of the Stevens--Johnson syndrome and
toxic epidermal necrolysis spectrum and requires urgent specialist care.

\synonyms
TEN

\medterm{Toxic Megacolon} Acute life-threatening dilation of the colon associated with systemic toxicity,
representing one of the most serious complications of inflammatory bowel disease,
particularly ulcerative colitis. It can also occur in infectious colitis caused by
Clostridioides difficile and cytomegalovirus.

\medterm{Toxic Metabolic Encephalopathy} Diffuse brain dysfunction caused by systemic metabolic or toxic disturbances, producing confusion, altered alertness, agitation, seizures, or coma.

\medterm{Toxic Nodular Goiter} An enlarged thyroid with nodules that produce excessive thyroid hormone.

\medterm{Toxic Shock Syndrome} A severe toxin-mediated illness, usually caused by Staphylococcus aureus or Streptococcus pyogenes, with fever, rash, hypotension, and organ dysfunction.
\medterm{Toxic Synovitis} Temporary inflammation of the lining of the hip joint, usually in a child after a viral illness. It can cause hip pain and limping and must be distinguished from a serious joint infection.

\medterm{Toxicity} The ability of a substance, treatment, or exposure to cause harmful effects in a living organism. Severity depends on the dose, route, duration, metabolism, target organs, and individual susceptibility; effects may be temporary or permanent.

\medterm{Toxicity Grade} A standardized rating of how severe a harmful effect is, often used for adverse effects of medicines or cancer treatment. Grades generally range from mild to life-threatening and help guide monitoring and management.
\medterm{Toxicity Test} A laboratory or clinical study that examines whether a substance causes harmful effects and identifies affected organs, dose ranges, and exposure duration. Results help estimate safety but cannot predict every person's response.

\medterm{Toxicogenomics} The study of how chemicals or other harmful exposures alter gene activity and cellular pathways. It helps researchers identify mechanisms of toxicity, possible biomarkers, susceptible populations, and potential target organs.
\medterm{Toxicokinetic Description} An account of how a harmful substance enters the body, is absorbed, distributed, changed by metabolism, and eliminated. These processes determine the concentration and duration of exposure at target organs.
\medterm{Toxicokinetics} The study of absorption, distribution, metabolism, and excretion of a toxic substance
and its relationship to concentration over time. Toxicokinetics helps estimate exposure,
timing, accumulation, and expected duration of toxicity.

\medterm{Toxicology} The study of poisons and harmful effects of chemicals, medicines, and other substances. It includes how exposure happens, how the body is affected, how poisoning is recognized, how it can be prevented, and how it is treated.
\medterm{Toxicology Screen} A test of urine, blood, or another sample for medicines, drugs, alcohol, or poisons. Results must be interpreted with the person's symptoms, timing, and prescribed medicines.

\medterm{Toxicology Specimen Preservation} Collection, labeling, storage, and transport of biological specimens to preserve
analytes and prevent contamination or degradation. Appropriate specimens may include
peripheral blood, urine, vitreous humor, bile, liver, hair, or gastric contents
depending on the suspected exposure and timing.

\medterm{Toxidrome} A recognizable pattern of signs and symptoms produced by a class of toxins. Important
patterns include opioid, sympathomimetic, anticholinergic, cholinergic,
sedative-hypnotic, and serotonin toxidromes; treatment combines supportive care,
targeted antidotes when available, and decontamination decisions.

\medterm{Toxin} A poisonous substance made by a living organism, such as a bacterium, plant, animal, or fungus. Toxins can damage specific tissues or disrupt nerve, muscle, or cellular functions.

\synonyms
Biological poison

\medterm{Toxin Ciguatera (Ciguatera Poisoning)} Alternate index label for Ciguatera Poisoning.

\medterm{Toxins} Poisonous substances produced by living organisms, including bacteria, fungi, plants, and animals. They can injure cells or disrupt organ function; the term is usually distinguished from toxicants, which are harmful substances of nonbiological origin.
\medterm{Toxocara} A genus of parasitic roundworms that commonly infect dogs and cats. People acquire infection by swallowing eggs from contaminated soil or food; migrating larvae can cause fever, enlarged organs, or eye or nervous-system disease.
\medterm{Toxocara Canis} A dog roundworm whose eggs can infect people. Larvae migrate through tissues and may cause toxocariasis, with fever, cough, enlarged liver, eye inflammation, or rarely nervous-system involvement.
\medterm{Toxocara Cati} A cat roundworm whose eggs can infect people. The larvae may migrate through organs and cause toxocariasis, including fever, cough, enlarged liver, eye inflammation, or rarely neurologic disease.
\medterm{Toxocariasis} Infection by dog roundworm (Toxocara canis) or cat (T. cati), acquired by ingesting infective eggs from soil contaminated by animal feces (pica in children a risk).
\medterm{Toxoid} A bacterial toxin that has been chemically or otherwise inactivated so it cannot cause its usual harm but can still stimulate protective immunity. Toxoids are used in vaccines against diseases such as tetanus and diphtheria.
\medterm{Toxoid Vaccines} Vaccines that protect against harmful bacterial toxins. These vaccines contain toxins
that have been detoxified and rendered inactive.

\medterm{Toxoplasma} A genus of intracellular parasites; Toxoplasma gondii is the species that infects people. Infection is often mild, but it can seriously affect a fetus or people with weakened immunity, especially the brain and eyes.
\medterm{Toxoplasma Blood Test} A blood test detecting antibodies or genetic material related to Toxoplasma gondii infection, interpreted with symptoms, immune status, pregnancy, and timing of exposure.

\medterm{Toxoplasma Gondii} An obligate intracellular protozoan parasite that causes toxoplasmosis. The definitive
host is the cat (family Felidae), which sheds environmentally resistant oocysts in
feces.

\medterm{Toxoplasmosis} An infection caused by the parasite Toxoplasma gondii, usually mild but potentially serious during pregnancy or weakened immunity.

\synonyms
Chorioretinitis Toxoplasma (Toxoplasmosis)
Meningoencephalitis Toxoplasma (Toxoplasmosis)

\medterm{Toxoplasmosis (Toxo)} Infection with Toxoplasma gondii, often mild in healthy people but potentially severe in pregnancy or in people with weakened immunity.

\medterm{TPA} Tissue plasminogen activator, a natural clot-dissolving protein or a laboratory-made medicine based on it. In selected emergencies it helps break down clots, but it can cause serious bleeding.

\synonyms
Tissue plasminogen activator

\medterm{TPMT Polymorphism} Thiopurine S-methyltransferase is an enzyme that metabolizes thiopurine drugs including
azathioprine, 6-mercaptopurine, and thioguanine, which are used in the treatment of
leukemia, inflammatory bowel disease, and autoimmune conditions.

\medterm{TPN} TPN, or total parenteral nutrition, supplies nutrients intravenously when the gastrointestinal tract cannot be used adequately.

\medterm{Tra Gene} A gene involved in sex determination and sexual differentiation in some organisms, especially insects; its function and sequence vary by species.

\medterm{Trabectedin} Trabectedin is an antineoplastic drug that binds DNA and alters transcription and repair, used for selected advanced soft-tissue sarcomas and ovarian cancer.

\medterm{Trabecula} A small beam, strut, or partition of tissue, especially within spongy bone. Trabeculae form a supportive internal network that helps bones bear forces while keeping them relatively light.
\medterm{Trabecular Bone} Light, honeycomb-like bone found inside many bones and near their ends and joints. It supports the body, contains bone marrow, and changes quickly when weight, hormones, medicines, or diseases such as osteoporosis affect bone strength.

\synonyms
Spongy bone; cancellous bone

\medterm{Trabecular Meshwork} A sieve-like tissue near the front of the eye through which the clear fluid called aqueous humor drains. If the drainage pathway becomes blocked, pressure inside the eye may rise and damage the optic nerve in glaucoma.

\medterm{Trabeculectomy} A conventional glaucoma filtration surgery that creates a new drainage pathway for
aqueous humor by removing a small piece of the trabecular meshwork and creating a
partial-thickness scleral flap under which aqueous humor filters from the anterior
chamber into the subconjunctival space, forming a filtering bleb.

\medterm{Trabeculoplasty} Laser treatment that improves aqueous-humor drainage through the trabecular meshwork.

\medterm{Trabulsiella} A genus of gram-negative bacteria related to Enterobacter that is rarely associated with opportunistic human infection.

\medterm{Trabulsiella Guamensis} A gram-negative bacterium rarely isolated from human clinical specimens; disease significance depends on a compatible infection rather than colonization.

\medterm{Trace Elements} Minerals needed in very small amounts for normal enzyme function, blood formation, thyroid activity, and other body processes.
\medterm{Trace Mineral} A mineral required by the body in very small amounts.

\medterm{Tracer} A detectable substance used to follow a biological, chemical, or physical process. Medical tracers can show blood flow, organ function, metabolism, or disease location through imaging or laboratory testing.
\medterm{Tracer Principle} Administered in amounts too small to produce pharmacological effect. Dose optimized for imaging, not therapeutic.
\medterm{Trachea} The windpipe connecting the larynx to the right and left main bronchi. Cartilage rings help keep it open so air can move to and from the lungs.
\medterm{Tracheal Agenesis With Bronchoesophageal Fistula} A very rare birth defect in which the windpipe is absent or severely underdeveloped and an abnormal passage connects the breathing system with the food pipe. Newborns have extreme breathing difficulty and need immediate specialist care.

\medterm{Tracheal Foreign Body} A foreign object lodged in the trachea, causing severe respiratory distress, stridor,
and complete airway obstruction. It is a medical emergency requiring immediate
bronchoscopic removal. In complete obstruction, emergency cricothyrotomy may be needed.

\medterm{Tracheal Rupture} Tracheal rupture is a tear in the windpipe caused by trauma, instrumentation, surgery, or pressure and can produce air leakage, breathing difficulty, and life-threatening complications.

\medterm{Tracheal Stenosis} Narrowing of the trachea, often a complication of prolonged endotracheal intubation or
tracheostomy. It presents with progressive dyspnea, stridor, and wheezing that may be
mistaken for asthma. Diagnosis is by bronchoscopy or CT. Treatment is balloon dilation,
stenting, or surgical resection.

\medterm{Tracheitis} Inflammation or infection of the windpipe. Mild cases may occur with a viral illness, but bacterial infection can rapidly block a child's airway and cause high fever, a harsh cough, noisy breathing, and severe illness. It requires emergency airway and antibiotic treatment.
\medterm{Trachelectomy} Trachelectomy is surgical removal of the cervix, sometimes with nearby tissue, used in selected cervical cancers or other conditions.

\medterm{Tracheobronchomalacia} Abnormal softening of the trachea and/or bronchi causing excessive airway collapse
during breathing. It may be congenital or acquired. Symptoms include chronic cough,
recurrent infections, and expiratory airflow obstruction. Diagnosis is by dynamic CT or
bronchoscopy. Treatment includes CPAP, stenting, or surgical reconstruction.

\medterm{Tracheobronchomegaly} Abnormal enlargement of the trachea and main bronchi, usually associated with
connective-tissue weakness. It may cause recurrent lower-respiratory infection,
bronchiectasis, ineffective cough, dyspnea, and difficulty clearing airway secretions.

\medterm{Tracheoesophageal Fistula} An abnormal connection between the trachea and esophagus, often associated with
esophageal atresia. Presents with recurrent aspiration, coughing with feedings, and
cyanosis during meals. The H-type fistula may present later in infancy with chronic
aspiration and failure to thrive. Diagnosis by esophagram or bronchoscopy.

\medterm{Tracheoesophageal Fistula And Esophageal Atresia Repair} Surgery in a newborn to close an abnormal connection between the windpipe and food pipe and join separated parts of the food pipe when possible. The operation is urgent; breathing problems, leakage, narrowing, reflux, and feeding difficulties may persist.

\medterm{Tracheomalacia} Softening of the tracheal cartilage causing excessive collapse of the trachea during
expiration. It may be congenital or acquired (prolonged intubation, relapsing
polychondritis). Symptoms include expiratory stridor and difficulty clearing secretions.
Diagnosis is by dynamic CT or bronchoscopy.

\medterm{Tracheomalacia - Acquired} Acquired tracheomalacia is weakening of the tracheal wall after injury, inflammation, compression, or prolonged ventilation, causing airway collapse and breathing symptoms.

\medterm{Tracheomalacia - Congenital} Congenital tracheomalacia is weakness of the infant tracheal cartilage causing dynamic airway collapse, noisy breathing, cough, or recurrent respiratory symptoms.

\medterm{Tracheostomy} A surgically created opening through the neck into the trachea, usually with insertion of a tube.
\medterm{Tracheostomy Care} Care of a breathing tube and neck opening, including cleaning the skin and tube, suctioning mucus when needed, checking airflow, and keeping emergency equipment nearby. Blockage, bleeding, infection, or accidental tube displacement requires prompt attention.

\medterm{Tracheostomy Tube} A tube inserted through a tracheostomy opening to maintain access to the airway and permit breathing,
suction, or mechanical ventilation.

\medterm{Tracheostomy Tube - Eating} Eating and drinking with a tracheostomy tube requires swallowing assessment and individualized precautions because aspiration risk depends on airway protection, secretions, and the underlying illness.

\medterm{Tracheostomy Tube - Speaking} Speech with a tracheostomy tube may be supported by cuff deflation, a speaking valve, or other methods when airflow and airway protection are adequate; specialist assessment is required.

\medterm{Tracheotomy} An operation that creates an opening through the neck into the trachea. The opening is called a tracheostomy.

\medterm{Trachoma} A chronic eye infection caused by Chlamydia trachomatis that can scar the conjunctiva and cornea.

\medterm{Trachyonychia} Rough, thin, ridged nails with many tiny pits, sometimes called twenty-nail dystrophy when most or all nails are affected. It may occur on its own or with psoriasis, alopecia areata, eczema, or lichen planus.
\medterm{Traction} A controlled pulling force used to align or immobilize a fracture, reduce a dislocation, or relieve pressure on tissues. It may be applied manually, with weights, or through a medical device.
\medterm{Tractional Retinal Detachment} Separation of the neurosensory retina caused by contraction of fibrovascular membranes,
commonly associated with proliferative diabetic retinopathy.

\medterm{Tractotomy} Surgical interruption of a nerve tract in the brain or spinal cord. Surgical interruption of a nerve tract, usually to reduce severe pain or other neurologic symptoms.
\medterm{Train-Of-Four Monitoring} A monitoring test used during anesthesia that gives four brief electrical pulses to a peripheral nerve and observes the muscle responses. The pattern and strength of the responses show how much a muscle-relaxing medicine is still working and whether recovery is adequate.

\medterm{Trained Immunity} Innate immune memory: epigenetic/metabolic reprogramming of monocytes/macrophages/NK
cells after initial stimulus (BCG, beta-glucan, Candida, oxidized LDL) -> enhanced
response to secondary challenge (heterologous protection).

\medterm{Trait} A trait is a measurable characteristic of an organism influenced by genes, environment, development, or their interaction.

\medterm{Trajectory} The path a projectile follows in flight and, in wound ballistics, the course it
traverses through tissue from entrance to exit or rest. Internal trajectory is
reconstructed from entrance and exit defects, the organ track, recovered fragments, and
clothing damage; angles disclose shooter position and victim posture.

\medterm{TRALI} Transfusion-Related Acute Lung Injury, a life-threatening complication of blood
transfusion characterized by acute hypoxemia and bilateral pulmonary edema within 6
hours of transfusion, without evidence of volume overload.

\medterm{Tram-Track Appearance} Histological finding on silver stain of MPGN representing GBM duplication. Results from
mesangial cell interposition into capillary wall with new basement membrane formation.
Creates double-contour pattern on Jones methenamine silver stain.

\medterm{Trance} An altered state of awareness with unusual attention, responsiveness, or sense of surroundings. It may occur during hypnosis, seizures, intoxication, sleep-related events, or psychological states and requires context for interpretation.
\medterm{Tranexamic Acid} An antifibrinolytic medicine that reduces bleeding by inhibiting plasminogen activation.
\medterm{Trans Fat} An unsaturated fat containing at least one trans-configured double bond. Industrial trans fats, formerly common in partially hydrogenated oils, can raise LDL cholesterol, lower HDL cholesterol, and increase cardiovascular risk; small amounts also occur naturally in some meat and dairy foods.

\medterm{Trans Fatty Acid} An unsaturated fat with a trans shape, often formed during industrial processing and associated with higher heart disease risk.

\synonyms
Trans fat

\medterm{Transabdominal Cerclage} Placement of a strong band around the upper cervix through abdominal surgery, usually before pregnancy or early in pregnancy, to reduce the risk of premature opening in severe cervical weakness. It requires cesarean birth and carries surgical, bleeding, infection, and pregnancy risks.

\synonyms
TAC; Cervico-Isthmic Cerclage

\medterm{Transabdominal Scan} An ultrasound scan performed with the transducer placed on the maternal abdomen. A full
bladder improves visualization of pelvic structures. Used in early pregnancy to confirm
viability and gestational age, and throughout pregnancy to assess fetal growth, anatomy,
amniotic fluid volume, and placental location.

\medterm{Transaminase, Serum Glutamic Pyruvic} SGPT, now called ALT, is a liver-associated enzyme measured in blood and often elevated when liver cells are injured.

\medterm{Transaminitis} Elevation of blood levels of aminotransferase enzymes, usually alanine aminotransferase and aspartate aminotransferase, indicating possible liver-cell injury but not identifying its cause.

\medterm{Transarterial Chemoembolization} A procedure in which chemotherapy and embolizing particles are infused selectively into
tumor-feeding hepatic arteries, combining regional chemotherapy with ischemic tumor
necrosis; used for unresectable hepatocellular carcinoma.

\medterm{Transbronchial Biopsy} Bronchoscopic sampling of lung parenchyma through the bronchial wall. It may evaluate interstitial or granulomatous disease, infection, transplant rejection, or peripheral lesions; complications include bleeding and pneumothorax.

\medterm{Transcatheter Aortic Valve Replacement} Replacement of a narrowed aortic heart valve using a folded artificial valve delivered through a catheter, usually from an artery in the groin. It avoids open-heart surgery for selected people, but bleeding, stroke, blood-vessel injury, rhythm problems, leakage, or need for a pacemaker can occur.

\synonyms
TAVR; TAVI; Transcatheter Aortic Valve Implantation

\medterm{Transcatheter Edge-To-Edge Repair} A catheter procedure that joins parts of the mitral-valve leaflets to reduce leakage. It may help selected people with severe symptoms who are not good candidates for open surgery; bleeding, infection, stroke, valve narrowing, residual leakage, or device problems are possible.

\synonyms
TEER; MitraClip Procedure; Percutaneous Mitral Valve Repair

\medterm{Transcervical Foley Catheter For Cervical Ripening} A mechanical method for cervical ripening: a Foley catheter (16-30 Fr) is inserted
through the cervix, the balloon is inflated with 30-60 mL of saline, and traction is
applied (or the catheter is taped to the thigh with gentle tension). The balloon applies
direct pressure to the internal os, stimulating endogenous prostaglandin release.

\medterm{Transcervical Resection Of Endometrium} A camera-guided procedure through the cervix that removes or destroys the uterine lining to reduce heavy menstrual bleeding. It is not suitable for future pregnancy and requires careful evaluation beforehand.
\medterm{Transcranial Doppler} Ultrasound measuring blood flow velocity in cerebral arteries. Used for vasospasm detection, stroke risk assessment.
\medterm{Transcranial Doppler Scanning} An ultrasound test measuring blood-flow speed and direction in major arteries inside the skull. It can help assess narrowing, blockage, vasospasm, emboli, or circulation after certain brain injuries.
\synonyms
Transcranial Doppler; TCD

\medterm{Transcranial Doppler Ultrasound} An ultrasound test that measures blood-flow speed and direction in major arteries inside the skull. It can help assess stroke risk, vessel narrowing, bleeding-related spasm, and blood flow changes.

\medterm{Transcription} Synthesis of RNA from DNA template by RNA polymerase. Eukaryotic: RNA Pol II transcribes
mRNA, rRNA, snRNA; RNA Pol I transcribes rRNA; RNA Pol III transcribes tRNA, 5S rRNA.
Steps: initiation (promoter recognition, transcription factors), elongation (5' to 3'),
termination. Promoter elements: TATA box, CAAT box, GC box.

\medterm{Transcription Factor} Protein binding to specific DNA sequences, regulating transcription. Activators enhance
transcription; repressors inhibit. Examples: p53, NF-kB, steroid hormone receptors.

\medterm{Transcutaneous Electrical Nerve Stimulation} TENS: low-voltage electrical pulses delivered through electrodes on the skin to reduce pain for some people. It does not treat every cause of pain and should be used according to clinical guidance.

\synonyms
TENS

\medterm{Transdermal} A route of drug administration delivering drug through intact skin into the systemic
circulation via a patch or gel, bypassing first-pass metabolism and providing steady
plasma concentrations over hours to days. Examples: nitroglycerin, fentanyl, nicotine,
estrogen, clonidine, rotigotine.

\medterm{Transfer Training} Teaching safe transfer techniques: bed to chair, chair to toilet, floor to standing. Use
of transfer boards, lift equipment.

\medterm{Transferase} An enzyme that moves a chemical group from one molecule to another. Transferases help cells make and break down substances, and some transferase levels in blood help identify tissue injury or disease.

\medterm{Transference} A psychoanalytic concept in which feelings, attitudes, and expectations originally
directed toward a significant figure from childhood are unconsciously redirected onto
the therapist. In psychodynamic therapy, transference is a therapeutic tool for
understanding the patient's relational patterns and resolving conflicts.

\medterm{Transferrin} An iron transport protein synthesized by the liver, carrying iron in plasma (2
iron-binding sites). In iron deficiency: transferrin is elevated with low ferritin,
giving high TIBC and low transferrin saturation; in iron overload/inflammation/chronic
disease: transferrin is low (acute phase).

\medterm{Transferrin Saturation} The percentage of transferrin's iron-binding sites occupied by iron. It is calculated from iron and total iron-binding capacity and helps evaluate iron deficiency and iron overload.

\synonyms
TSAT

\medterm{Transfusion} Delivery of donated blood or blood components into a person’s bloodstream. It can replace lost red cells, platelets, or clotting factors but carries risks including reactions, infection, fluid overload, and lung injury.


\medterm{Transfusion Of Neonate} Administration of a blood product to a newborn for clinically significant anemia, bleeding, or
another specific indication. The decision, product characteristics, volume, and rate depend
on the infant’s symptoms, laboratory results, gestational age, and local transfusion
protocol.

\medterm{Transfusion Reaction} An adverse reaction to transfusion of blood products. Acute hemolytic transfusion
reaction: ABO incompatibility, immediate fever/chills, hypotension, back pain, DIC,
renal failure; management: STOP transfusion, maintain BP/urine output, supportive care.

\medterm{Transient Familial Hyperbilirubinemia} Transient familial hyperbilirubinemia is inherited temporary elevation of bilirubin, often from reduced conjugation, causing episodic jaundice without major liver injury.

\medterm{Transient Global Amnesia} A temporary episode of sudden inability to form new memories, with repetitive questioning, that usually resolves within 24 hours.

\synonyms
Amnesia, Transient Global

\medterm{Transient Hypogammaglobulinemia Of Infancy} A temporary delay in normal immunoglobulin production during infancy, causing low antibody levels and increased susceptibility to infections before immune function improves spontaneously.

\synonyms
TGA

\medterm{Transient Ischaemic Attack} A brief episode of neurological dysfunction caused by temporary focal brain, spinal-cord, or retinal ischaemia without persistent tissue infarction.

\synonyms
TIA

\medterm{Transient Ischaemic Attacks Or Episodes} Short episodes of weakness, numbness, speech trouble, vision loss, or other neurological symptoms caused by temporarily reduced blood flow to the brain. Symptoms resolve without lasting brain injury, but a TIA is an emergency warning of possible stroke.
\medterm{Transient Ischemic Attack} A transient ischemic attack is a temporary episode of neurological dysfunction
caused by focal brain, spinal cord, or retinal ischemia without acute infarction.

\medterm{Transient Ischemic Attack (TIA, Mini-Stroke)} Alternate index label for Transient Ischemic Attack (TIA, Mini-Stroke).
\medterm{Transient Pain} Pain that is brief and resolves rather than continuing or recurring persistently. Even brief pain may need assessment if it is severe, recurrent, or associated with warning symptoms.

\medterm{Transient Synovitis} Temporary inflammation of the hip lining, usually in children ages 3 to 8 and often after a viral illness. It causes hip, groin, thigh, or knee pain and a limp, but usually little fever or illness. Rest and pain medicine generally help. Doctors must distinguish it from a dangerous joint infection, especially when fever, severe pain, or refusal to walk is present.

\synonyms
Toxic Synovitis; Irritable Hip; Transient Synovitis of the Hip

\medterm{Transient Tachypnea - Newborn} Temporary rapid breathing in a newborn caused by delayed clearance of fetal lung fluid, usually beginning soon after birth and improving within one to three days with supportive care.

\medterm{Transient Tachypnea Of The Newborn} Temporary rapid breathing after birth, usually caused by delayed clearance of fetal lung fluid. It is more common after caesarean birth and typically improves within one to three days with supportive care.

\medterm{Transillumination} Transillumination shines light through a body part or structure to reveal fluid, cavities, or other differences in tissue density.

\medterm{Transition To Adulthood} The process of moving from child-focused to adult healthcare, education, work, and community services. Planning helps young people develop self-management skills and maintain continuity of care.

\medterm{Transitional Care} Coordinated care provided when a person moves between healthcare settings or levels of care, such
as from hospital to home.

\medterm{Transitional Epithelium} Specialized stretchable epithelium lining the renal pelvis, ureters, bladder, and proximal urethra, allowing urinary organs to expand while maintaining a protective barrier.

\synonyms
Transjugular Intrahepatic Portosystemic Shunt

\medterm{Transjugular Intrahepatic Portosystemic Shunt} A radiologically created channel between a portal vein branch and a hepatic vein that lowers portal pressure, used for complications such as variceal bleeding or refractory ascites; encephalopathy and shunt dysfunction can occur.

\medterm{Translational Medicine} An interdisciplinary field that turns findings from laboratory, clinical, and community research into better methods of disease prevention, diagnosis, treatment, and public health. It connects research findings with patient care and community health.

\synonyms
Translational Medicines

\medterm{Translation} The cellular process of making a protein from instructions carried by messenger RNA. A ribosome reads the message in order, while transfer molecules bring the matching amino acids. The process starts, builds the chain, and stops at a signal; the finished chain folds into a working protein.

\medterm{Translation Initiation} The first stage of protein synthesis, when a ribosome recognizes messenger RNA, selects the start codon, and assembles the initiator transfer RNA and ribosomal subunits.

\medterm{Translocation} Movement of a chromosome segment to a new location or movement of a molecule, cell, or organism between compartments; chromosomal translocations can create fusion genes or alter gene dosage.
\medterm{Transmembrane Transport} Movement of ions, nutrients, drugs, or other molecules across a cell membrane through channels, carriers, pumps, or vesicles, sometimes requiring energy.

\medterm{Transmission} Passage of an infection, signal, genetic trait, force, or substance from one source to another; in infectious disease, route, dose, and host susceptibility determine risk.

\medterm{Transmural Infarction} A heart attack damaging the full thickness of part of the heart wall, usually from prolonged blocked blood supply.

\medterm{Transplacental} Passing from the pregnant person's blood across the placenta to the fetus, or in the opposite direction. Antibodies, medicines, chemicals, and some infections can cross the placenta, with effects that depend on the substance and timing.

\medterm{Transplant} Transfer of cells, tissue, or an organ from one site or person to another to replace damaged tissue or restore function. Donor matching, rejection prevention, infection risk, and long-term follow-up depend on the type of transplant.

\medterm{Transplant Immunology} The study of how the immune system reacts to donated organs, tissues, or cells. Blood-group and tissue matching, antibody testing, and immunosuppressive medicines help reduce rejection, but they also increase infection risk and require regular monitoring of the transplant and medicine effects.

\medterm{Transplant Rejection} Immune-mediated destruction of a transplanted organ or tissue, classified by timing and
mechanism. Hyperacute: preformed donor-specific antibodies, minutes to hours
post-transplant, immediate graft thrombosis, irreversible; prevention by cross-matching.

\medterm{Transplant Services} The coordinated medical care involved in organ or tissue transplantation, including donor assessment, recipient evaluation, surgery, rejection prevention, infection monitoring, rehabilitation, and long-term follow-up.

\medterm{Transplant Surgery} Transplant surgery involves replacing a failing organ with a healthy donor organ. Kidney
transplant: living donors (laparoscopic nephrectomy, left kidney preferred) or deceased
donors (standard criteria, expanded criteria, DCD).

\medterm{Transplantation} Transfer of an organ, tissue, or cells from a donor to a recipient to replace or restore failing function.

\synonyms
Transplant

\medterm{Transplantation Conditioning} Treatment given before some stem-cell transplants to suppress the recipient's bone marrow and immune system and make room for donated cells. It may use chemotherapy, radiation, or both and can cause infection, bleeding, infertility, or organ injury.

\medterm{Transplantation Genetics} The study and clinical use of inherited and tissue genetic information in transplantation. It can help assess donor compatibility, inherited disease, rejection risk, and the chance that a donor organ carries a genetic condition.

\medterm{Transplantation Immunology} The study of how the immune system recognizes and responds to transplanted cells, tissues, or organs. It explains rejection and tolerance and guides donor matching, antibody testing, and medicines that suppress harmful immune responses.

\medterm{Transport} Regulations for shipping radioactive materials. Package types: excepted, A, B.

\medterm{Transport Protein} A membrane or soluble protein that carries a molecule or ion across a membrane or through body fluids, often determining absorption, distribution, and cellular concentration.

\medterm{Transport Vesicle} A transport vesicle is a membrane-bound carrier that moves proteins, lipids, or other cargo between cellular compartments.

\medterm{Transposition Of Great Vessels} Alternate index label; see \textit{Transposition of the great arteries}.

\medterm{Transposition Of The Great Arteries} A congenital heart defect in which the aorta arises from the right ventricle and the pulmonary
artery from the left ventricle. This creates parallel circulations, so survival after birth
depends on adequate mixing through an associated opening or ductus until definitive
treatment is provided.

\synonyms
Transposition Of Great Vessels

\medterm{Transposons} Mobile genetic elements that can move between different locations within the bacterial
genome or between the chromosome and plasmids. Transposons carry genes including
antibiotic resistance genes and can facilitate the spread of resistance.

\medterm{Transrectal Ultrasonography} Ultrasound imaging of the prostate and nearby structures using a probe placed in the rectum. It can guide biopsies and assess size or abnormalities, but may cause discomfort, bleeding, or infection.

\synonyms
TRUS

\medterm{Transtracheal} Transtracheal means passing through or relating to the trachea, such as a catheter, oxygen delivery route, or diagnostic procedure.

\medterm{Transthyretin Amyloidosis} A disease in which abnormal transthyretin protein builds up in organs, especially the heart or nerves. It may be inherited or develop with age without an inherited gene change.
\synonyms
ATTR amyloidosis; ATTR-CM; Transthyretin-related amyloidosis

\medterm{Transtympanic} Transtympanic means passing through the tympanic membrane, such as medicine delivered into the middle ear or a procedure crossing the eardrum.

\medterm{Transudate} A transudate is a type of edema fluid with low protein content and few cells, resulting
from imbalances in hydrostatic and oncotic pressures across capillary walls rather than
from inflammation or increased vascular permeability.

\medterm{Transudation} Movement of low-protein fluid through a membrane when pressure or protein balance shifts, often causing certain edema or effusions.
\medterm{Transudative Pleural Effusion} From systemic factors altering hydrostatic/oncotic pressures. Most commonly from heart
failure. Low protein, low LDH. Treatment addresses the underlying cause.

\medterm{Transurethral Incision Of The Prostate} An operation making small cuts in the prostate through an instrument passed in the urethra. It can improve urine flow in selected people with prostate enlargement while preserving most prostate tissue.

\synonyms
TUIP

\medterm{Transurethral Microwave Thermotherapy} A treatment for some people with an enlarged prostate in which a catheter uses microwave heat to shrink or destroy prostate tissue blocking urine flow. It is less commonly used now and can cause pain, bleeding, infection, or temporary urinary problems.

\synonyms
TUMT

\medterm{Transurethral Needle Ablation} A treatment using needles passed through the urethra to deliver radiofrequency heat into enlarged prostate tissue. The heat destroys selected tissue and may improve urine flow, but relief develops gradually and complications can occur.

\synonyms
TUNA

\medterm{Transurethral} Performed through the urethra, usually using an instrument to examine, treat, or remove tissue from the urinary tract.
\medterm{Transurethral Resection} Removal of tissue through an instrument passed along the urethra, commonly for prostate or bladder disease.
\medterm{Transurethral Resection Of The Prostate} An endoscopic operation that removes obstructing prostate tissue to
improve urine flow in selected people with benign prostatic enlargement. Alternatives and
indications depend on symptoms, prostate size, complications, and patient preference.


\medterm{Transvaginal} Performed through the vagina, commonly describing an ultrasound examination or a procedure involving the uterus, cervix, ovaries, or other pelvic structures.
\medterm{Transvaginal Ultrasound} Transvaginal ultrasound uses a covered probe in the vagina to produce images of the uterus, ovaries, cervix, pregnancy, or nearby pelvic structures.

\medterm{Transverse} Running across or perpendicular to the long axis or midline.
\medterm{Transverse Colon} Longest, most mobile colonic segment from hepatic to splenic flexure, approximately 45
cm, with complete mesentery (transverse mesocolon). Crosses anterior to duodenum,
pancreas, and major vessels. Splenic flexure is the highest colon point and common gas
trapping site. Most common diverticular disease site along with the sigmoid.

\medterm{Transverse Lie} A fetal lie where the long axis of the fetus is perpendicular to the long axis of the
mother. The presenting part is a shoulder or arm. Occurs in approximately 1\% of term
pregnancies. Risk of cord prolapse and obstructed labor. Cesarean delivery required as
vaginal delivery is not possible.

\medterm{Transverse Myelitis} An inflammatory disorder of the spinal cord causing acute or subacute motor, sensory,
and autonomic dysfunction.

\medterm{Transverse Plane} A horizontal plane that divides the body or an organ into upper (superior) and lower
(inferior) portions.

\medterm{Transverse Vaginal Septum} A congenital band of tissue crossing the vagina and partly or completely blocking it. A complete blockage can cause absent periods and monthly pelvic pain after puberty; imaging confirms the condition, and surgery may create an open vaginal passage.

\medterm{Transversus Abdominis Muscle} Deepest flat abdominal muscle with horizontal fibers, from inguinal ligament, iliac
crest, thoracolumbar fascia, and inner lower six costal cartilages. Innervated by T7-T12
and L1. Compresses abdomen without trunk movement, providing core stability and
increased intra-abdominal pressure during coughing, lifting, Valsalva.

\medterm{Transvestic Disorder} A paraphilic disorder involving recurrent sexual arousal from cross-dressing that causes marked distress,
impairment, or harm, or is acted on without consent. It is distinct from gender dysphoria,
which concerns gender identity rather than sexual arousal.

\medterm{Tranylcypromine} An antidepressant that increases certain brain chemicals by blocking their breakdown. It can interact dangerously with many medicines and foods high in tyramine, so careful prescribing and a medicine-free waiting period before or after some treatments are required.
\medterm{TRAP Sequence} Twin Reversed Arterial Perfusion sequence, a complication of monochorionic twin
gestations where one twin (the pump twin) perfuses the other (the acardiac twin) via
arterioarterial anastomoses.

\medterm{Trapezium (Multangulum Majus)} A small wrist bone at the base of the thumb. It forms the thumb's highly mobile joint with the hand and can be affected by fractures or arthritis that causes pain during gripping and pinching.
\medterm{Trapezius Muscle} Large triangular muscle from occipital bone and C7-T12 spinous processes to clavicle,
acromion, and scapular spine. Innervated by CN XI and C3-C4. Upper fibers elevate
scapula, middle fibers retract, lower fibers depress. Essential for shoulder stability
and scapulothoracic motion.

\medterm{Trapezoid} Smallest distal row carpal bone between trapezium and capitate, articulating with the
second metacarpal. Wedge shape provides stability for the firmly fixed second
metacarpal. Fractures rare (less than 1 percent of carpal fractures). Contributes to the
radial column wrist stability.

\medterm{Trapezoid Bone} A small distal-row carpal bone between the trapezium and capitate that articulates mainly with the index metacarpal and helps stabilize the radial side of the hand.

\medterm{Trastuzumab} Trastuzumab is a monoclonal antibody targeting HER2 used for HER2-positive breast, gastric, and other cancers; infusion reactions and potentially reversible cardiomyopathy require cardiac monitoring.

\medterm{Trauma} Physical injury caused by an outside force, or the mental response to a deeply distressing event. Physical trauma may damage organs, bones, or blood vessels; psychological trauma can affect emotions, sleep, and daily life.
\medterm{Trauma (Vascular)} Injury to a blood vessel caused by blunt force, penetrating injury, fracture, medical procedures, or
other mechanisms.

\synonyms
Vascular

\medterm{Trauma Center} Hospital designated to provide comprehensive care for injured patients. Levels: Level I
(complete care, research, prevention), Level II (similar but limited resources), Level
III (stabilization, transfer), Level IV (advanced trauma life support, transfer).
Transfer criteria: need for higher level of care, specific injuries.

\medterm{Trauma Surgery} Trauma surgery encompasses the acute care of injured patients, guided by ATLS
principles: Airway maintenance with cervical spine protection, Breathing and
ventilation, Circulation with hemorrhage control, Disability/neurologic assessment, and
Exposure/environmental control.

\medterm{Traumatic Amputation} Traumatic amputation is loss of a body part from injury and requires hemorrhage control, preservation of the amputated part when possible, urgent surgery, rehabilitation, and psychological support.

\medterm{Traumatic Bleb} A filtering bleb formed after glaucoma filtration surgery (trabeculectomy) that develops
complications from trauma, or a thin-walled, elevated bleb that leaks or is at risk of
rupture. Traumatic bleb failure or overfiltering can lead to hypotony, choroidal
detachment, and endophthalmitis risk.

\medterm{Traumatic Brain Injury} Disruption of brain function caused by an external mechanical force. Severity: mild (GCS
13-15, concussion), moderate (GCS 9-12), severe (GCS 3-8). Primary injury: direct damage
at moment of impact (contusion, diffuse axonal injury, intracranial hemorrhage).

\synonyms
TBI

\medterm{Traumatic Events And Children} Traumatic events can affect children’s safety, behavior, sleep, learning, and emotions; supportive caregivers, stable routines, and professional assessment help identify persistent trauma responses.

\medterm{Traumatic Grief} Alternate index label; see \textit{Complicated grief}.

\medterm{Traumatic Injury} Damage to the body caused by an external force, such as a fall, collision, blow, cut, burn, or penetrating object. Injury may affect skin, bones, organs, blood vessels, nerves, or the brain.
\medterm{Traumatic Injury Of The Bladder And Urethra} Damage to the bladder or urethra caused by an accident, surgery, or penetrating injury. Blood in the urine, inability to urinate, pelvic pain, or shock requires urgent evaluation.

\medterm{Traumatic Shock} A life-threatening state in which severe injury causes inadequate blood flow and oxygen delivery to tissues.

\medterm{Traumatic Spondylolisthesis Of The Axis (Hangman’S Fracture)} A bilateral fracture of the pars interarticularis of the C2 vertebra (axis), typically
caused by hyperextension of the head. Classically associated with judicial hanging but
also occurs in motor vehicle accidents and falls. May present with neck pain,
torticollis, or neurological deficits if the spinal cord is compressed.

\synonyms
hangman’s fracture

\medterm{Traumatology} The study and treatment of physical injury and its consequences.
\medterm{Traumeel S} Traumeel S is a multi-ingredient complementary product marketed for pain and inflammation; evidence, formulation, and safety depend on the product and it should not replace needed care.

\medterm{Travel Sickness (Motion Sickness / Kinetosis)} Nausea, dizziness, sweating, or vomiting caused by conflicting signals from the eyes and balance organs during travel or other motion. Looking toward the horizon, fresh air, gradual exposure, and suitable preventive medicines may help; severe or persistent symptoms need assessment.


\medterm{Travelers' Diarrhea} Acute diarrhea occurring during or shortly after travel, usually caused by contaminated food or water and commonly involving bacterial, viral, or parasitic infection.

\medterm{Travoprost} Travoprost is a prostaglandin analogue eye drop that increases uveoscleral outflow and lowers intraocular pressure in glaucoma or ocular hypertension; iris pigmentation and eyelash changes may occur.

\medterm{Trazodone} Trazodone is a serotonin antagonist and reuptake inhibitor used for depression and sometimes insomnia; sedation, orthostatic hypotension, serotonin syndrome, and rare priapism are important adverse effects.

\medterm{Trazodone Overdose} Poisoning from taking too much trazodone. It can cause severe sleepiness, vomiting, low blood pressure, seizures, abnormal heart rhythms, or a prolonged painful erection.

\medterm{Trb Gene} A gene involved in biological processes that vary by species; its precise identity and function must be interpreted from the organism and gene nomenclature used.

\medterm{Trd Gene} The gene encoding the delta chain of the T-cell receptor in humans, contributing to gamma-delta T-cell antigen recognition.

\medterm{Treacher Collins Syndrome} An inherited condition affecting development of the face and skull. It may cause small cheekbones and jaw, downward-slanting eyes, malformed ears, and hearing loss; breathing or feeding may also be difficult. Intelligence is usually normal. Care may include airway support, hearing treatment, and reconstructive surgery.

\medterm{Treacle} A thick sugar syrup with little current medical use. Older medical writing sometimes used the word for a sweetened mixture or supposed antidote; it is not a modern treatment or diagnosis.

\medterm{Treat-To-Target Strategy} A management approach in which treatment is adjusted at regular intervals toward a
predefined goal such as remission or low disease activity. It requires validated
disease-activity measures, shared decisions, monitoring, and timely escalation or
modification.

supportive.

\medterm{Treatment Epoch} A treatment epoch is a defined period in a clinical study or treatment plan during which a particular therapy or treatment phase is given.


\medterm{Treatment-Resistant Depression} Depression that does not improve adequately after appropriate trials of standard
treatments. It requires reassessment of diagnosis, adherence, comorbidity, and treatment
options.

\medterm{Trematoda} A class of parasitic flatworms called flukes; human infections include schistosomiasis and liver or lung fluke disease.

\medterm{Trematosphaeria Grisea} An environmental dematiaceous fungus rarely associated with human infection; clinical importance depends on recovery from a compatible lesion or sterile site.

\medterm{Tremelimumab} A monoclonal antibody that blocks CTLA-4 and enhances T-cell activity; it is used with durvalumab for selected unresectable hepatocellular carcinoma and other cancers, with immune-related adverse effects possible.

\medterm{Tremor} An involuntary, rhythmic shaking of a body part. It may occur at rest, while holding a position, or during movement and can result from medicines, caffeine, anxiety, essential tremor, Parkinson disease, or other neurologic or metabolic conditions.

\medterm{Trench Fever} A louse-borne infection caused by Bartonella quintana, causing recurring fever, headache, weakness, and bone pain.
\medterm{Trench Foot} An injury caused by keeping feet cold and wet for a long time. It can cause numbness, swelling, pain, skin color changes, blisters, infection, and tissue death without gradual warming and medical care.
\medterm{Trench Mouth} Trench mouth is acute necrotizing ulcerative gingivitis causing painful bleeding gums, foul breath, and ulceration, often associated with plaque, smoking, stress, or poor nutrition.

\medterm{Trendelenberg Position} A position in which the patient lies supine with the feet higher than the head; clinical uses and risks depend on context.
\medterm{Trendelenburg} Trendelenburg position places the body supine with the feet higher than the head; its use is procedure-specific and may affect breathing, pressure, and circulation.

\medterm{Trendelenburg Position} A position lying on the back with the head and upper body lower than the feet. It is used for selected procedures but can worsen breathing, pressure in the head, or blood pressure in some patients.
\medterm{Treosulfan} An alkylating chemotherapy drug used in conditioning before allogeneic hematopoietic stem-cell transplantation; infection, mucositis, infertility, and organ toxicities may occur.

\medterm{Trepanning} Drilling or cutting a round opening in a bone, especially the skull. It was used historically for many reasons and is now done only for selected medical procedures, such as relieving pressure, reaching the brain, or treating a bone problem.
\medterm{Trephine} A circular cutting instrument used to remove a disk of tissue or bone, such as during a bone-marrow biopsy or corneal procedure.

\medterm{Trephining} Another term for trepanning, creating a circular opening in bone.
\medterm{Trephining (Trepanation)} A surgical procedure in which a hole is drilled or scraped into the skull (cranium),
exposing the dura mater. Historically, trephining dates back to Neolithic times (10,000
BCE) and was performed for ritual, therapeutic (releasing 'evil spirits'), and
post-traumatic indications.

\synonyms
Trepanation

\medterm{Treponema} A group of spiral-shaped bacteria, including Treponema pallidum, which causes syphilis. Different species can live in people or the environment, and identifying the species helps determine whether disease and treatment are present.
\medterm{Treponema Pallidum} A spirochete that causes syphilis, a sexually transmitted infection. The organism is
transmitted through sexual contact or transplacentally (congenital syphilis). Primary
syphilis presents with a painless genital ulcer (chancre). Secondary syphilis presents
with rash, mucous patches, and condylomata lata.

\medterm{Treponema Pallidum Infection} Infection with the bacterium Treponema pallidum, which causes syphilis and can affect the
skin, nervous system, cardiovascular system, and other organs.

\medterm{Treponemal Infections} Infections caused by spiral-shaped Treponema bacteria. They include syphilis and related diseases and may affect the skin, nerves, heart, eyes, or other organs; diagnosis usually combines clinical findings with blood or lesion tests.
\medterm{Trepopnea} Dyspnea that occurs in one lateral decubitus position but not the other. It is typically
worse when lying on the side of the affected lung. It may be caused by unilateral lung
disease, pleural effusion, or mediastinal mass.

\medterm{Tretinoin} Tretinoin is all-trans retinoic acid used topically for acne and photoaging and orally with arsenic or chemotherapy for acute promyelocytic leukaemia; skin irritation and systemic differentiation syndrome are important risks.

\medterm{TRG Gene} A gene encoding the gamma chain of the T-cell receptor, contributing to gamma-delta T-cell development and antigen recognition.

\medterm{Triiodothyronine} The thyroid hormone T3, which contains three iodine atoms and has stronger biological activity than T4. It helps regulate metabolism, heart rate, temperature, and growth.
\medterm{Triage} Sorting patients by urgency and resource needs so those at greatest risk receive timely care.
\medterm{Trial Indication} The disease, condition, symptom, or clinical purpose for which an investigational intervention is being studied in a clinical trial.

\medterm{Trial Of Labor After Cesarean} An attempt to give birth vaginally after a previous cesarean delivery. Whether it is appropriate
depends on the prior uterine incision, obstetric history, current pregnancy, facility
resources, and the person’s preferences. Counseling includes the possibility of successful
VBAC, repeat cesarean delivery, and uterine rupture.

\medterm{Trial Screening} The process of checking potential participants against a clinical trial’s eligibility criteria, including medical history, examination, laboratory tests, and informed consent.

\medterm{Triamcinolone} A corticosteroid medicine used on skin, by injection, or in other forms to reduce inflammation and itching.
\medterm{Triamcinolone Acetonide} A synthetic corticosteroid used topically, by injection, or in other preparations to reduce inflammation; prolonged or extensive use can cause local and systemic steroid effects.

\medterm{Triamterene} Triamterene is a potassium-sparing diuretic that blocks epithelial sodium channels in the collecting duct and is used for hypertension or oedema; hyperkalaemia and kidney stones are important risks.

\medterm{Triazenes} A group of chemical compounds that includes some anticancer agents; their effects and toxicity depend on the specific compound.

\medterm{Triazine} A chemical ring structure found in some medicines, herbicides, and industrial compounds; health effects depend on the specific substance and exposure.

\medterm{Triaziquone} An older alkylating anticancer drug that cross-links DNA; its clinical use is limited by toxicity and availability.

\medterm{Triazolam} Triazolam is a short-acting benzodiazepine hypnotic for short-term insomnia; amnesia, falls, dependence, respiratory depression, and dangerous interactions with opioids or other sedatives can occur.

\medterm{Triazole} A five-membered nitrogen-containing chemical ring found in several antifungal medicines and other drugs; clinical effects depend on the specific compound.

\medterm{Triazophos} An organophosphate insecticide that blocks acetylcholinesterase, allowing nerve signals to remain active. Significant exposure can cause sweating, vomiting, diarrhea, small pupils, breathing difficulty, weakness, seizures, or life-threatening poisoning.

\medterm{Tributyrin} A triglyceride that releases the short-chain fatty acid butyrate during digestion; it is studied experimentally for intestinal and metabolic effects.

\medterm{Triceps} A muscle with three heads, especially the triceps brachii at the back of the upper arm. It straightens the elbow and helps stabilize the shoulder; injury or nerve problems can weaken pushing and lifting.
\medterm{Triceps Brachii Muscle} Three-headed posterior arm muscle. Long head from infraglenoid tubercle, lateral head from posterior humerus above radial groove, medial head below. All insert on the olecranon. Innervated by radial nerve (C7-C8). Primary elbow extensor; assists shoulder extension/adduction.

\medterm{Trichiasis} Inward misdirection of eyelashes so they rub against the eye.
\medterm{Trichilemmal Carcinoma} A rare malignant follicular tumor showing outer-root-sheath differentiation, usually occurring on sun-damaged skin of older adults as a papule or nodule; complete excision is generally required.

\medterm{Trichilemmoma} A benign follicular tumor with outer-root-sheath differentiation, usually presenting as a small facial papule; multiple lesions may be associated with Cowden syndrome.

\medterm{Trichiniasis} An infection caused by Trichinella larvae, acquired from undercooked meat.
\medterm{Trichinosis} Infection with \textit{Trichinella} worms, usually acquired by eating undercooked infected pork or wild game. It may cause diarrhea, fever, muscle pain, and swelling around the eyes.
\medterm{Trichobezoar} A mass of swallowed hair in the digestive tract, usually the stomach, which may cause pain, vomiting, or blockage.

\medterm{Trichoblastoma} A rare benign follicular tumor of the scalp or face, usually presenting as a solitary dermal or subcutaneous nodule; it can resemble basal-cell carcinoma and is diagnosed histologically.

\medterm{Trichoepithelioma} A rare benign follicular adnexal tumor that usually presents as a skin-coloured papule or nodule on the face or scalp; it can resemble basal-cell carcinoma and may require histopathologic confirmation.
\medterm{Trichomonas} A flagellated protozoan, Trichomonas vaginalis, that infects the urogenital tract and causes trichomoniasis.
It is transmitted mainly through sexual contact and may cause discharge, irritation, or no
symptoms.

\medterm{Trichomonas Vaginalis} A flagellated protozoan that causes trichomoniasis, a sexually transmitted infection.
\medterm{Trichomoniasis} A sexually transmitted infection caused by the parasite \textit{Trichomonas vaginalis}. It may cause genital irritation and discharge or no symptoms; testing and treatment of partners help prevent reinfection.
\medterm{Trichophyton} A group of fungi that grow in keratin and cause ringworm infections of skin, hair, or nails. They spread through contact with infected people, animals, or objects and are treated with antifungal medicines.
\medterm{Trichorrhexis Nodosa} A hair-shaft defect in which weakened points split into brush-like ends, caused by inherited disease, trauma, chemicals, heat, or nutritional or metabolic disorders.

\medterm{Trichothiodystrophy Syndrome} A group of inherited disorders characterized by brittle sulfur-deficient hair with a
distinctive banded appearance, often accompanied by photosensitivity, ichthyosis,
developmental delay, short stature, or recurrent infection.

\medterm{Trichotillomania} Recurrent compulsive pulling out of one's hair, causing hair loss and distress or impairment.
\medterm{Trichuriasis} Intestinal infection caused by the whipworm \textit{Trichuris trichiura}, usually acquired by swallowing eggs from contaminated soil or food. Heavy infection can cause abdominal pain, diarrhea, anemia, or rectal prolapse.
\medterm{Tricuspid Annular Plane Systolic Excursion} A measurement made during an echocardiogram that shows how far part of the tricuspid valve moves as the right ventricle contracts. A low value may indicate reduced right-heart pumping and is interpreted with the complete scan and clinical findings.

\synonyms
TAPSE; Tricuspid Annular Excursion; RV Longitudinal Systolic Function

\medterm{Tricuspid Atresia} Absence of the tricuspid valve, resulting in no direct communication between the right
atrium and right ventricle. Blood flow depends on an atrial septal defect and
ventricular septal defect for mixing. Presents with cyanosis.

\medterm{Tricuspid Incompetence} Incomplete closure of the tricuspid heart valve, allowing blood to leak backward into the right upper chamber. Severe leakage can cause swelling, liver congestion, irregular heartbeat, fatigue, and breathlessness.
\medterm{Tricuspid Regurgitation} Incomplete closure of the tricuspid valve allowing blood to flow backward from the right
ventricle to the right atrium during systole. Causes include right ventricular dilation
(from pulmonary hypertension or left heart failure), endocarditis, rheumatic disease,
and carcinoid syndrome.

\medterm{Tricuspid Stenosis} Narrowing of the tricuspid valve opening, making it harder for blood to move from the right upper chamber to the right lower chamber. It can cause fatigue, swelling, enlarged liver, or neck-vein fullness.
\medterm{Tricuspid Valve} The heart valve between the right upper and right lower chambers. It opens to let blood move forward and closes to prevent backflow; leakage or narrowing can strain the heart.
\medterm{Tricyclic Antidepressant Drugs} Antidepressants that inhibit norepinephrine and serotonin reuptake and can cause anticholinergic, cardiac, and toxic effects.
\medterm{Tricyclic Antidepressants} Older antidepressants that non-selectively inhibit the reuptake of serotonin and
norepinephrine. They include amitriptyline, nortriptyline, imipramine, clomipramine, and
desipramine.

\medterm{Trifluoperazine} A phenothiazine antipsychotic used for selected psychotic disorders. It can cause drowsiness, movement disorders, and other serious effects, so treatment requires medical supervision.
\medterm{Trigeminal Nerve} The fifth cranial nerve, carrying sensation from the face and controlling muscles used for chewing.
\medterm{Trigeminal Neuralgia} Sudden, severe, electric-shock-like pain in one side of the face along one or more branches of the trigeminal nerve. Attacks may be triggered by touching the face, brushing teeth, eating, speaking, or a light breeze.

\synonyms
Neuralgia, Trigeminal
Tic Douloureux (Trigeminal Neuralgia)

\medterm{Trigger} A trigger is an event, exposure, stimulus, or condition that initiates or worsens a symptom, disease episode, behavior, or physiologic response.

\medterm{Trigger Finger} A condition in which a finger tendon catches as it moves through a tight covering. It can cause pain, clicking, stiffness, or locking and may improve with rest, injections, splinting, or surgery.

\synonyms
Stenosing Tenosynovitis
Trigger Thumb (Trigger Finger)

\medterm{Trigger Thumb} Stenosing tenosynovitis causing painful catching or locking of the thumb.
\medterm{Trigger Thumb (Trigger Finger)} Alternate index label for Trigger Finger.


\medterm{Triglyceride} A lipid composed of glycerol and three fatty acids, transported in blood and stored in adipose tissue.
\medterm{Triglyceride Level} A triglyceride level measures a type of fat in the blood; persistently high levels can increase cardiovascular risk and, when very high, cause pancreatitis.

\medterm{Triglycerides} Triglycerides are blood fats made of glycerol and three fatty acids; elevated levels contribute to cardiovascular risk and, when very high, pancreatitis risk.


\medterm{Trigone} A triangular area, especially the smooth region at the base of the urinary bladder between the ureter openings and urine outlet. It helps direct urine flow and can be affected by inflammation or tumors.
\medterm{Trimester} One of the three approximately three-month periods of pregnancy. Each trimester has typical developmental milestones and screening or care needs.
\medterm{Trimethoprim} An antibiotic that blocks bacterial folate production and is often combined with sulfamethoxazole. It treats selected infections but can cause rash, blood-cell changes, high potassium, kidney effects, or medicine interactions.
\medterm{Trimethoprim-Sulfamethoxazole} A fixed-dose combination of trimethoprim and sulfamethoxazole that synergistically
inhibits folate synthesis. TMP-SMX is indicated for urinary tract infections,
Pneumocystis jirovecii pneumonia, and community-acquired pneumonia. The combination is
bactericidal against many gram-positive and gram-negative organisms.

\medterm{Trimethylaminuria} A metabolic disorder in which impaired oxidation of trimethylamine causes a
characteristic fish-like body odor in sweat, urine, and breath. It may be inherited or
acquired and can cause substantial psychosocial distress despite limited physical
illness.

\medterm{Trimipramine} A sedating tricyclic antidepressant used for depression in some countries. Its anticholinergic and cardiovascular effects can cause dry mouth, constipation, blurred vision, drowsiness, orthostatic hypotension, or dangerous toxicity in overdose; interactions and suicide risk require clinical monitoring.
\medterm{Trinitrin} Another name for glyceryl trinitrate or nitroglycerin, a nitrate vasodilator.
\medterm{Trinitrophenol} Picric acid, a corrosive and potentially explosive chemical once used in medicine, laboratories, and industry.
\medterm{Trinucleotide Repeat} Expansion of three-nucleotide sequences causes diseases: Huntington, Fragile X
(CGG), myotonic dystrophy (CTG), Friedreich ataxia (GAA). Anticipation common.

\medterm{Trip} A trip is a journey or an accidental stumble; in drug-related language it can also mean an episode of altered perception after a hallucinogen.

\medterm{Tripe Palms} Diffuse, velvety, thickened hyperkeratosis of the palmar skin with a cerebriform,
corrugated appearance resembling the lining of a cow's stomach (tripe), which is
strongly associated with underlying malignancy, most commonly gastric adenocarcinoma. It
may occur alone or in conjunction with acanthosis nigricans.

\medterm{Triple X Syndrome} A sex-chromosome variation in which a female has an extra X chromosome. Many affected people have typical sexual development and fertility, while some have tall stature, speech or learning difficulties, coordination problems, or anxiety; severity varies widely.

\synonyms
Trisomy X; 47,XXX Syndrome

\medterm{Triple-Negative Breast Cancer} Breast cancer lacking ER, PR, HER2. Aggressive, limited treatment options.

\medterm{Triplets} Three offspring born from the same pregnancy or three infants born in a multiple birth.
\medterm{Triploidy Syndrome} A chromosomal condition in which cells contain three complete haploid chromosome sets.
It causes severe fetal growth restriction and multiple congenital anomalies and is
usually incompatible with prolonged survival; partial or mosaic forms vary.

\medterm{Triptan} A class of medications used for acute migraine treatment, acting as 5-HT1B/1D receptor
agonists. 5-HT1B receptors are located on cranial blood vessels; activation causes
vasoconstriction of dilated meningeal arteries.

\medterm{Triptans (Triptan)} Triptans are serotonin-receptor agonists used to abort migraine attacks by constricting cranial vessels and reducing trigeminal signaling.

\medterm{Triquetrum} Pyramidal proximal row carpal bone on the ulnar side, articulating with lunate, hamate,
pisiform, and TFCC. Second most commonly fractured carpal bone. Contributes to proximal
carpal row and ulnar wrist stability. Pisiform articulates with its volar surface.

\medterm{Trismus} Restricted ability to open the mouth, commonly caused by tetanus, dental infection, or muscle spasm.
\medterm{Trisodium Phosphate Poisoning} Poisoning from swallowing or contacting trisodium phosphate, a strong alkaline cleaner. It can burn the mouth, throat, eyes, and skin and may cause vomiting or breathing problems.

\medterm{Trisomy} Presence of three copies of a chromosome instead of two. Trisomy 21 (Down syndrome):
intellectual disability, characteristic facies, heart defects. Trisomy 18 (Edwards):
severe intellectual disability, multiple anomalies. Trisomy 13 (Patau): severe
intellectual disability, holoprosencephaly, cleft lip/palate.

\medterm{Trisomy 13} Patau syndrome, a chromosomal disorder caused by an extra chromosome 13 and associated with severe developmental, brain, facial, heart, kidney, and limb abnormalities.

\medterm{Trisomy 18} Edwards syndrome, a chromosomal disorder caused by an extra chromosome 18 and associated with growth restriction, characteristic craniofacial and limb findings, heart disease, and severe developmental impairment.

\medterm{Trisomy 18 (Edwards Syndrome)} Alternate index label for Edwards Syndrome.

\medterm{Trisomy 21} Down syndrome: the most common chromosomal disorder, caused by an extra chromosome 21
(trisomy 21; ~95 percent meiotic nondisjunction, mosaicism, translocation).

\medterm{Trisomy 21 (Down Syndrome Overview)} Alternate index label for Down Syndrome Overview.

\medterm{Trisomy 21 Syndrome} Alternate index label; see \textit{Down syndrome}.

\medterm{Trisomy X} Alternate index label; see \textit{Triple X syndrome}.

\medterm{Trocar} A pointed instrument placed inside a hollow tube to enter a body cavity or tissue. It is used in minimally invasive procedures to create an opening for cameras, instruments, or drainage tubes.
\medterm{Trochanter} A bony projection on the thigh bone, especially the larger outer or smaller inner projection. Muscles and tendons attach there, and irritation around the greater trochanter can cause outer-hip pain.
\medterm{Trochanteric Bursitis (Hip Bursitis)} Alternate index label for Hip Bursitis.

\medterm{Troche} A troche is a small medicated lozenge that dissolves slowly in the mouth or throat for local or systemic drug delivery.

\medterm{Trochlear Nerve} Cranial nerve IV, which supplies the superior oblique eye muscle. Injury can cause vertical or tilted double vision, often worse when looking down or inward.
\medterm{Trophic} Relating to the nutrition, growth, development, or maintenance of tissue. Trophic changes may occur when nerve supply or blood flow is impaired.
\medterm{Trophoblast} The outer cell layer of an early embryo that attaches to the uterus and helps form the placenta. It supports implantation, exchanges nutrients, and produces hormones needed to maintain early pregnancy.
\medterm{Trophoderm} The outer cell layer of a blastocyst, the early stage before implantation. It attaches the embryo to the uterus and develops into tissues that help form the placenta and support early pregnancy.

\medterm{Trophozoite} The active, feeding, and multiplying stage of a protozoan parasite.
\medterm{Tropical Chronic Pancreatitis} A historically recognized form of chronic pancreatitis reported mainly in tropical and subtropical regions, often beginning at a younger age and associated with recurrent abdominal pain, pancreatic duct stones, and diabetes. Its causes are multifactorial and may include genetic susceptibility, nutritional factors, and environmental exposures; the term should not be applied solely because a patient lives in a tropical region.

\medterm{Tropical Diseases} Illnesses more common in warm tropical or subtropical regions, including malaria, dengue, schistosomiasis, and other neglected infections.
\medterm{Tropical Medicine} The medical specialty concerned with diseases common in tropical and subtropical regions and illnesses related to travel. It includes infections, parasites, insect-borne diseases, environmental hazards, prevention, diagnosis, and treatment.
\medterm{Tropical Oil} Tropical oils are plant-derived oils from warm regions, such as coconut or palm oil; their nutritional effects depend on the specific oil and amount consumed.

\medterm{Tropical Sprue} A disorder causing long-lasting diarrhea and poor absorption of nutrients, reported mainly in tropical and subtropical areas. It can lead to weight loss and deficiencies of folate, iron, and vitamin B12. Testing helps exclude infection and other causes; treatment usually includes antibiotics and replacement of missing nutrients.

\medterm{Tropical Ulcer} A chronic ulcer occurring in tropical regions, often associated with infection, trauma, or poor wound healing.
\medterm{Troponin} A protein released into the blood when heart muscle is injured. A raised level can support a heart-attack diagnosis, but it may also result from other heart, lung, kidney, or severe systemic illnesses.

\medterm{Troponin Test} A troponin test measures heart-muscle proteins in the blood and helps assess possible heart injury, especially during suspected heart attack.

\medterm{Trouble Swallowing} Difficulty moving food or liquid from the mouth to the stomach. It may cause coughing, choking, pain, or a feeling that food is stuck in the throat or chest.

\medterm{Trousseau Sign Of Latent Tetany} A sensitive and specific clinical sign of hypocalcemia. Elicited by inflating a blood
pressure cuff on the upper arm to a pressure greater than systemic systolic blood
pressure for 3 minutes.

\medterm{Trousseau Syndrome} Recurrent or migratory venous thrombosis associated with active malignancy, particularly
mucin-producing adenocarcinoma. It reflects tumor-associated activation of coagulation
and is distinct from Trousseau sign of latent tetany.

\medterm{Trp} Trp is the standard abbreviation for tryptophan, an essential amino acid used in protein synthesis and serotonin and niacin pathways.

\medterm{True Ribs} Ribs 1-7 articulating directly with the sternum via their own costal cartilages. Each
has a head, neck, tubercle, and shaft. The first two are short and strongly curved for
neck/shoulder muscle attachment. Costal cartilages provide thoracic cage elasticity for
respiratory expansion.

\medterm{Truncation} A shortened or incomplete form caused by removal of part of a sequence, structure, chromosome, protein, or process; the effect depends on what was lost and where.

\medterm{Truncus Arteriosus} A birth defect in which one large blood vessel leaves the heart instead of separate vessels for the body and lungs. A large hole between the lower chambers usually allows oxygen-poor and oxygen-rich blood to mix, causing blue skin and heart failure. Early surgery is usually required.

\medterm{Trunk} The central part of the body excluding the head and limbs, or a major nerve or vessel trunk.
\medterm{Truss} A supportive appliance used historically for a hernia; definitive management depends on the type and complications.
\medterm{Trypanosoma} A genus of protozoan parasites causing African and American trypanosomiasis.
\medterm{Trypanosoma Brucei} A protozoan parasite transmitted by tsetse flies that causes African trypanosomiasis, or sleeping sickness.
The human-infecting subspecies differ in geographic distribution and typical disease
course; infection can progress from fever and lymph-node enlargement to neurologic
involvement.

\medterm{Trypanosomiasis} Disease caused by Trypanosoma parasites, including African sleeping sickness and Chagas disease. Infected insects spread them, and illness may affect blood, lymph nodes, the nervous system, heart, or digestive organs.
\medterm{Trypanosomiasis, American} Alternate index label; see \textit{Chagas disease}.

\medterm{Trypsin} A pancreatic proteolytic enzyme secreted as trypsinogen and activated in the small intestine, where it helps digest dietary proteins and activates other pancreatic enzymes.
\medterm{Trypsin And Chymotrypsin In Stool} Stool trypsin and chymotrypsin testing measures digestive enzymes and may help assess pancreatic enzyme production, although other tests are often preferred.

\medterm{Trypsinogen Test} A trypsinogen test measures a pancreatic enzyme precursor and may help evaluate pancreatic function or newborn screening for cystic fibrosis.

\medterm{Tryptophan} An amino acid the body cannot make and must obtain from food. It is used to make proteins and can be changed into serotonin, melatonin, and niacin, which help regulate mood, sleep, and other body functions.
\medterm{Tsetse Fly} A blood-feeding fly found in parts of sub-Saharan Africa. Its bite can spread Trypanosoma parasites, which cause African sleeping sickness; infection may cause fever, headaches, behavior changes, sleep disturbance, and serious nervous-system disease.
\medterm{TSH} Thyroid-stimulating hormone, a pituitary hormone that stimulates thyroid hormone production; its blood concentration is the main initial test for many thyroid disorders.

\medterm{TSH Test} A TSH test measures thyroid-stimulating hormone in the blood and is a common initial test of thyroid function.

\medterm{TSH-Receptor Antibody} An antibody that binds the thyrotropin receptor. Stimulating antibodies cause Graves
disease, while blocking antibodies can cause hypothyroidism; measurement supports
diagnosis and can help assess fetal or neonatal risk in pregnancy.

\medterm{TSI Test} A TSI test detects thyroid-stimulating immunoglobulins and can support diagnosis or monitoring of Graves disease.
\medterm{Tsutsugamushi Fever} Scrub typhus, a mite-borne infection caused by Orientia tsutsugamushi. It may cause fever, headache, rash, an eschar, and organ complications and is treated with an appropriate antibiotic.
\medterm{Tubal Ligation} Surgical sterilization by occluding the fallopian tubes, preventing the ovum from
reaching the uterus. Postpartum tubal ligation: after vaginal or cesarean delivery
(Pomeroy technique, Parkland). Interval: laparoscopic (bands, clips, bipolar cautery) or
mini-laparotomy (Hulka or Filshie clip).


\medterm{Tubal Ligation Reversal} Microsurgical reconnection of previously divided or blocked fallopian-tube segments to restore a possible path for sperm and egg; success depends on remaining tube length, age, scarring, and original method.

\medterm{Tubal Occlusion Procedure, Selective} Selective tubal occlusion blocks a fallopian tube to prevent sperm and egg from meeting and is a permanent contraceptive procedure.

\medterm{Tubal Pregnancy} An ectopic pregnancy that implants in a fallopian tube instead of the uterus. It can rupture and cause internal bleeding; one-sided pain, vaginal bleeding, dizziness, or fainting needs urgent evaluation.

\medterm{Tuber} A rounded swelling or raised part of a bone or other body structure. In anatomy, tubers provide attachment for muscles or ligaments; the word can also mean a plant storage organ and is not itself a diagnosis.

\medterm{Tubercle} A small rounded projection or nodule; it may refer to bone anatomy or a tuberculous lesion.
\medterm{Tuberculide} A skin eruption caused by an immune reaction to tuberculosis elsewhere in the body. The skin lesions usually contain no bacteria and may improve when the underlying tuberculosis is treated.
\medterm{Tuberculin} Protein material from Mycobacterium tuberculosis used in a skin test for immune sensitization. A positive result suggests exposure or vaccination but does not by itself prove active tuberculosis or show when exposure occurred.
\medterm{Tuberculin Skin Test} Intradermal injection of purified protein derivative to test for Mycobacterium
tuberculosis infection. Read at 48-72 hours measuring induration (not erythema).
Positive: greater than 5 mm in HIV or close contacts; greater than 10 mm in high-risk
groups; greater than 15 mm in low-risk populations.

\medterm{Tuberculin Test} A skin test that checks whether the immune system has reacted to tuberculosis proteins. A positive result suggests infection or previous vaccination and does not by itself prove active disease.

\medterm{Tuberculoid Leprosy} The paucibacillary pole of the leprosy spectrum characterised by strong cell-mediated
immunity to Mycobacterium leprae, resulting in well-formed granulomas and few or no
organisms in tissue.


\medterm{Tuberculosis} An infection caused by Mycobacterium tuberculosis, most often affecting the lungs. It spreads through airborne particles and may cause cough, fever, night sweats, weight loss, or silent infection without symptoms.

\synonyms
TB


\medterm{Tuberculosis In Children} Infection with Mycobacterium tuberculosis, often from household contacts. Latent TB
infection (LTBI) positive tuberculin skin test or interferon-gamma release assay with no
symptoms or radiographic findings. Active TB presents with cough, fever, weight loss,
and lymphadenopathy. Diagnosis by PPD, chest X-ray, and sputum culture.

\medterm{Tuberculosis, MDR} CEA is a blood tumor marker used mainly to help monitor some colorectal and other cancers; it is not specific enough to diagnose cancer alone.
\medterm{Tuberculous Meningitis} The most severe, life-threatening form of extrapulmonary tuberculosis, resulting from hematogenous dissemination of \textit{Mycobacterium tuberculosis} and the rupture of subependymal Rich foci into the subarachnoid space. Causes thick gelatinous basilar exudate, cranial nerve palsies, communicating hydrocephalus, and cerebral vasculitis with stroke.


\medterm{Tuberous Sclerosis} A rare genetic disorder, also called tuberous sclerosis complex, that causes growths in the brain and
other organs, including the skin, kidneys, heart, and lungs. It may cause seizures, developmental or
behavioral problems, and skin changes; the severity varies widely.
\medterm{Tuberous Sclerosis Complex} A genetic condition in which noncancerous growths develop in several organs, especially the brain, skin, kidneys, heart, and lungs. It may cause seizures, developmental or learning difficulties, skin changes, kidney problems, or breathing disease. Lifelong monitoring and treatment are tailored to the organs affected.

\medterm{Tubes} Hollow structures that carry air, fluid, food, or medical instruments. Examples include breathing tubes, fallopian tubes, drainage tubes, and blood vessels; the specific structure must be named for the term to be clinically precise.

\medterm{Tubes Ear Problems (Eustachian Tube Problems)} Alternate index label for Eustachian Tube Problems.

\medterm{Tubo-Ovarian Abscess} A complex infected mass involving the fallopian tube and ovary, usually arising as a complication of pelvic inflammatory disease. It may cause pelvic pain, fever, and sepsis and requires prompt antimicrobial and sometimes drainage treatment.

\medterm{Tubocurarine} A historical muscle-paralyzing medicine derived from curare. It blocks nerve signals at skeletal-muscle junctions, including signals needed for breathing, but does not cause unconsciousness or relieve pain; safer modern drugs have largely replaced it.
\medterm{Tubography} Contrast imaging of the fallopian tubes, usually as part of hysterosalpingography, to see whether they are open. It may help assess infertility but can cause temporary cramping, bleeding, infection, or rarely injury.

\medterm{Tubular Adenoma} A benign glandular polyp, most often in the colon, that can develop dysplastic changes and is considered a precursor lesion for some colorectal cancers.

\medterm{Tubular Reabsorption} The process by which kidney tubules return useful water and dissolved substances from filtered fluid to the blood. This prevents the loss of most water, salt, glucose, and other substances in urine.

\medterm{Tubular Secretion} The transport of substances from the peritubular capillary blood into the tubular lumen,
a second mechanism (besides filtration) by which the kidney eliminates waste and
regulates ions.

\medterm{Tubule} A small tube in the body. In the kidney, tubules change filtered fluid by taking useful water and substances back into the blood and leaving waste to form urine; other tubules carry fluids in different organs.
\medterm{Tubuloglomerular Feedback} A kidney control system that adjusts blood flow and filtration according to the amount of salt reaching a sensing area in the tubule. It helps keep urine production and body fluid balance stable.

\medterm{Tularaemia} A zoonotic infection caused by Francisella tularensis, acquired from animals, ticks, contaminated water, or aerosols.
\medterm{Tularemia} Tularemia is infection with Francisella tularensis acquired from ticks, animals, contaminated water, or aerosols and causing ulcer-glandular, pulmonary, or other disease.

\medterm{Tularemia Blood Test} A tularemia blood test looks for antibodies or other evidence of infection with Francisella tularensis; antibodies may not appear early in illness.

\medterm{Tulle Gras} A paraffin-impregnated gauze dressing used to cover wounds and burns.
\medterm{Tumescence} Swelling or enlargement caused by increased blood flow, fluid, or tissue growth. The term commonly describes erection but can also refer to temporary or abnormal enlargement in other body tissues.
\medterm{Tumor} A tumor is an abnormal mass of cells caused by excessive or uncontrolled growth; it may be benign or malignant.

\medterm{Tumor Adrenal Gland (Pheochromocytoma)} Alternate index label for Pheochromocytoma.

\medterm{Tumor Board} A multidisciplinary conference of oncologists, surgeons, radiologists, and pathologists
who review individual cancer cases to coordinate diagnosis and treatment planning.

\medterm{Tumor Brain Cancer (Brain Cancer)} Alternate index label for Brain Cancer.

\medterm{Tumor Burden} The total amount of cancer in the body, estimated from the number, size, and extent of
tumor deposits; higher burden is generally associated with worse prognosis.

\medterm{Tumor Differentiation} How closely tumor cells resemble the normal cells from which they arose. Well-differentiated tumors look and behave more like normal tissue, while poorly differentiated tumors often behave more aggressively.
\medterm{Tumor Embolism} Tumor embolism occurs when malignant cells, tumor fragments, or tumor thrombi enter the
venous circulation and travel to distant sites. The pulmonary vasculature is the most
common site of embolization, as venous blood from most organs drains to the lungs
through the pulmonary arteries.

\medterm{Tumor Grade} A histologic measure of how much tumor cells resemble normal tissue; low-grade
(well-differentiated) tumors behave less aggressively than high-grade (poorly
differentiated) tumors. Grading is a prognostic, not staging, feature.

\medterm{Tumor Grading} A description of a tumor based on how abnormal the cancer cells look under a microscope
compared to normal cells; a lower grade indicates slower-growing cancer.

\medterm{Tumor Heterogeneity} The presence of genetically, epigenetically, or phenotypically distinct subpopulations
within the same tumor or between primary and metastatic sites. Heterogeneity can
influence prognosis, sampling, treatment response, and resistance.

\medterm{Tumor Lysis Syndrome} A medical emergency caused by rapid cancer-cell breakdown, releasing substances that disturb blood chemistry. It can damage kidneys, trigger seizures or abnormal heart rhythms, and requires prevention, monitoring, and urgent treatment.

\medterm{Tumor Marker} A substance found in blood, urine, tissue, or another body fluid that may be produced by cancer or by the body’s response to it. Tumor markers can help monitor treatment or recurrence and sometimes guide diagnosis, but most are not specific enough to prove cancer alone.

\medterm{Tumor Marker, CEA} CEA is a blood tumor marker used mainly to help monitor some colorectal and other cancers; it is not specific enough to diagnose cancer alone.

\medterm{Tumor Marker, NSE} NSE is a protein marker that may rise in some neuroendocrine tumors or nerve-cell injury but is not diagnostic by itself.

\medterm{Tumor Microenvironment} The cells, blood vessels, connective tissue, and chemical signals surrounding a tumor. These nearby tissues can support or limit tumor growth, affect spread, influence the immune response, and change how well treatment works.
\medterm{Tumor Mutational Burden} The number of somatic mutations per unit of coding DNA in a tumor specimen. TMB may
correlate with production of neoantigens and response to some immune checkpoint
inhibitors, but interpretation requires a validated assay and disease-specific context.

\medterm{Tumor Necrosis Factor} A proinflammatory cytokine (TNF-alpha) produced mainly by macrophages, acting on
TNFR1/TNFR2 receptors. Functions: activation of endothelium and leukocytes, fever,
cachexia, apoptosis, and coordination of inflammation and the acute-phase response.

\medterm{Tumor Necrosis Factor Receptor-Associated Periodic Syndrome} A rare inherited autoinflammatory disorder caused by variants in the TNFRSF1A gene, producing recurrent prolonged fever, muscle pain, abdominal pain, rash, and inflammation; it is also called TRAPS.

\medterm{Tumor Staging} A system of describing the size of a cancer and how far it has spread throughout the
body from its original location (e.g., using the TNM system).

\medterm{Tumor Suppressor} A gene or protein that helps limit cell growth, repair DNA damage, or promote removal of abnormal cells. Loss of its function can
contribute to cancer.

\medterm{Tumor Suppressor Gene} A gene that helps prevent uncontrolled cell growth by slowing division, repairing DNA damage, or directing abnormal cells to die. If both copies or the gene’s function are lost or disabled, cells may grow abnormally and cancer risk can increase.

\medterm{Tumor Testicle (Testicular Disorders)} Alternate index label for Testicular Disorders.

\medterm{Tumor, Salivary Gland} Alternate index label; see \textit{Salivary gland tumors}.

\medterm{Tumor, Spinal} Alternate index label; see \textit{Spinal tumor and spinal mass}.

\medterm{Tumor-Suppressor Gene} A gene whose normal product restrains cell proliferation, promotes DNA repair, induces
apoptosis, or maintains genomic stability. Loss or inactivation of both functional
copies is common in the development of many cancers, although mechanisms vary by gene.

\medterm{Tumors Uterine (Uterine Fibroids)} Alternate index label for Uterine Fibroids.

\medterm{Tumour} An abnormal mass caused by uncontrolled or inappropriate cell growth; it may be noncancerous or cancerous.
\medterm{Tumour (Neoplasm)} An abnormal mass of tissue produced by uncontrolled or inappropriate cell growth; it may be benign or malignant.

\synonyms
Neoplasm

\medterm{Tungiasis} A skin infestation caused by the sand flea Tunga penetrans, producing an itchy or painful lesion, usually on
the feet.

\medterm{Tunica Adventitia} The outer connective-tissue layer of a blood-vessel wall. It contains collagen and
elastic fibers that provide structural support, along with nerves and, in larger
vessels, small vessels called vasa vasorum that supply the outer wall. It is relatively
prominent in veins and helps anchor vessels to surrounding tissues.

\medterm{Tunica Media} The middle layer of a blood-vessel wall, located between the tunica intima and tunica
adventitia. It consists primarily of smooth muscle and elastic fibers; its contraction
or relaxation changes vascular diameter, regulating blood pressure and tissue perfusion.
It is thickest in arteries and thinner in veins.

\medterm{Tunnel Vision} Loss of peripheral visual field leaving a narrowed central field, caused by glaucoma, retinal disease, optic-nerve or brain lesions, severe hypoxia, or other conditions.

\medterm{Turbinate Surgery} An operation that reduces or reshapes swollen tissues along the side walls of the nose to improve airflow. It may help persistent nasal blockage when other treatment fails; bleeding, crusting, dryness, scarring, or ongoing obstruction are possible.

\medterm{Turbinates} Curved bony shelves on the lateral walls of the nasal cavity, covered by vascular mucosa. They increase surface area and help
warm, humidify, and filter inhaled air.

\medterm{Turbinates Hypertrophy} Chronic enlargement of the nasal turbinates, most commonly the inferior turbinates,
resulting in nasal airway obstruction, which can be caused by chronic inflammation from
allergic rhinitis, non-allergic rhinitis, chronic sinusitis, hormonal changes, or
vasomotor dysfunction.

\medterm{Turcot Syndrome} A rare inherited cancer-predisposition pattern involving colon polyps or colon cancer together with a brain tumor. It can result from mismatch-repair gene changes or the APC gene change seen in familial adenomatous polyposis, so genetic counseling and screening are important.

\synonyms
Brain Tumor-Polyposis Syndrome

\medterm{Turgor} The firmness of tissue produced by fluid pressure within cells; skin turgor can help assess dehydration.
\medterm{Turmeric} Turmeric is a plant spice containing curcuminoids and used in foods and supplements; evidence varies by condition and interactions or liver injury are possible.

\medterm{Turner Hypoplasia} Enamel defect on permanent tooth from infection of overlying primary tooth. Localized hypoplasia. A defect in permanent tooth enamel caused by infection or injury involving the baby tooth above it.
\medterm{Turner Syndrome} A chromosome condition affecting females in which one X chromosome is missing or altered. It commonly causes short stature and reduced ovarian function, with delayed or absent puberty and infertility. Heart, kidney, hearing, and thyroid problems may also occur, so lifelong medical follow-up is important.

\medterm{Turner'S Syndrome} A chromosome condition affecting females in which one X chromosome is missing or altered. It commonly causes short stature and reduced ovarian function, with delayed or absent puberty and infertility. Heart, kidney, hearing, and thyroid problems may also occur, so lifelong medical follow-up is important.
\medterm{Turning Patients Over In Bed} Repositioning a patient at planned intervals to protect skin, improve comfort and breathing, prevent pressure injuries and contractures, and reduce complications of immobility.

\medterm{Turpentine Oil Poisoning} Poisoning from swallowing or inhaling turpentine oil. It can irritate or damage the lungs, stomach, and nervous system and may cause chemical pneumonia.

\medterm{Tussis} Cough; a forceful reflex expiration that helps clear secretions, irritants, or foreign material from the airways.

\medterm{TVS} TVS commonly means transvaginal sonography, an ultrasound examination using a probe in the vagina to assess pelvic organs and early pregnancy.


\medterm{Twilight Sleep} Twilight sleep is an older term for a sedated, partly conscious state used during procedures; modern analgesia and anesthesia specify the medicines and monitoring used.

\medterm{Twin} A twin is one of two offspring from the same pregnancy, either identical from one fertilized egg or fraternal from two eggs and two sperm.

\medterm{Twin Placenta} Placental anatomy in a twin pregnancy, classified by chorionicity and amnionicity; these features determine risks such as twin-to-twin transfusion and guide antenatal monitoring.

\medterm{Twin Pregnancy} A pregnancy involving two fetuses. Clinical risk depends on chorionicity and amnionicity, with increased risks including preterm birth, fetal growth problems, and hypertensive disorders.

\medterm{Twin-To-Twin Transfusion Syndrome} A complication of monochorionic twin pregnancy in which shared placental vessels cause
unequal blood flow between the twins.

\medterm{Twin-Twin Transfusion Syndrome} A complication of monochorionic diamniotic twin pregnancies where placental vascular
anastomoses (arteriovenous connections) cause unbalanced blood flow from one twin (the
donor) to the other (the recipient). The donor twin becomes hypovolemic, anemic, and
growth-restricted with oligohydramnios ("stuck twin").

\medterm{Twins} Two offspring from the same pregnancy, either monozygotic or dizygotic.
\medterm{Twitch} A brief involuntary muscle contraction or movement, which may be benign, medication-related, stress-related, or a sign of nerve, muscle, or brain disease.

\medterm{Twitching} Repeated involuntary muscle contractions or fasciculations, commonly caused by fatigue, stimulants, electrolyte disturbance, nerve irritation, medicines, or neurologic disease.


\medterm{Tympanic} Tympanic means relating to the eardrum or middle ear, or resembling a drum in shape or sound.

\medterm{Tympanic Bulla} The tympanic bulla is a hollow or expanded bony chamber surrounding the middle ear, especially prominent in many mammals.

\medterm{Tympanic Cavity} The air-filled space of the middle ear containing the malleus, incus, and stapes. The air-filled middle-ear chamber containing three small bones that transmit sound vibrations toward the inner ear.
\medterm{Tympanic Hyperectasis} Abnormal outward bulging or retraction of the tympanic membrane from chronic middle ear
pressure changes, potentially causing conductive hearing loss and requiring tympanostomy
tube placement.

\medterm{Tympanic Membrane} The thin eardrum separating the outer ear canal from the middle ear. Sound makes it vibrate, and those vibrations move the middle-ear bones; infection, injury, or pressure can tear it.
\medterm{Tympanic Membrane Perforation} A tear or hole in the eardrum caused by infection, trauma, pressure changes, or instrumentation, potentially producing pain, drainage, or conductive hearing loss.

\medterm{Tympanites} Abdominal distension caused by excessive gas in the stomach or intestines, producing a drum-like sound on percussion and sometimes pain, belching, or altered bowel movements.
\medterm{Tympanometry} An objective test that varies air pressure in the ear canal while measuring movement of the tympanic membrane and middle-ear pressure. It helps assess middle-ear fluid, eustachian-tube dysfunction, perforation, and ossicular or tympanic-membrane abnormalities.

\medterm{Tympanoplasty} Surgical repair or reconstruction of the tympanic membrane and, when necessary, the middle-ear ossicles to close a perforation or improve hearing.

\medterm{Tympanosclerosis} Scarring and calcification of the tympanic membrane or middle-ear lining, often following recurrent infection or previous inflammation.
It may be asymptomatic or cause conductive hearing loss when the ossicles are affected.

\synonyms
Tympanic-membrane sclerosis; myringosclerosis when limited to the tympanic membrane.

\medterm{Tympanostomy Tubes} Small tubes inserted through the tympanic membrane to ventilate the middle ear and
prevent fluid accumulation. Indicated for recurrent acute otitis media (3 episodes in 6
months or 4 in 12 months) or persistent otitis media with effusion greater than 3 months
with hearing loss. Performed as outpatient procedure under general anesthesia.

\medterm{Tympanum} The middle-ear cavity or, in some contexts, the tympanic membrane.
\medterm{Tympany} Tympany is a drumlike sound on percussion, often produced by gas in a hollow organ such as the stomach or intestine.

\medterm{Type 1 Diabetes} Type 1 diabetes is autoimmune destruction of pancreatic beta cells causing absolute insulin deficiency, hyperglycemia, and risk of diabetic ketoacidosis without insulin replacement.

\synonyms
Diabetes, Type 1
Diabetes Type 1 (Type 1 Diabetes)


\medterm{Type 1 Diabetes In Children} Type 1 diabetes diagnosed during childhood, causing high blood sugar and requiring insulin, monitoring, nutrition planning, and education.

\medterm{Type 1 Diabetes Mellitus} Type 1 diabetes is an autoimmune disease in which the immune system attacks and destroys
the insulin-producing beta cells of the pancreas, resulting in absolute insulin
deficiency.

\medterm{Type 2 Diabetes} Type 2 diabetes is chronic hyperglycemia caused by insulin resistance and progressive beta-cell dysfunction, with risks to the eyes, kidneys, nerves, heart, and blood vessels.

\synonyms
Diabetes, Type 2
Diabetes Type 2 (Type 2 Diabetes)

\medterm{Type 2 Diabetes - Oral Medicines} Oral medicines for type 2 diabetes lower glucose through effects on insulin, liver glucose production, kidney reabsorption, or intestinal absorption; choice depends on health and treatment goals.



\medterm{Type 2 Diabetes In Children} Type 2 diabetes beginning during childhood or adolescence, often involving insulin resistance and requiring lifestyle and medical management.

\medterm{Type 2 Diabetes Mellitus} Type 2 diabetes is a chronic metabolic disorder characterized by insulin resistance and
progressive pancreatic beta-cell dysfunction, leading to hyperglycemia.

\medterm{Type 3B Diabetes} Type 3b diabetes is a form of diabetes caused by disease or removal of the pancreas, which can impair both insulin and digestive-enzyme production.

\medterm{Type I Hypersensitivity} An immediate (anaphylactic) hypersensitivity reaction mediated by allergen-specific IgE
bound to mast cells and basophils; re-exposure cross-links IgE, triggering degranulation
and release of histamine and mediators within minutes. Manifestations: allergic
rhinitis, asthma, urticaria, food allergy, and anaphylaxis.

\medterm{Type II Hypersensitivity} An antibody-mediated (cytotoxic) hypersensitivity reaction in which IgG or IgM binds to
cell-surface or matrix antigens, causing cell destruction by complement activation
(MAC), opsonization/phagocytosis, or antibody-dependent cellular cytotoxicity.

\medterm{Type III Hypersensitivity} An immune complex-mediated hypersensitivity reaction in which circulating
antigen-antibody complexes deposit in tissues and vessels, activating complement and
attracting neutrophils, causing inflammation and tissue damage. Local forms: Arthus
reaction, serum sickness.

\medterm{Type III Hypersensitivity Reaction} Tissue injury caused when antigen-antibody clusters deposit in tissues and activate inflammation. This reaction can damage blood vessels, joints, kidneys, skin, or other organs, as in serum sickness or lupus.

\medterm{Type IV Hypersensitivity} A delayed-type hypersensitivity (cell-mediated) reaction mediated by sensitized T cells
(not antibody), occurring 24-72 hours after antigen exposure.

\medterm{Type V Glycogen Storage Disease} An inherited muscle disorder, also called McArdle disease, in which muscle cells cannot properly break down glycogen for energy. Exercise can cause muscle pain, cramps, and sometimes dark urine from muscle injury.


\medterm{Types Of Health Care Providers} Health professionals and services that diagnose, treat, prevent, or support health needs, including physicians, nurses, pharmacists, therapists, dentists, mental-health professionals, and community services.


\medterm{Types Of Ileostomy} Types of ileostomy include end, loop, and continent procedures, which divert stool through an opening in the abdominal wall when the colon or rectum cannot be used normally.

\medterm{Typhlitis} Inflammation of the cecum in immunocompromised patients particularly those with
neutropenia from chemotherapy. Also known as neutropenic enterocolitis. Presents with
right lower quadrant pain, fever, and diarrhea. Diagnosis is by CT showing cecal wall
thickening.

\synonyms
neutropenic enterocolitis.

\medterm{Typhoid} A systemic infection caused by Salmonella enterica serovar Typhi, usually acquired through contaminated food or
water and characterized by prolonged fever and gastrointestinal or systemic symptoms.

\medterm{Typhoid Fever} A potentially serious whole-body infection caused by Salmonella Typhi, usually spread through contaminated food or water.

\synonyms
Enteric Fever (Typhoid Fever)

\medterm{Typhus} Typhus refers to several rickettsial infections transmitted by lice, fleas, or mites and characterized by fever, headache, rash, and variable systemic complications.
arthropod vectors. Three main forms: (1) Epidemic (louse-borne) typhus—Rickettsia
prowazekii, transmitted by the human body louse (Pediculus humanus corporis). Associated
with overcrowding, poverty, cold climate, and refugees ('war typhus').

\medterm{Typhus Fever} A group of infections caused by Rickettsia species and transmitted mainly by lice or fleas.
\medterm{Typical Antipsychotics} Typical antipsychotics, also known as first-generation antipsychotics, primarily block
dopamine D2 receptors in the mesolimbic pathway, reducing positive symptoms of
schizophrenia including hallucinations, delusions, and disorganized thinking.

\synonyms
first-generation antipsychotics, primarily block.

\medterm{Tyramine} A naturally occurring amine found in some aged or fermented foods. In people taking monoamine oxidase inhibitors, excessive tyramine can cause a dangerous rise in blood pressure.
\medterm{Tyrosinase Peptide} A short protein fragment related to tyrosinase, the enzyme that helps make melanin, the pigment in skin, hair, and eyes. Such peptides are mainly used in research.

\medterm{Tyrosine} An amino acid used to build proteins and make thyroid hormones, melanin, and adrenaline-like chemicals. The body usually makes enough from phenylalanine, but inherited metabolic disorders can disrupt this pathway.

\synonyms
Full term

\medterm{Tyrosinemia} An inherited disorder of tyrosine metabolism that causes accumulation of tyrosine or its metabolites. Type I can cause severe liver and kidney disease and is associated with fumarylacetoacetate hydrolase deficiency.

\medterm{Tzanck Smear} A microscope test of cells scraped from the base of a skin blister or sore. It may show changes suggesting herpesvirus infection or pemphigus, but it cannot reliably identify the cause and has largely been replaced by more specific tests.

\medterm{Tamiflu (Oseltamivir)} Tamiflu is the brand name for oseltamivir, an antiviral medicine used to treat or help prevent influenza. It works best when started early and does not replace flu vaccination.

\medterm{Tamoxifen} Tamoxifen is a selective estrogen-receptor modulator used mainly to treat or prevent some estrogen-sensitive breast cancers. It can increase the risk of blood clots and, rarely, uterine cancer.

\medterm{Tarlov Cysts} Tarlov cysts are fluid-filled sacs around the nerve roots, most often in the sacral spine. Many cause no symptoms, but larger or symptomatic cysts may cause pain, numbness, or bladder problems.

\medterm{Taste Receptor} A taste receptor is a protein on taste-bud cells that detects chemical substances in food, helping create taste sensations such as sweet, salty, sour, bitter, and umami.

\medterm{Tattoos} Tattoos are permanent designs made by placing pigment into the skin. Risks include infection, allergic reactions, scarring, and difficulty removing the pigment.

\medterm{Taurine} Taurine is a sulfur-containing compound involved in normal cell, nerve, muscle, and bile-acid function. It is present in foods and added to some energy drinks; high caffeine intake in such drinks can be harmful.

\medterm{Taxotere} Taxotere is the brand name for docetaxel, a chemotherapy medicine that interferes with cell division. It is used for several cancers and can cause low blood counts, infection risk, hair loss, and other side effects.

\medterm{T-Cell} A T-cell is a type of white blood cell that helps coordinate immune responses or directly destroys infected or abnormal cells. T cells mature in the thymus.

\medterm{Tea} Tea is a drink made by steeping plant leaves or other plant material in water. Some teas contain caffeine, and concentrated herbal products can interact with medicines.

\medterm{Teeth Whitening} Teeth whitening lightens tooth color by removing or chemically breaking down stains. Professional or properly used products are generally safe, but temporary sensitivity and gum irritation can occur.

\medterm{Teledermatology} Teledermatology is the use of video, photographs, or other digital communication to assess and manage skin conditions remotely. Some problems still require an in-person examination or biopsy.

\medterm{Tendinosis} Tendinosis is chronic degeneration and disorganization of tendon tissue, usually related to repeated loading rather than acute inflammation. Treatment may include activity modification and rehabilitation.

\medterm{Testicular Failure} Testicular failure occurs when the testes do not produce enough testosterone, sperm, or both. Causes include genetic conditions, injury, infection, medicines, and damage from chemotherapy or radiation.

\medterm{The Microbiome Project} The Microbiome Project refers to research that studies the microorganisms living in and on the human body and their genes. This research helps investigate links between microbes, health, and disease.

\medterm{Thimerosal} Thimerosal is a mercury-containing preservative used in some multidose vaccines and other medical products. The mercury form it contains is ethylmercury, which is cleared from the body more quickly than methylmercury.

\medterm{Threadworms} Threadworms, also called pinworms, are small intestinal worms that commonly cause itching around the anus at night. Treatment usually includes an antihelminthic medicine and hygiene measures for household contacts.

\medterm{Thymine} Thymine is one of the four main bases in DNA. It pairs with adenine through hydrogen bonds and helps store genetic information.

\medterm{Tick-borne Diseases} Tick-borne diseases are infections or other illnesses transmitted by ticks, including Lyme disease, babesiosis, ehrlichiosis, and some viral diseases. Removing ticks promptly can reduce risk.

\medterm{Tokophobia} Tokophobia is an intense fear of pregnancy or childbirth that can cause substantial distress or avoidance. Support may include psychological therapy, birth planning, and specialist maternity care.

\medterm{Tooth Brushing} Tooth brushing removes plaque and food from teeth and gums. Brushing twice daily with fluoride toothpaste helps prevent tooth decay and gum disease.

\medterm{Tooth Erosion} Tooth erosion is the chemical loss of tooth enamel caused by acids, not by bacteria. Acidic foods, drinks, reflux, and repeated vomiting can contribute to it.

\medterm{Traditional Medicine} Traditional medicine includes health practices, remedies, and healing systems developed within cultural traditions. Products and practices vary in evidence and safety, so they should not replace effective care without professional advice.

\medterm{Travel Vaccinations} Travel vaccinations are vaccines recommended or required for protection against infections encountered in particular countries or regions. Recommendations depend on the destination, activities, health conditions, and previous vaccinations.

\medterm{Trichinellosis} Trichinellosis is an infection caused by Trichinella worms, usually after eating undercooked infected meat. It may cause diarrhea followed by fever, muscle pain, and swelling around the eyes.

\medterm{Triclosan} Triclosan is an antimicrobial chemical formerly used in some soaps, cosmetics, and consumer products. Routine use in consumer products has limited benefit and may contribute to environmental and antimicrobial concerns.

\medterm{Tridecanoic Acid} Tridecanoic acid is a saturated fatty acid containing 13 carbon atoms. It occurs in small amounts in some foods and is mainly studied in nutrition and lipid research.

\medterm{Tropical Disease} A tropical disease is an infection or other health condition more common in tropical and subtropical regions, often because of local climate, vectors, or living conditions. Travel history helps guide diagnosis.

\medterm{Tubulointerstitial Nephropathy} Tubulointerstitial nephropathy is kidney damage involving the renal tubules and surrounding tissue. It may result from medicines, infection, toxins, autoimmune disease, or inherited disorders and can impair kidney function.
